\documentclass[10pt,letterpaper]{bp-template-concise}

\usepackage{colortbl}
\usepackage{makecell}
\colorlet{dkblue}{cyan!50!gray}
\definecolor{dkgreen}{rgb}{0,0.6,0}
\usetikzlibrary{arrows.meta,calc}
\newtcolorbox{stagecard2}[1][]{%
  enhanced,
  colback=SPred!4,
  colframe=SPred!45,
  coltitle=white,
  fonttitle=\bfseries\small,
  attach boxed title to top left={yshift=-2mm, xshift=4mm},
  boxed title style={colback=SPred, colframe=SPred,
                     sharp corners, boxrule=0pt,
                     inner ysep=3pt, inner xsep=6pt},
  sharp corners,
  left=6pt, right=6pt, top=10pt, bottom=6pt,
  #1
}

% Numbers for Section 4.6 (limited attention), generated by scripts/compute_attention.py
% Generated by scripts/compute_attention.py -- do not edit by hand.
% Values are provisional (band-midpoint energy costs; current registry match).
\newcommand{\attNDgtBApt}{37}
\newcommand{\attNMarketsApt}{45}
\newcommand{\attExcessApt}{4.5}
\newcommand{\attReachApt}{1.7}
\newcommand{\attShareUninfApt}{95}
\newcommand{\attCursedApt}{0.1}
\newcommand{\attLevelGapPreApt}{2.6}
\newcommand{\attLevelGapPostApt}{7.0}
\newcommand{\attEtaApt}{2.32}
\newcommand{\attRatioApt}{2.10}
\newcommand{\attCstarApt}{31}
\newcommand{\attCsilApt}{244}
\newcommand{\attPsiApt}{0.13}
\newcommand{\attPsiLoApt}{\ensuremath{-}0.12}
\newcommand{\attPsiHiApt}{0.37}
\newcommand{\attPsiSeApt}{0.13}
\newcommand{\attPsiOLSApt}{0.21}
\newcommand{\attPsiOLSseApt}{0.11}
\newcommand{\attPsiMedApt}{0.03}
\newcommand{\attPsiConstApt}{0.52}
\newcommand{\attIcptApt}{2.7}
\newcommand{\attIcptSeApt}{1.9}
\newcommand{\attTfullApt}{6.9}
\newcommand{\attIntervalHiApt}{0.13}
\newcommand{\attSlackApt}{0.90}
\newcommand{\attSlackTApt}{1.0}
\newcommand{\attCstarLoApt}{\ensuremath{-}31}
\newcommand{\attCstarHiApt}{89}
\newcommand{\attShareInBandApt}{0}
\newcommand{\attShareViolBottomApt}{67}
\newcommand{\attNViolApt}{21}
\newcommand{\attShareViolTopApt}{40}
\newcommand{\attBthreeApt}{4.1}
\newcommand{\attDthreeApt}{8.5}
\newcommand{\attLamApt}{6.1}
\newcommand{\attLamTApt}{6.78}
\newcommand{\attCzeroBaseApt}{2.4}
\newcommand{\attBandLoApt}{8.45}
\newcommand{\attBandHiApt}{8.54}
\newcommand{\attBandWidthApt}{0.09}
\newcommand{\attBandWidthBaseApt}{0.68}
\newcommand{\attFittedPremiumApt}{8.53}
\newcommand{\attSigmaOneApt}{6.3}
\newcommand{\attSigmaZeroApt}{8.0}
\newcommand{\attHighLossApt}{\ensuremath{-}4.5}
\newcommand{\attSalienceApt}{6.1}
\newcommand{\attEnforceSilApt}{0.02}
\newcommand{\attSensEApt}{0.07}
\newcommand{\attPsiLevROMApt}{0.08}
\newcommand{\attPsiLevSlopeApt}{0.81}
\newcommand{\attCzeroFitApt}{8.5}
\newcommand{\attCzeroCoreMeanApt}{2.4}
\newcommand{\attRmseCoreApt}{8.6}
\newcommand{\attRmseExtApt}{5.9}
\newcommand{\attCorrCoreApt}{0.40}
\newcommand{\attCorrExtApt}{0.43}
\newcommand{\attResidCoreApt}{6.1}
\newcommand{\attResidExtApt}{0.0}
\newcommand{\attRbarGapApt}{0.12}
\newcommand{\attSigmaObsApt}{15.2}
\newcommand{\attSigmaFullApt}{0.0}
\newcommand{\attShareZeroFullApt}{100}
\newcommand{\attPsiZeroApt}{0.13}
\newcommand{\attSdResidCoreApt}{6.0}
\newcommand{\attSdResidExtApt}{5.9}
\newcommand{\attCfRevealedApt}{23,443}
\newcommand{\attCfGainAApt}{0.00}
\newcommand{\attCfGainHiEurApt}{117}
\newcommand{\attCfGainHiEurLwApt}{74}
\newcommand{\attCfGainMidEurApt}{44}
\newcommand{\attCfCostHiTwoFiftyApt}{0.10}
\newcommand{\attCfCostHiFiveHundredApt}{8.73}
\newcommand{\attCfCostMidTwoFiftyApt}{1.21}
\newcommand{\attCfCostMidFiveHundredApt}{10.42}
\newcommand{\attCfSharePosHiTwoFiftyApt}{60}
\newcommand{\attCfNCasHiApt}{8}
\newcommand{\attCfSigmaHiApt}{20.9}
\newcommand{\attCostInfExtTwoFiftyApt}{0.41}
\newcommand{\attCostInfExtFiveHundredApt}{2.66}
\newcommand{\attRbarCoreEurTwoFiftyApt}{118}
\newcommand{\attRbarExtEurTwoFiftyApt}{122}
\newcommand{\attCostCoreTwoFiftyApt}{2.53}
\newcommand{\attCostExtTwoFiftyApt}{2.24}
\newcommand{\attCostCoreFiveHundredApt}{11.86}
\newcommand{\attCostExtFiveHundredApt}{11.57}
\newcommand{\attSharePosCoreTwoFiftyApt}{98}
\newcommand{\attSharePosExtTwoFiftyApt}{98}
\newcommand{\attOneMinusPsiYoneApt}{0.92}
\newcommand{\attOneMinusPsiYtwoApt}{0.95}
\newcommand{\attOneMinusPsiYthreeApt}{0.80}
\newcommand{\attPsiGfourtwentyApt}{0.16}
\newcommand{\attPsiGsixhundredApt}{0.09}
\newcommand{\attNViolGfourtwentyApt}{20}
\newcommand{\attNViolGsixhundredApt}{23}
\newcommand{\attScBaseRbarApt}{122}
\newcommand{\attScBaseSigmaApt}{15.2}
\newcommand{\attScBaseNCdApt}{0}
\newcommand{\attScBaseNPpApt}{45}
\newcommand{\attScBaseNCasApt}{0}
\newcommand{\attScBaseCostTwoFiftyApt}{0.12}
\newcommand{\attScBaseShareTwoFiftyApt}{98}
\newcommand{\attScBaseCostFiveHundredApt}{0.76}
\newcommand{\attScBaseShareFiveHundredApt}{100}
\newcommand{\attScSalRbarApt}{122}
\newcommand{\attScSalSigmaApt}{0.0}
\newcommand{\attScSalNCdApt}{45}
\newcommand{\attScSalNPpApt}{0}
\newcommand{\attScSalNCasApt}{0}
\newcommand{\attScSalCostTwoFiftyApt}{0.12}
\newcommand{\attScSalShareTwoFiftyApt}{98}
\newcommand{\attScSalCostFiveHundredApt}{0.76}
\newcommand{\attScSalShareFiveHundredApt}{100}
\newcommand{\attScPiOneRbarApt}{122}
\newcommand{\attScPiOneSigmaApt}{0.0}
\newcommand{\attScPiOneNCdApt}{45}
\newcommand{\attScPiOneNPpApt}{0}
\newcommand{\attScPiOneNCasApt}{0}
\newcommand{\attScPiOneCostTwoFiftyApt}{0.23}
\newcommand{\attScPiOneShareTwoFiftyApt}{98}
\newcommand{\attScPiOneCostFiveHundredApt}{1.49}
\newcommand{\attScPiOneShareFiveHundredApt}{100}
\newcommand{\attScHiRbarApt}{239}
\newcommand{\attScHiSigmaApt}{20.9}
\newcommand{\attScHiNCdApt}{33}
\newcommand{\attScHiNPpApt}{4}
\newcommand{\attScHiNCasApt}{8}
\newcommand{\attScHiCostTwoFiftyApt}{0.01}
\newcommand{\attScHiShareTwoFiftyApt}{60}
\newcommand{\attScHiCostFiveHundredApt}{0.60}
\newcommand{\attScHiShareFiveHundredApt}{100}
\newcommand{\attBaseSMedTwoFiftyApt}{127}
\newcommand{\attBaseSMaxTwoFiftyApt}{246}
\newcommand{\attBaseUnitTwoFiftyApt}{8,016}
\newcommand{\attBaseCostShareTwoFiftyApt}{53}
\newcommand{\attBaseNoSubsidyTwoFiftyApt}{2.2}
\newcommand{\attBaseSMedFiveHundredApt}{375}
\newcommand{\attBaseSMaxFiveHundredApt}{496}
\newcommand{\attBaseUnitFiveHundredApt}{23,618}
\newcommand{\attBaseCostShareFiveHundredApt}{76}
\newcommand{\attBaseNoSubsidyFiveHundredApt}{0.0}
\newcommand{\attSavingUbTwoFiftyApt}{0.077}
\newcommand{\attSavingUbDTwoFiftyApt}{0.085}
\newcommand{\attSavingUbObsTwoFiftyApt}{0.070}
\newcommand{\attSavingUbRegTwoFiftyApt}{0.044}
\newcommand{\attSavingUbFiveHundredApt}{0.077}
\newcommand{\attSavingUbDFiveHundredApt}{0.117}
\newcommand{\attSavingUbObsFiveHundredApt}{0.070}
\newcommand{\attSavingUbRegFiveHundredApt}{0.044}
\newcommand{\attHideLetterEApt}{7.0}
\newcommand{\attHideLetterFApt}{8.7}
\newcommand{\attHideLetterGApt}{11.1}
\newcommand{\attSavingTopShareApt}{83}
\newcommand{\attLamFitMeanApt}{6.1}
\newcommand{\attSigmaRegMeanApt}{5.8}
\newcommand{\attPiOneMarginApt}{4}
\newcommand{\attPsiBarApt}{0.77}
\newcommand{\attPsiBarOneApt}{0.71}
\newcommand{\attPsiBarThetaApt}{0.78}
\newcommand{\attNBetweenApt}{5}
\newcommand{\attTopMarketShareApt}{77}
\newcommand{\attRawNCdApt}{25}
\newcommand{\attRawNPpApt}{0}
\newcommand{\attRawNCasApt}{20}
\newcommand{\attRawNFlagApt}{9}
\newcommand{\attRawNNegApt}{3}
\newcommand{\attRawRbarRatioApt}{0.66}
\newcommand{\attBaseRbarMedApt}{125}
\newcommand{\attBaseRbarMinApt}{4}
\newcommand{\attBaseRbarMaxApt}{258}
\newcommand{\attBaseAreaMedApt}{63}
\newcommand{\attNListingsTotalApt}{633,399}
\newcommand{\attScBasePopApt}{44,120}
\newcommand{\attScPiOnePopApt}{86,137}
\newcommand{\attGeCorrPctBApt}{0.4}
\newcommand{\attGeCorrEurApt}{0.5}
\newcommand{\attGeDisclosedShareApt}{85}
\newcommand{\attNDgtBHouse}{33}
\newcommand{\attNMarketsHouse}{97}
\newcommand{\attExcessHouse}{\ensuremath{-}1.9}
\newcommand{\attReachHouse}{2.4}
\newcommand{\attShareUninfHouse}{98}
\newcommand{\attCursedHouse}{0.1}
\newcommand{\attLevelGapPreHouse}{\ensuremath{-}3.6}
\newcommand{\attLevelGapPostHouse}{\ensuremath{-}5.5}
\newcommand{\attEtaHouse}{2.11}
\newcommand{\attRatioHouse}{0.67}
\newcommand{\attCstarHouse}{233}
\newcommand{\attCsilHouse}{245}
\newcommand{\attPsiHouse}{0.95}
\newcommand{\attPsiLoHouse}{0.86}
\newcommand{\attPsiHiHouse}{1.04}
\newcommand{\attPsiSeHouse}{0.05}
\newcommand{\attPsiOLSHouse}{0.96}
\newcommand{\attPsiOLSseHouse}{0.04}
\newcommand{\attPsiMedHouse}{0.96}
\newcommand{\attPsiConstHouse}{0.99}
\newcommand{\attIcptHouse}{0.3}
\newcommand{\attIcptSeHouse}{0.8}
\newcommand{\attTfullHouse}{1.1}
\newcommand{\attIntervalHiHouse}{0.95}
\newcommand{\attSlackHouse}{8.37}
\newcommand{\attSlackTHouse}{13.0}
\newcommand{\attCstarLoHouse}{210}
\newcommand{\attCstarHiHouse}{256}
\newcommand{\attShareInBandHouse}{11}
\newcommand{\attShareViolBottomHouse}{22}
\newcommand{\attNViolHouse}{9}
\newcommand{\attShareViolTopHouse}{0}
\newcommand{\attBthreeHouse}{5.8}
\newcommand{\attDthreeHouse}{3.9}
\newcommand{\attLamHouse}{0.4}
\newcommand{\attLamTHouse}{1.10}
\newcommand{\attCzeroBaseHouse}{3.5}
\newcommand{\attBandLoHouse}{3.23}
\newcommand{\attBandHiHouse}{3.95}
\newcommand{\attBandWidthHouse}{0.72}
\newcommand{\attBandWidthBaseHouse}{0.75}
\newcommand{\attFittedPremiumHouse}{3.89}
\newcommand{\attSigmaOneHouse}{2.4}
\newcommand{\attSigmaZeroHouse}{2.9}
\newcommand{\attHighLossHouse}{1.9}
\newcommand{\attSalienceHouse}{0.4}
\newcommand{\attEnforceSilHouse}{0.06}
\newcommand{\attSensEHouse}{0.63}
\newcommand{\attPsiLevROMHouse}{1.01}
\newcommand{\attPsiLevSlopeHouse}{1.24}
\newcommand{\attCzeroFitHouse}{3.9}
\newcommand{\attCzeroCoreMeanHouse}{3.5}
\newcommand{\attRmseCoreHouse}{3.9}
\newcommand{\attRmseExtHouse}{3.9}
\newcommand{\attCorrCoreHouse}{0.48}
\newcommand{\attCorrExtHouse}{0.47}
\newcommand{\attResidCoreHouse}{0.4}
\newcommand{\attResidExtHouse}{0.0}
\newcommand{\attRbarGapHouse}{0.18}
\newcommand{\attSigmaObsHouse}{6.8}
\newcommand{\attSigmaFullHouse}{0.0}
\newcommand{\attShareZeroFullHouse}{100}
\newcommand{\attPsiZeroHouse}{0.96}
\newcommand{\attSdResidCoreHouse}{3.9}
\newcommand{\attSdResidExtHouse}{3.9}
\newcommand{\attCfRevealedHouse}{5,275}
\newcommand{\attCfGainAHouse}{0.00}
\newcommand{\attCfGainHiEurHouse}{100}
\newcommand{\attCfGainHiEurLwHouse}{96}
\newcommand{\attCfGainMidEurHouse}{43}
\newcommand{\attCfCostHiTwoFiftyHouse}{1.96}
\newcommand{\attCfCostHiFiveHundredHouse}{15.60}
\newcommand{\attCfCostMidTwoFiftyHouse}{4.65}
\newcommand{\attCfCostMidFiveHundredHouse}{18.44}
\newcommand{\attCfSharePosHiTwoFiftyHouse}{67}
\newcommand{\attCfNCasHiHouse}{80}
\newcommand{\attCfSigmaHiHouse}{86.2}
\newcommand{\attCostInfExtTwoFiftyHouse}{1.15}
\newcommand{\attCostInfExtFiveHundredHouse}{3.34}
\newcommand{\attRbarCoreEurTwoFiftyHouse}{124}
\newcommand{\attRbarExtEurTwoFiftyHouse}{128}
\newcommand{\attCostCoreTwoFiftyHouse}{7.31}
\newcommand{\attCostExtTwoFiftyHouse}{7.14}
\newcommand{\attCostCoreFiveHundredHouse}{21.18}
\newcommand{\attCostExtFiveHundredHouse}{20.98}
\newcommand{\attSharePosCoreTwoFiftyHouse}{91}
\newcommand{\attSharePosExtTwoFiftyHouse}{90}
\newcommand{\attOneMinusPsiYoneHouse}{0.15}
\newcommand{\attOneMinusPsiYtwoHouse}{0.03}
\newcommand{\attOneMinusPsiYthreeHouse}{0.02}
\newcommand{\attPsiGfourtwentyHouse}{0.95}
\newcommand{\attPsiGsixhundredHouse}{0.95}
\newcommand{\attNViolGfourtwentyHouse}{8}
\newcommand{\attNViolGsixhundredHouse}{10}
\newcommand{\attScBaseRbarHouse}{128}
\newcommand{\attScBaseSigmaHouse}{6.8}
\newcommand{\attScBaseNCdHouse}{0}
\newcommand{\attScBaseNPpHouse}{97}
\newcommand{\attScBaseNCasHouse}{0}
\newcommand{\attScBaseCostTwoFiftyHouse}{0.33}
\newcommand{\attScBaseShareTwoFiftyHouse}{90}
\newcommand{\attScBaseCostFiveHundredHouse}{0.95}
\newcommand{\attScBaseShareFiveHundredHouse}{100}
\newcommand{\attScSalRbarHouse}{128}
\newcommand{\attScSalSigmaHouse}{0.0}
\newcommand{\attScSalNCdHouse}{97}
\newcommand{\attScSalNPpHouse}{0}
\newcommand{\attScSalNCasHouse}{0}
\newcommand{\attScSalCostTwoFiftyHouse}{0.33}
\newcommand{\attScSalShareTwoFiftyHouse}{90}
\newcommand{\attScSalCostFiveHundredHouse}{0.95}
\newcommand{\attScSalShareFiveHundredHouse}{100}
\newcommand{\attScPiOneRbarHouse}{128}
\newcommand{\attScPiOneSigmaHouse}{0.0}
\newcommand{\attScPiOneNCdHouse}{97}
\newcommand{\attScPiOneNPpHouse}{0}
\newcommand{\attScPiOneNCasHouse}{0}
\newcommand{\attScPiOneCostTwoFiftyHouse}{0.82}
\newcommand{\attScPiOneShareTwoFiftyHouse}{90}
\newcommand{\attScPiOneCostFiveHundredHouse}{2.37}
\newcommand{\attScPiOneShareFiveHundredHouse}{100}
\newcommand{\attScHiRbarHouse}{227}
\newcommand{\attScHiSigmaHouse}{86.2}
\newcommand{\attScHiNCdHouse}{9}
\newcommand{\attScHiNPpHouse}{8}
\newcommand{\attScHiNCasHouse}{80}
\newcommand{\attScHiCostTwoFiftyHouse}{0.08}
\newcommand{\attScHiShareTwoFiftyHouse}{67}
\newcommand{\attScHiCostFiveHundredHouse}{0.70}
\newcommand{\attScHiShareFiveHundredHouse}{100}
\newcommand{\attBaseSMedTwoFiftyHouse}{145}
\newcommand{\attBaseSMaxTwoFiftyHouse}{248}
\newcommand{\attBaseUnitTwoFiftyHouse}{16,510}
\newcommand{\attBaseCostShareTwoFiftyHouse}{56}
\newcommand{\attBaseNoSubsidyTwoFiftyHouse}{10.3}
\newcommand{\attBaseSMedFiveHundredHouse}{388}
\newcommand{\attBaseSMaxFiveHundredHouse}{498}
\newcommand{\attBaseUnitFiveHundredHouse}{44,252}
\newcommand{\attBaseCostShareFiveHundredHouse}{74}
\newcommand{\attBaseNoSubsidyFiveHundredHouse}{0.0}
\newcommand{\attSavingUbTwoFiftyHouse}{0.002}
\newcommand{\attSavingUbDTwoFiftyHouse}{0.002}
\newcommand{\attSavingUbObsTwoFiftyHouse}{0.004}
\newcommand{\attSavingUbRegTwoFiftyHouse}{0.001}
\newcommand{\attSavingUbFiveHundredHouse}{0.002}
\newcommand{\attSavingUbDFiveHundredHouse}{0.002}
\newcommand{\attSavingUbObsFiveHundredHouse}{0.004}
\newcommand{\attSavingUbRegFiveHundredHouse}{0.001}
\newcommand{\attHideLetterEHouse}{2.4}
\newcommand{\attHideLetterFHouse}{3.1}
\newcommand{\attHideLetterGHouse}{4.6}
\newcommand{\attSavingTopShareHouse}{28}
\newcommand{\attLamFitMeanHouse}{0.4}
\newcommand{\attSigmaRegMeanHouse}{2.7}
\newcommand{\attPiOneMarginHouse}{84}
\newcommand{\attPsiBarHouse}{0.73}
\newcommand{\attPsiBarOneHouse}{0.68}
\newcommand{\attPsiBarThetaHouse}{0.74}
\newcommand{\attNBetweenHouse}{43}
\newcommand{\attTopMarketShareHouse}{37}
\newcommand{\attRawNCdHouse}{45}
\newcommand{\attRawNPpHouse}{11}
\newcommand{\attRawNCasHouse}{41}
\newcommand{\attRawNFlagHouse}{12}
\newcommand{\attRawNNegHouse}{21}
\newcommand{\attRawRbarRatioHouse}{0.76}
\newcommand{\attBaseRbarMedHouse}{112}
\newcommand{\attBaseRbarMinHouse}{2}
\newcommand{\attBaseRbarMaxHouse}{326}
\newcommand{\attBaseAreaMedHouse}{114}
\newcommand{\attNListingsTotalHouse}{500,415}
\newcommand{\attScBasePopHouse}{22,575}
\newcommand{\attScPiOnePopHouse}{55,907}
\newcommand{\attGeCorrPctBHouse}{0.9}
\newcommand{\attGeCorrEurHouse}{1.3}
\newcommand{\attGeDisclosedShareHouse}{93}
\newcommand{\attDiffT}{6.1}
\newcommand{\attDiffCstarT}{6.1}
\newcommand{\attSensRatio}{8}
\newcommand{\attCzeroRatio}{3.5}
\newcommand{\attNBoot}{4,999}
\newcommand{\attSampleYears}{3.5}
\newcommand{\attBudgetLo}{4}
\newcommand{\attBudgetHi}{5}
\newcommand{\attSavingUbTotTwoFifty}{0.079}
\newcommand{\attSavingUbTotTwoFiftyLo}{1.6}
\newcommand{\attSavingUbTotTwoFiftyHi}{2.0}
\newcommand{\attSavingUbTotFiveHundred}{0.079}
\newcommand{\attSavingUbTotFiveHundredLo}{1.6}
\newcommand{\attSavingUbTotFiveHundredHi}{2.0}
\newcommand{\attScCostBaseTotTwoFifty}{0.44}
\newcommand{\attScCostBaseTotFiveHundred}{1.72}
\newcommand{\attScCostBaseTotTwoFiftyLo}{9}
\newcommand{\attScCostBaseTotFiveHundredLo}{34}
\newcommand{\attScCostBaseTotTwoFiftyHi}{11}
\newcommand{\attScCostBaseTotFiveHundredHi}{43}

% Generated by scripts/mc_attention_estimator.py -- do not edit by hand.
\newcommand{\attMCtrue}{0.15}
\newcommand{\attMCmarkets}{81}
\newcommand{\attMCshareLo}{25}
\newcommand{\attMCshareHi}{75}
\newcommand{\attMColsLo}{0.33}
\newcommand{\attMColsHi}{0.70}
\newcommand{\attMCromLo}{0.10}
\newcommand{\attMCromHi}{0.14}
\newcommand{\attMColsSelLo}{0.30}
\newcommand{\attMColsSelHi}{0.60}
\newcommand{\attMCromSelLo}{0.27}
\newcommand{\attMCromSelHi}{0.49}


\begin{document}

% ============================================================
% TITLE
% ============================================================
\leftalignedtitle{When Silence Pays: Mandatory Disclosure and Strategic Concealment in
  Housing Markets
}

\leftalignedauthor{
    Bastien Patras, \href{mailto:bastien.patras@sciencespo.fr}{bastien.patras@sciencespo.fr}$^\dagger$\\
    {\small $^\dagger$PhD Student, Sciences Po Department of Economics, Paris, France}\\
}
\vspace{1em}

\abs{
\textbf{Abstract}\\[0.5em]
Standard unraveling theory predicts that mandatory disclosure separates high-quality from low-quality products: silence signals bad news and non-disclosers should trade at-or-below the worst disclosed category. We document the opposite. Following France's 2021 enforcement of Energy Performance Certificate disclosure, properties listed without a certified score trade at a premium of 6.9 percent above disclosed low-efficiency properties for apartments and 3.4 percent for houses. The puzzle resolves because non-disclosure pools informed sellers strategically concealing unfavorable scores with uninformed sellers who genuinely lack a certified audit. Unable to distinguish between them, buyers treat silence as ambiguous rather than damning. We build and structurally estimate a model of the disclosure game, bringing it to the data through a variable unobservable from the buyer's side: private listing data from SeLoger, France's leading residential platform, matched to the ADEME registry of certified energy audits, give us direct observation of each seller's information status at listing time. We derive a closed-form expression for the non-disclosure premium that, in the partial-pooling equilibrium, collapses in theory to the structural concealment cost, invariant to market composition; we test this restriction directly and find it holds for houses but is rejected for apartments, a divergence we trace to a measurable over-identifying wedge, reported directly rather than resolved by construction. An extension in which buyers attend less to energy on listings that show no rating reconciles the two segments with one attention parameter per segment, rejecting full attention for apartments but not for houses. Attention changes who hides, not who renovates: raising attention to silence leaves the minimum renovation subsidy unchanged. % [AUTHOR REWRITE]
 We extend the same framework to endogenous renovation, deriving a local minimum incentive-compatible subsidy and showing that renovation subsidies, enforcement against hiding, and audit-infrastructure investment act on distinct margins. We conclude by showing that the required cost-share varies widely across French housing markets, so that a uniform national programme overpays some markets and underpays others.
}

\vspace{1em}
\keywords{Energy Performance Certificates, Mandatory Disclosure, Housing Markets, Unraveling, Information Asymmetry} \\
\vspace{0.2em}
\jel{Q48, R31, Q52, D83, R38, H23}

% =====================================================================
% §1 INTRODUCTION — reframed draft for real-estate/finance audiences
% Follows the 14-paragraph funnel in
%   skills/realestate-finance-writing/references/introduction.md
% Companion change log: 01-introduction-changes.md
% =====================================================================

\section{Introduction}
\label{sec:intro}

% ---------- ¶1  puzzle, then the theoretical benchmark ----------
Mandatory disclosure is the standard regulatory response to asymmetric information about
product quality. Yet the case for mandating disclosure is contested, because theory suggests
markets may not need it. When quality is verifiable, the unraveling result of
\citet{grossman_informational_1981} and \citet{milgrom_good_1981} holds that sellers disclose
voluntarily, since buyers rationally read silence as the worst possible news. All but the very
worst sellers therefore disclose, and disclosure becomes effectively universal without a
mandate. Whether this result describes real markets remains unsettled. It rests on strong
assumptions, and it is difficult to test, because in most settings where quality is verifiable
disclosure has long been compulsory.

% ---------- ¶2  "to learn about X, we examine Y" ----------
To examine whether unraveling describes a market where quality is verifiable but silence is
common, we study the French residential housing market around the July~2021 enforcement of
Energy Performance Certificate (EPC) disclosure. Under a 2018 reform that took effect on
July~1, 2021, a property's \emph{diagnostic de performance \'energ\'etique} (DPE) became
legally opposable.\footnote{Article 179 III of \emph{loi n\textdegree{}2018-1021 du 23 novembre
2018} (loi \'ELAN) amended Articles L271-4 and L134-3-1 of the \emph{Code de la construction et
de l'habitation}, with opposability effective July~1, 2021; the accompanying methodology
overhaul (replacing the earlier invoice-based calculation with the current
physical-characteristics-based method) was set by the \emph{Arr\^et\'e du 31 mars 2021}. The
\emph{loi Climat et R\'esilience} (\emph{loi n\textdegree{}2021-1104 du 22 ao\^ut 2021}) is a
separate, later reform; we invoke it elsewhere in the paper for the F/G rental-ban calendar it
introduced, not for DPE opposability.} A seller who advertises an energy score is now bound by
it, and buyers acquired legal recourse against misstatement. The reform changed what a disclosed
score means without changing the underlying housing stock or the legal consequences of remaining
silent. It therefore isolates the demand-side value of the signal, and it does so in a market
where the certificate is issued and archived by a national agency rather than by the seller or
the listing platform.

% ---------- ¶3  what theory predicts here, and why this market adjudicates ----------
Applied to this setting, unraveling makes a sharp prediction. Once the DPE is credible, sellers
of energy-efficient properties should publish their score to separate themselves; sellers of
mediocre properties should then publish to avoid being pooled with worse ones; and the process
should continue until listings without a score are confined to the very bottom of the quality
distribution, priced at or below the worst disclosed grade. French housing is a suitable
testbed for this prediction because the disclosed attribute is a single certified index with a
clear mapping into a running cost buyers face directly, because the certificate is verifiable
against an independent public registry, and because a large stock of properties entered the
post-reform market without one.

% ---------- ¶4  the assumption audit ----------
The unraveling prediction rests on four assumptions: (a)~buyers know sellers possess private
information about quality, (b)~disclosure is costless, (c)~disclosure is credible, and
(d)~buyers are homogeneous in how much they value the disclosed attribute. Three of the four are
either directly satisfied in this setting or can be dispensed with. Assumption~(c) holds by
construction after July~2021, since the DPE is issued by a certified third-party diagnostician,
archived in the registry of the \emph{Agence de la Transition \'Ecologique} (ADEME), and legally
binding on the seller. Assumption~(b) is close to satisfied, because a DPE is a one-off outlay of
roughly 100--250~\euro{} according to ADEME's own estimates, negligible relative to
residential transaction values, and it is the seller's obligation rather than the buyer's.
Assumption~(a), which the theoretical literature treats as a maintained hypothesis, we do not
need to assume at all. By matching listings to the ADEME registry we observe directly which
sellers hold a certificate at the moment of listing, and among the apartment listings we can
match at the address level, 48.8\% hold none.
Silence in this market is therefore not always a choice, a possibility first identified by
\citet{dye_disclosure_1985} and formalised by \citet{jung_kwon_1988}, and here it is measured
rather than inferred.

% ---------- ¶5  isolate the assumption at stake ----------
The descriptiveness of unraveling in this market accordingly hinges on assumption~(d).
Unlike disclosure costs and disclosure credibility, which a large literature has examined,
the homogeneity of buyer valuations has received little attention, even though it sets the size of
the disclosure incentive. A seller's penalty for withholding is proportional to the share of
buyers who price the withheld attribute. Housing makes the
assumption particularly implausible. A buyer of a rarely-occupied second home in a mild climate
may rationally ignore an energy score that a buyer of a primary residence in a cold region
prices directly, as may an investor facing the impending rental prohibition on the worst-rated
properties. Evidence across many regulated markets shows that disclosure interventions fail
precisely when recipients are inattentive or heterogeneous in their response
\citep{loewenstein_disclosure_2014, benshahar_failure_2011}, and
\citet{fishman_mandatory_2003} show formally that mandated disclosure can be ineffective, or
welfare-reducing, when buyers differ in their capacity to use the disclosed information. We
therefore let the share of buyers who price the EPC vary across local markets and estimate it.

% ---------- ¶6  data and vantage point  [PRESERVED from current ¶6] ----------
Testing a theory of strategic concealment requires observing both sides of the information
asymmetry at once, the seller's true information, the seller's disclosure choice and the
market's response. Most empirical work observes only one or two. Studies of information
asymmetry in housing infer agent advantages from equilibrium price deviations
\citep{levitt_syverson_2008}; analyses of certification intermediaries characterise incentives
without observing individual information endowments \citep{lizzeri_1999}; and work on the
energy-efficiency gap proxies for seller awareness through indirect variation in energy prices
\citep{myers_asymmetric_2020}. We observe all three. Listing records from SeLoger, a major
French residential platform, give the listing price, property characteristics, the seller's
disclosure decision, and, when a score is displayed, the letter itself. The registry of the
\emph{Agence de la Transition \'Ecologique} (ADEME) is the official national database of
certified DPE audits, reported by accredited third-party assessors and compiled independently
of the listing market. We match listings to the registry on address, property type, area, and
floor, which identifies whether a certified audit existed at the time of listing, the one
thing buyers can never observe. Because the match requires a geocoded address, the matched
sample is the subset of listings that can be located at the address level rather than the full
listing universe. Within that matched sample, every listing falls into one of four groups, namely
truthful disclosers, strategic hiders who hold a certificate and omit it, genuinely uninformed
sellers for whom no certificate exists, and a small residual displaying a score without a
verified audit. The information structure of the disclosure game is therefore observed rather
than inferred, a vantage point that to our knowledge is unique in this literature. Transaction
records enter separately and are not matched to individual listings. We use them to estimate
the mapping from listing to realised sale price, which supports treating the listing price as
the relevant outcome, since the negotiation margin in French residential sales is small and
most properties transact at or near list price
\citep{merlo_bargaining_2004, genesove_loss_2001, friggit_marche_2013}.

% ---------- ¶7  descriptive findings, enumerated ----------
The data yield four descriptive findings, drawn from two distinct samples that we keep separate
throughout. First, silence commands a premium rather than a penalty. Estimated on the panel of
market-level hedonic price indices, and relative to the pre-reform baseline, listings without a
score trade approximately 6.9\% above disclosed energy-intensive properties for apartments and
3.4\% for houses, the opposite of the sign unraveling predicts. Second, the premium widened
abruptly at the reform rather than emerging from trend. The composition-adjusted
difference-in-differences estimates evolve in parallel before July~2021, averaging $+0.002$ for
apartments and $+0.008$ for houses with no discernible drift, and then shift sharply upward,
stabilising around $+0.044$ and $+0.037$ by mid-2024. For apartments the level ordering already
favoured non-disclosers beforehand, so what the reform changes is the magnitude of the premium
rather than its sign. Third, turning to the address-matched listing sample, non-disclosure is
overwhelmingly involuntary. Among apartment listings 79.8\% disclose truthfully, 16.7\% are
genuinely uninformed, and only 3.5\% hold a certificate and withhold it, with the corresponding
house shares 90.4\%, 8.5\% and 1.1\%.
Fourth, the non-disclosure rate varies widely across local markets, ranging among
apartments from 6.5\% to 40.9\% across d\'epartements, around a median of 16.5\%. Together
these facts are inconsistent with unraveling as a description of this market, while pointing to
the mechanism that defeats it. The pool of silent sellers is dominated by the uninformed, which
blunts the inference buyers can draw from silence and shelters the few who conceal
deliberately.

% ---------- ¶8  model motivation ----------
To recover the forces behind these patterns, we build and estimate a model of the disclosure
game that allows for, but does not impose, strategic concealment. Sellers differ in whether they
hold a certified audit and in the quality it would reveal; buyers differ in whether they price
energy performance. Informed sellers of poor properties choose between publishing a bad score
and accepting the price buyers assign to silence, and that choice depends on how contaminated
the silent pool is by sellers who simply have no certificate. Because the composition of the
silent pool is itself the outcome of those choices, the model is closed by a fixed point in
buyer beliefs, and it nests the classical benchmark. When audit coverage and buyer sensitivity
are both high, the silent pool empties and the Grossman--Milgrom outcome is recovered. For
intermediate values, which is where the French data sit, the model predicts partial pooling in
which some informed sellers conceal and silence carries a premium.

% ---------- ¶9  what the model finds ----------
The estimates deliver two structural magnitudes and one specification test. Identifying buyer
energy-sensitivity separately from the concealment cost requires a source of variation external
to prices, since the two are shown to be identified only as their product by price moments
alone (Appendix~\ref{app:rho_identification}); we construct one from the capitalised energy-cost
saving a sensitive buyer attaches to verified high- versus low-quality certification, built from
ADEME-certified consumption bands rather than from the price data itself. The resulting share of
buyers who price energy performance averages 49.5\% across apartment markets and 45.8\% across
house markets, and the cost of withholding a certificate settles at 2.4\% of property value for
apartments and 3.5\% for houses, on the order of a third of the notary bill a French buyer
separately pays on the same transaction. Both are economically substantial magnitudes, and of the
same order in two segments that trade very differently. The model predicts, in the partial-pooling
equilibrium that describes the bulk of French housing markets, that the non-disclosure premium
should collapse to this concealment cost exactly, invariant to market composition; we test this
prediction directly rather than impose it. It holds for houses: the gap between the observed
premium and the identified concealment cost is small and statistically indistinguishable from
zero. It does not hold for apartments: the premium exceeds the concealment cost by a wide and
precisely estimated margin, a specification tension the baseline cannot absorb, because no belief
about silence can price a silent listing above a verified efficient one.
% [AUTHOR REWRITE]
Section~\ref{ss:attention} shows that one change on the buyer side can account for it: if buyers attend
to energy when a rating is displayed and much less when it is not, the model fits houses as before,
rejects full attention in apartments, and leaves the estimated buyer sensitivity unchanged.
We maintain, rather than test, the model's prediction that informed sellers of high-quality
properties never conceal; Appendix~\ref{app:robustness} relaxes it by reading the withheld grade
from the registry and finds the structural estimates essentially unchanged. Under the limited-attention
extension the prediction survives whenever not every informed low-quality seller hides (Corollary~\ref{cor:att_hq}).

% ---------- ¶11  contribution 1 ----------
Our paper contributes to the literature on disclosure by testing the unraveling result in a
market where the informed and uninformed can be told apart. Prior work has examined the
consequences of uncertainty about information endowment \citep{dye_disclosure_1985,
jung_kwon_1988}, of disclosure costs, and of disclosure credibility, but has generally had to
treat sellers' information as latent. By observing it directly we can decompose silence into
inability and choice, and show that the first dominates the second by an order of magnitude
while the second is what sustains the price premium. The finding that silence is rewarded rather
than punished is not a failure of the theory so much as a demonstration of how sensitive its
central prediction is to an assumption that is rarely examined.

% ---------- ¶12  contribution 2 + scope disclaimer ----------
Our paper also provides an information-based account of the residential energy-efficiency gap
\citep{allcott_is_2012, gerarden_assessing_2017}. The standard diagnosis appeals to split
incentives, credit constraints, and inattention to future energy costs \citep{allcott_consumers_2011, sallee_rational_2014, allcott_evaluating_2015}. We identify a
complementary channel internal to the disclosure game. A seller holding a poor rating can
renovate and publish a good one, earning the quality premium, or stay silent and earn the
non-disclosure premium instead, so the premium to opacity looks like an opportunity cost of
renovation that raises the private hurdle above engineering cost. In equilibrium it does not, except in the corner where every informed low-quality seller hides. This has an unexpected policy
implication. Because the equilibrium absorbs enforcement through the fraction of sellers who
conceal, the renovation threshold depends only on the price gap between disclosed quality
levels and is independent of enforcement intensity. Renovation subsidies, enforcement against
concealment, and investment in audit coverage act on three distinct margins and do not substitute for one another,
contrary to the conventional treatment in
\citet{gillingham_palmer_2014}. The result holds outside the corner in which every informed low-quality seller hides, and with one concealment cost per market. % [AUTHOR REWRITE] The extension of Section~\ref{ss:attention} brings the inattention
diagnosis inside the disclosure game: buyers who attend less to energy when no rating is shown let
silence earn a premium that no belief about quality can produce. We do not claim to have estimated the one true model of housing
disclosure; we propose one way of reading these data through a workhorse framework, and we are
explicit throughout about which parameters are point-identified and which are bounded.

% ---------- ¶12.5  policy extension results ----------
The model also speaks to renovation policy. Extending the framework to an endogenous upgrade
decision, we show the minimum subsidy needed to make renovation privately worthwhile reduces, in
the empirically dominant partial-pooling regime, to the price gap between verified high and low
quality alone, invariant to enforcement against hiding and to how many sellers are informed,
which act on separate margins instead. Forcing silent listings to be read as carefully as rated ones would change who hides but not who renovates, and only a higher weight on displayed ratings could raise the renovation threshold, within a bound that the data do not pin down. % [AUTHOR REWRITE]
A uniform national subsidy is accordingly poorly targeted. The
required cost-share the state must cover varies widely across
markets, averaging \attBaseCostShareTwoFiftyApt{} percent in apartment markets that need support at all at the light benchmark, and aggregate programme costs among the sellers who can renovate run to
\euro{}\attScCostBaseTotTwoFifty{}--\attScCostBaseTotFiveHundred{} billion a year in our listing sample at plausible renovation-cost benchmarks, for scale between \attScCostBaseTotTwoFiftyLo{} and
\attScCostBaseTotFiveHundredHi{} percent of MaPrimeR\'enov's current annual budget. Efficient targeting requires conditioning subsidy design
on local market composition rather than applying one national rate. Enforcement against strategic hiding and investment in audit
infrastructure act on separate margins: the first changes who hides, the second who is eligible to renovate, and neither lowers the subsidy a given seller needs.

% ---------- ¶13  related literature, three strands ----------
Our paper draws on three strands of work. It belongs to the literature that estimates
disclosure models structurally in order to assess the descriptiveness of classical disclosure
theory and to recover primitives such as information endowment, disclosure costs, and beliefs;
closest to us is \citet{bourveau_learning_2025}, who structurally examine whether the unraveling
result describes an unregulated nineteenth-century market and find that full disclosure is
reached only gradually. Where that literature and the experimental work on disclosure ask how far
receivers infer the worst from silence \citep{jin_is_2021, hagenbach_motivated_2025}, our extension
asks whether they attend to silence at all, drawing on evidence that limited attention changes how
disclosed and undisclosed information is priced \citep{hirshleifer_limited_2003, dellavigna_investor_2009}.
It is closely related to the extensive evidence on the capitalisation of energy labels in
housing, which documents robust premia for efficient properties wherever labels have been
studied \citep{brounen_kok_2011, hyland_lyons_2013, fuerst_energy_2016, frondel_power_2018,
myers_mandatory_2022}; that literature establishes that labels are priced, whereas our question
is why their \emph{absence} is priced favourably. And it builds on work on information
asymmetry in real-estate markets, where the informational advantage of one side of the
transaction has been inferred from prices and time-on-market \citep{levitt_syverson_2008,
lizzeri_1999}.

% ---------- ¶14  generalizability, conceded then defended ----------
The features that make this setting informative also bound it. French EPC disclosure is a
single mandate, in one country, covering one verifiable attribute, and the low audit coverage
that lets us separate inability from choice is itself a feature of a young disclosure regime.
The mechanism, however, is not specific to energy labels. Wherever a disclosure mandate is
introduced before the infrastructure that makes disclosure feasible is universal, a pool of
genuinely uninformed non-disclosers will exist, and that pool will shelter strategic
concealment and dull the incentive to disclose. That configuration characterises a broad class
of emerging disclosure regimes, from climate-risk and emissions reporting to supply-chain and
sustainability disclosure, in which mandates routinely precede measurement capacity, including
the French carbon-disclosure mandate studied by \citet{bourveau_carbon_2024}, where the same
gap between obligation and measurement infrastructure arises.


% =====================================================================
% §2 CONCEPTUAL SECTION — new material
% Pattern: economic problem -> theory in prose, no math -> assumptions
%          enumerated -> closes on the empirical question
% See skills/realestate-finance-writing/references/structure.md
%
% WIRING DEPENDENCY: this file contains fig:equilibrium_outcomes, moved
% here from §3.1. Deleting the float from §3.1 is required in the same
% commit that inputs this file, or the label will collide.
% =====================================================================

\section{Unraveling with Possibly Uninformed Sellers}
\label{sec:unraveling}

\subsection{The Unraveling Benchmark}

Asymmetric information about product quality can prevent trade from happening at the right
prices. If buyers cannot verify quality before purchasing, they bid no more than the value of
what they expect to receive on average. Sellers of genuinely good products are then underpaid,
and may withdraw from the market altogether, leaving buyers facing a pool of worse quality than
they had hoped for. The outcome is bad for both sides: buyers pay for uncertainty they would
rather not bear, and the best sellers cannot capture the value of what they own.

Disclosure appears to resolve this problem without any regulatory intervention. Suppose quality
is verifiable, so that a seller who makes a claim about it can be held to that claim. A seller
of the highest-quality product will disclose, because staying silent means being valued at the
average of the pool rather than at what the product is actually worth. Buyers, anticipating
this, understand that any seller who remains silent is not of the highest quality, and revise
downward what they will pay for silence. That revision gives the second-highest-quality seller
the same incentive, and disclosing becomes worthwhile for them too. The argument repeats. Each
round of disclosure worsens the expected quality of the remaining silent pool, which strengthens
the incentive for the best of those who remain to separate themselves. The process continues
until all sellers but possibly the very worst have disclosed. This is the unraveling result of
\citet{grossman_informational_1981} and \citet{milgrom_good_1981}: even sellers with bad news
disclose, because they anticipate that silence will be read as worse news still. If it holds,
market forces alone deliver near-complete disclosure and a mandate is redundant.

The argument is powerful but it is not free. It rests on four assumptions. First, (a)~buyers know
that sellers possess private information about quality; if a seller may simply not know, silence
no longer implies unwillingness to reveal. Second, (b)~disclosure is costless, or at least cheap
relative to the price consequences of silence; a large literature examines how disclosure costs
truncate the unraveling process and leave a residual pool of non-disclosers
\citep{verrecchia_1983, jovanovic_1982}. Third, (c)~disclosure is credible, so that a claim about
quality cannot be fabricated; without verifiability, disclosure conveys nothing and the
unraveling logic never starts. Fourth, (d)~buyers are homogeneous in how much they value the
disclosed attribute, so that the price penalty for silence is felt by every seller with equal
force.

The first three assumptions have received sustained attention. The fourth has not, even though
it governs the strength of the entire mechanism: what drives a seller to disclose is the revenue
lost when buyers price silence pessimistically, and that loss is proportional to the share of
buyers who care about the attribute in the first place. Whether unraveling describes any
particular market is therefore an empirical question, and answering it requires knowing which of
these assumptions fails there.

\subsection{Two Departures and Why They Compound}

Two of the four assumptions are implausible in markets where quality is certified by a
third party and buyers differ in their exposure to the certified attribute.

The first departure concerns assumption~(a). When some sellers genuinely lack private
information, silence pools those who \emph{will not} disclose with those who \emph{cannot}.
\citet{dye_disclosure_1985} showed that this possibility breaks the buyer inference on which
unraveling depends: a rational buyer confronted with silence must now form a posterior over two
distinct populations, and how pessimistic that posterior should be depends on their relative
masses rather than on the act of silence itself. \citet{jung_kwon_1988} characterised the
resulting equilibria and showed that partial pooling, in which some informed sellers of poor
quality withhold while the uninformed are silent by necessity, is generically stable. Disclosure
does not unravel to the bottom; it stops somewhere in the interior.

The second departure concerns assumption~(d). If only a fraction of buyers price the disclosed
attribute, a seller who withholds it faces a diluted penalty, because part of the demand side is
indifferent to what was withheld. \citet{fishman_mandatory_2003} show formally that mandated
disclosure can be ineffective, and even welfare-reducing, when buyers differ in their capacity to
use disclosed information. The disclosure incentive is thus scaled by the composition of demand,
and where that composition is unfavourable the incentive can be too weak to sustain separation
even among sellers who are perfectly well informed.

The two departures do not merely coexist; they compound. A larger mass of sellers who cannot
disclose dilutes the adverse inference buyers draw from silence, which raises the expected quality
of the silent pool and makes withholding more attractive to informed sellers of poor quality.
Those sellers then join the pool, and while their entry worsens its true average quality, buyers
cannot separate them from the involuntarily silent. Meanwhile, a demand side in which many buyers
do not price the attribute weakens the penalty that would otherwise discipline them. Put together,
the consequence is stronger than persistence of non-disclosure: silence can command a
\emph{premium} over the worst disclosed quality, because the silent pool contains a substantial
mass of sellers whose silence carries no adverse information at all. The sellers who cannot
disclose do not intend to shelter anyone. Their presence is involuntary. But it is precisely what
makes silence ambiguous, and ambiguity is what makes concealment privately rational.

Because the two departures operate on different primitives, they map out a space of possible
market outcomes rather than a single prediction. Figure~\ref{fig:equilibrium_outcomes} organises
that space along the two dimensions: the share of sellers who hold verifiable information, and
the share of buyers who price it. Full unraveling requires both to be high. Where either is
interior, silence survives in equilibrium, and where the informed share is low while buyer
sensitivity is not, silence can be rewarded rather than punished.

\begin{figure}[H]
    \centering
    \includegraphics[width=\textwidth]{fig/visuals/equilibirum_outcomes.png}
    \caption{Equilibrium outcomes under seller information and buyer sensitivity heterogeneity}
    \label{fig:equilibrium_outcomes}
\end{figure}

Two features of this taxonomy matter for what follows. It implies that the informational content
of silence is a market-level object rather than a universal one, so that two markets facing the
same disclosure mandate can sit in different cells and sustain opposite price patterns. And it
implies that the two primitives are separately movable by policy: the share of sellers holding
verifiable information can be raised by investment in certification infrastructure, while the
share of buyers who price the attribute responds to the legal weight the certificate carries and
to the cost of the attribute itself. Which cell a given market occupies, however, cannot be
settled by theory. It requires measuring the composition of the silent pool and the sensitivity
of buyers separately, which is what the data described in the next section allow us to do.


% =====================================================================
% §3 SETTING AND DATA — restructured
% Corpus pattern: institutional background -> data -> stylized facts,
% all before the model. Assembled programmatically: 3.1 and 3.2's framing
% prose is new; the institutional paragraph, the data subsections, and the
% four stylized facts are relocated verbatim from elsewhere in the paper.
%
% WIRING DEPENDENCIES (all in the same commit that inputs this file):
%   1. delete the old \section{Setting and Stylized Facts} block
%   2. delete the old \section{Data} block  (its \label{sec:data} is
%      carried onto §3.2 below -- 4 live \ref{sec:data} calls depend on it)
%   3. delete the institutional paragraph from the identification section
%      ("Before July 2021, the French DPE was a paper exercise...")
% =====================================================================

\section{The French Energy-Performance Disclosure Regime}
\label{sec:setting}

This section sets out the institution, the data, and four features of the market that any
candidate theory of EPC disclosure has to accommodate. We lead with apartments throughout and
treat houses as a corroborating segment. The two display the same qualitative patterns, but the
apartment market has the lower audit coverage, the larger share of strategic concealment, and the
stability across the reform that the identification strategy exploits.

\subsection{Institutional Background}
\label{ss:institutions}

The French \emph{diagnostic de performance \'energ\'etique} (DPE) is a standardised energy audit of
a dwelling, carried out by an accredited third-party assessor and recorded in a national registry
maintained by the \emph{Agence de la Transition \'Ecologique} (ADEME). It assigns a letter grade
from A to G summarising modelled energy consumption per square metre. Two features make it
unusual as a disclosure instrument. It is certified rather than self-reported, so a displayed
grade can be checked against the registry. And the registry is centralised and public, so the
\emph{existence} of an audit is verifiable independently of both the seller and the listing
platform. That second property is what makes the empirical exercise in this paper possible.

Before July 2021, the French DPE was a paper exercise. Scores were methodologically inconsistent across assessors, carried no legal weight, and were routinely ignored in negotiations: a buyer who later discovered a disclosed rating was misleading had no legal recourse. The reform of July 1, 2021 changed the regime fundamentally. It standardized the methodology, imposed stricter accreditation requirements on assessors, and, most importantly, rendered disclosed EPC ratings legally opposable: a buyer can now use a disclosed rating as grounds for price renegotiation or legal action. The DPE became, for the first time, a document with contractual force.

Two things the reform did not change matter as much as what it did. First, the obligation to
commission a DPE before marketing a property predates it. The requirement, and the liability
attached to marketing a property without the required diagnostic, were established by \emph{loi
ALUR} (Article L271-4 of the \emph{Code de la construction et de l'habitation}) well before 2021.
The legal exposure a seller incurs by coming to market without a certificate is therefore governed
by a mandate that straddles the reform date unchanged, while what July 2021 altered is how much a
\emph{disclosed} grade is worth to the buyer. Second, the reform sharpened a pre-existing
regulatory threshold rather than inventing one. Under \emph{loi Climat et R\'esilience}, dwellings
rated E, F, or G face a phased prohibition on new lettings. This draws a bright line between the
bottom of the grade distribution and the rest, and gives the binary quality distinction used below
a concrete regulatory anchor rather than an arbitrary cutoff.

\subsection{Data and Measurement}
\label{sec:data}

Testing a theory of strategic concealment requires observing something buyers cannot, namely whether
a silent seller holds a certificate. We recover it by matching listings against the ADEME registry.

\subsubsection{Sources and Sample Construction}

The empirical analysis draws on two sources matched at the listing level, together with transaction records used separately. Listing records are sourced from SeLoger, a major French residential real-estate platform covering sales across metropolitan France. For each listing, we observe the listing price, property characteristics ($X_i$: area, number of rooms, floor, construction era, location), the disclosure indicator $d\in\{0,1\}$, and, when $d=1$, the displayed EPC letter (A through G). Transaction records are not matched to individual listings; we use them to estimate the mapping from listing price to realized sale price ($\gamma_0$ and $\gamma_1$ in equation~\eqref{eq:payoff}), which supports treating the listing price as the relevant outcome for the disclosure decision. The ADEME registry is the official national database of certified Diagnostic de Performance Énergétique (DPE) audits, maintained by France's environment and energy management agency and reported by accredited third-party assessors. We match listings to the ADEME registry using address, property type, area, and floor, assigning to each listing an indicator $K=1$ if a certified audit record exists at the time of listing, and $K=0$ otherwise. For the structural estimation, listings for which no address is available are dropped, as the ADEME match requires a geocoded address; these account for a small share of the total sample.

The full listing sample covers 2021Q3 to 2024Q4, the post-reform period during which EPC disclosure is legally binding and DPE scores are credible to buyers. The apartment sample comprises 843,621 listings and the house sample 687,878 listings. For structural estimation we restrict to 2022 to 2024: the H2~2021 period (the six months immediately following the reform) is excluded because the market had not yet fully repriced the DPE (hedonic estimates of $\hat\beta_1$ are near zero or negative in that window), yielding 767,000 apartment listings and 633,000 house listings for the structural sample. For 51.2\% of apartment listings and 40.3\% of house listings, the ADEME match yields $K=1$; the remainder are classified as genuinely uninformed ($K=0$, no verified audit on record). For identification of market-level price indices, we use the panel of hedonic price indices constructed in \citet{patras_mandatory_2024}, covering 110 markets (attraction zones, \textit{zones d'emploi}) with monthly observations from October 2018 through July 2024.

\subsubsection{Four-Group Classification}

Combining the listing disclosure decision $d$ with the ADEME information status $K$ generates four mutually exclusive groups, summarized in Table~\ref{tab:groups}.

\begin{table}[htbp]
\centering
\caption{Classification of listings and sample composition}
\label{tab:groups}
\small\renewcommand{\arraystretch}{1.3}
\begin{tabular}{llcccc}
\toprule\toprule
\textbf{Group} & \textbf{Interpretation} & $d$ & $K$ & \textbf{Apartments} & \textbf{Houses} \\
\midrule
Truthful disclosers  & Informed seller reveals $\theta$        & 1 & 1 & 47.7\% & 39.2\% \\
Strategic hiders     & Informed seller conceals $\theta$     & 0 & 1 & \textbf{3.5\%}  & \textbf{1.1\%}  \\
Genuinely uninformed & Silent; no registry record              & 0 & 0 & 16.7\% & 8.6\% \\
Disclosed, unmatched & Displays EPC; no post-2021 registry record & 1 & 0 & 32.2\% & 51.2\% \\
\midrule
Total                &                                         &   &   & 843,621 & 687,878 \\
\bottomrule\bottomrule
\end{tabular}
\vspace{0.5em}
\begin{minipage}{\linewidth}
\footnotesize\textit{Notes.} Shares of the full listing sample of each property type, computed from the
listing--registry match. The ``disclosed, unmatched'' group displays an EPC score for which no record
appears in the public ADEME registry, which contains only certificates issued from July~2021 onward;
these listings are classified as truthful disclosers under the no-falsification assumption. Sellers
with no post-2021 registry record ($K=0$) therefore account for 48.8\% of apartment and 61.1\% of
house listings, the large majority of whom do disclose a score. Columns may not sum to exactly 100\%
because of rounding.
\end{minipage}
\end{table}

\noindent Three features stand out. First, strategic hiding is rare but economically meaningful: 3.5\% of apartment listings and 1.1\% of house listings come from informed sellers who chose not to disclose. Pooling across listings, the hiding rate among informed sellers ($K=1$) is 6.9\% for apartments and 2.7\% for houses; the corresponding unweighted means across markets, reported as $\hat\sigma_{0,m}$ in the results, are 7.7\% and 3.7\%. Second, silence is uncommon relative to non-verification: only 16.7\% of apartment listings and 8.6\% of house listings are silent sellers with no registry record, and these are the non-disclosers who provide involuntary cover for strategic hiders. Third, and quantitatively dominant, 32.2\% of apartment listings and 51.2\% of house listings display a score with no matching registry record. Taken together, sellers without a post-2021 registry record account for 48.8\% of apartment listings and 61.1\% of house listings, but the large majority of them do disclose a score.

Three features of this match bear on its credibility. It is a match to an administrative registry
rather than to a survey or a self-report, so the information-status indicator does not depend on
seller cooperation and cannot be strategically misreported. It requires a geocoded address, so the
matched sample is the subset of listings that can be located at the address level rather than the
full listing universe; sample scopes are stated alongside the statistics reported below. And, most
consequentially, $K=0$ must be read as \emph{no post-July-2021 certificate in the public registry}
rather than as \emph{no certificate exists}. The registry was populated from the reform date
onward, while a certificate issued earlier could remain legally valid at the time of listing, so a
seller could hold and display a valid certificate that the registry cannot corroborate. The data
bear this out. The share of listings displaying a score with no registry record falls monotonically
as the registry accumulates, from $62.5\%$ of apartment listings in 2021Q3 to $18.4\%$ in 2024Q4,
mirrored by a rise in verified disclosers from $7.4\%$ to roughly one half.

This coverage gap has a known direction. A seller holding a still-valid pre-reform certificate who
withholds it is recorded as genuinely uninformed rather than as a strategic hider, so the measured
hiding rate $\hat\sigma_{0,m}$ understates the true one. Because the equilibrium belief
$e(\sigma_0)$ is decreasing in $\sigma_0$, an understated hiding rate implies an overstated belief
and hence an overstated concealment cost. The coverage gap therefore reinforces, rather than
threatens, the interpretation of $\hat c_{0,m}$ as an upper bound. It is also one reason the
structural sample begins in 2022, since the second half of 2021, where the gap is most severe, is
excluded.

\subsection{Stylized Facts}
\label{sec:facts}

Four features of the market discipline any candidate theory. Each is documented for both segments,
apartments first.

\subsubsection{Non-Disclosing Listings Command a Price Premium, Widened by the 2021 Reform}
\label{ss:fact1}

The starting puzzle is that silence about energy performance commands a price premium relative to the worst certified scores in both property types. Figure~\ref{fig:price_levels} plots the volume-weighted average of market-level hedonic price indices per square metre for each EPC disclosure category from January 2019 to June 2024, where each index holds property characteristics constant within a local market so that comparisons across disclosure categories reflect differences in the disclosure decision rather than differences in property type or location. For apartments (panel~(a)), non-disclosing listings occupy the highest price tier across the \emph{entire} sample window: the ordering Non-disclosed $>$ Efficient (A/B/C) $>$ Energy-intensive (E/F/G) is visible as early as 2019, well before the reform, and the gap widens sharply after July 2021. For houses (panel~(b)), the pre-reform picture differs: non-disclosers begin near the bottom of the price distribution and rise toward the middle post-reform, reflecting the higher baseline audit coverage in the house segment documented in Section~\ref{ss:fact3}. The pre-reform level premium for apartments is consistent with the concentration of non-audited stock in high-price urban markets. Figure~\ref{fig:fact1_premium} isolates the \emph{change} in the gap: it plots difference-in-differences estimates of the hedonic-controlled log price gap between non-disclosing and energy-intensive listings, purging observable composition differences and measured relative to the pre-reform baseline. Because the estimates are normalised to that baseline, the pre-reform coefficients are informative about trends rather than levels. They average $+0.002$ log points for apartments and $+0.008$ for houses and display no discernible drift, indicating that the two series evolved in parallel before July 2021. This is precisely the identifying assumption the design requires; it is not a claim that non-disclosing and energy-intensive listings were priced alike, and Figure~\ref{fig:price_levels} shows that for apartments they were not. After July 2021, when DPE certificates became legally opposable under \emph{loi Climat et R\'esilience}, the estimates shift abruptly upward and stabilize around $+0.044$ log points for apartments and $+0.037$ log points for houses by mid-2024. The reform therefore does not create the level difference so much as widen it sharply, in both the raw series and the composition-adjusted estimates. Classical unraveling predicts the opposite sign: silence should signal bad news and trade \emph{below} verified low-quality listings.

\begin{figure}[tbp]
    \centering
    \caption{Price per m\textsuperscript{2} by EPC disclosure category}
    \label{fig:price_levels}
    \begin{subfigure}{0.49\textwidth}
        \centering
        \caption{Apartments}
        \includegraphics[width=\textwidth]{fig/stylized_facts/apartment/fig_price_levels.pdf}
    \end{subfigure}
    \begin{subfigure}{0.49\textwidth}
        \centering
        \caption{Houses}
        \includegraphics[width=\textwidth]{fig/stylized_facts/house/fig_price_levels.pdf}
    \end{subfigure}
    \begin{minipage}{0.95\textwidth}
        \footnotesize\textit{Notes.} Volume-weighted national average of market-level hedonic price indices per square metre (\euro) by EPC disclosure category, January 2019 to June 2024. Price indices are constructed in \citet{patras_mandatory_2024} using hedonic regressions estimated separately within each \emph{zone d'emploi} (INSEE clusters of municipalities), covering 110 local markets. National series are obtained by aggregating market-level indices by listing volume. Series are smoothed with a centred three-month moving average. The vertical dashed line marks the July 2021 DPE reform. Source: \citet{patras_mandatory_2024}; authors' calculations.
    \end{minipage}
\end{figure}

\begin{figure}[tbp]
    \centering
    \caption{Non-disclosure price premium around the July 2021 DPE reform}
    \label{fig:fact1_premium}
    \begin{subfigure}{0.49\textwidth}
        \centering
        \caption{Apartments}
        \includegraphics[width=\textwidth]{fig/stylized_facts/apartment/fig_fact1_nondiscl_premium_eventstudy.pdf}
    \end{subfigure}
    \begin{subfigure}{0.49\textwidth}
        \centering
        \caption{Houses}
        \includegraphics[width=\textwidth]{fig/stylized_facts/house/fig_fact1_nondiscl_premium_eventstudy.pdf}
    \end{subfigure}
    \begin{minipage}{0.95\textwidth}
        \footnotesize\textit{Notes.} Monthly volume-weighted log gap between the price index of non-disclosing listings and certified energy-intensive (E/F/G) listings, January 2019 to June 2024. The vertical dashed line marks the legal opposability of the DPE (July 2021). The thick line is a centred three-month moving average. Source: SeLoger residential listings matched against the ADEME registry of certified energy audits; authors' calculations.
    \end{minipage}
\end{figure}

\subsubsection{Non-Disclosure Persists Long After Mandatory Disclosure}
\label{ss:fact2}

Classical unraveling further predicts that, once disclosure is mandatory and verifiable, the share of non-disclosing listings should converge rapidly toward zero. The French data reject this prediction in both segments, though with different magnitudes. Figure~\ref{fig:fact2_missing} plots the monthly share of listings without a certified EPC from October 2018 through July 2024. For apartments (panel~(a)), the share averages $35.9\%$ before the reform and $33.8\%$ afterward; over the final twelve months of the sample, it stabilizes at $28.9\%$. For houses (panel~(b)), the corresponding numbers are $24.6\%$ pre-reform, $14.4\%$ post-reform, and a final-twelve-month mean of $9.7\%$. Houses thus display better baseline audit coverage and faster post-reform convergence, but in both segments the share remains structurally bounded away from zero more than three years after disclosure became legally required. Almost three-tenths of apartment listings and one-tenth of house listings still come to market without a certified audit, indicating that a substantial fraction of sellers cannot disclose, because they do not possess a certified audit at the time of listing.

\begin{figure}[htbp]
    \centering
    \caption{Persistence of non-disclosure after the DPE reform}
    \label{fig:fact2_missing}
    \begin{subfigure}{0.49\textwidth}
        \centering
        \caption{Apartments}
        \includegraphics[width=\textwidth]{fig/stylized_facts/apartment/fig_fact2_missing_share_timeseries.pdf}
    \end{subfigure}
    \begin{subfigure}{0.49\textwidth}
        \centering
        \caption{Houses}
        \includegraphics[width=\textwidth]{fig/stylized_facts/house/fig_fact2_missing_share_timeseries.pdf}
    \end{subfigure}
    \begin{minipage}{0.95\textwidth}
        \footnotesize\textit{Notes.} Monthly share of listings without a certified EPC, October 2018 to July 2024. The horizontal dashed line marks each segment's final-twelve-month mean ($0.29$ apartments, $0.10$ houses). The vertical dashed line marks the July 2021 reform date. Volumes are summed across markets within each segment. Source: SeLoger residential listings matched against the ADEME registry of certified energy audits; authors' calculations.
    \end{minipage}
\end{figure}

\subsubsection{Direct Decomposition of the Listing Pool by Information Status}
\label{ss:fact3}

A central methodological contribution of this paper is that, by matching every listing against the ADEME registry, we can directly observe whether a seller is informed ($K=1$, audit exists) or uninformed ($K=0$, no audit on record). Crossing this latent information status with the seller's disclosure choice $d \in \{0,1\}$ yields a four-cell classification of which only three cells are economically distinct: truthful disclosers ($K=1, d=1$), strategically hiding sellers ($K=1, d=0$), and genuinely uninformed sellers ($K=0, d=0$).\footnote{The fourth cell ($K=0$, $d=1$) comprises sellers who display an EPC score for which no matching record appears in the public ADEME registry. These listings are treated as truthful disclosers, consistent with the no-falsification assumption: a score published by the owner is taken at face value, and only omission is treated as a strategic act. Two institutional features explain the size of this group. First, the public registry contains only certificates issued from July 2021 onward, whereas a certificate issued earlier could still be legally valid at the time of listing. Under \emph{d\'ecret} n\textsuperscript{o}~2020-1610 of 17 December 2020, certificates issued between 1 January 2013 and 31 December 2017 remained valid until 31 December 2022, and those issued between 1 January 2018 and 30 June 2021 until 31 December 2024, while certificates issued from 1 July 2021 are valid for ten years. A seller listing during our sample window could therefore display a still-valid certificate that is absent from the registry by construction. Second, the model assumes no falsification. Consistent with the first explanation, the share of listings in this cell falls monotonically as the registry accumulates post-reform certificates: from $62.5\%$ of apartment listings in 2021Q3 to $18.4\%$ in 2024Q4, and from $73.1\%$ to $43.7\%$ for houses.} Figure~\ref{fig:fact3_decomposition} reports the resulting decomposition for both segments. In apartments (panel~(a)), $79.8\%$ of listings are truthful disclosers, $16.7\%$ are genuinely uninformed, and $3.5\%$ are strategically hiding. In houses (panel~(b)), the corresponding shares are $90.4\%$, $8.5\%$, and $1.1\%$. Houses systematically display higher baseline disclosure and a smaller strategic-hiding share, but the qualitative structure is identical: a non-trivial mass of strategic hiders coexists with a much larger mass of genuinely uninformed sellers in both segments.

\begin{figure}[htbp]
    \centering
    \caption{Three-group decomposition of the listing pool}
    \label{fig:fact3_decomposition}
    \begin{subfigure}{0.49\textwidth}
        \centering
        \caption{Apartments}
        \includegraphics[width=\textwidth]{fig/stylized_facts/apartment/fig_fact3_three_group_decomposition.pdf}
    \end{subfigure}
    \begin{subfigure}{0.49\textwidth}
        \centering
        \caption{Houses}
        \includegraphics[width=\textwidth]{fig/stylized_facts/house/fig_fact3_three_group_decomposition.pdf}
    \end{subfigure}
    \begin{minipage}{0.95\textwidth}
        \footnotesize\textit{Notes.} Shares from the estimation sample of each property type: 834{,}816 apartment listings and 677{,}011 house listings, after registry-match cleaning and market-size filtering. Source: SeLoger residential listings matched against the ADEME registry of certified energy audits; authors' calculations.
    \end{minipage}
\end{figure}

\noindent Two features of this decomposition are central to the analysis. First, the non-disclosing pool ($d=0$) is mechanically composite in both segments: roughly five-sixths of apartment non-disclosers and seven-eighths of house non-disclosers are genuinely uninformed, while the remainder are strategically concealing. The vast majority of silent listings are silent by constraint, not by choice. Second, the relative weight of these two subgroups varies substantially across property types and, as we show below, across local markets, with direct implications for the equilibrium informativeness of non-disclosure.

\subsubsection{Substantial Cross-Market Heterogeneity in Non-Disclosure}
\label{ss:fact4}

Figure~\ref{fig:fact4_map} maps the apartment (panel~(a)) and house (panel~(b)) non-disclosure rates at the d\'epartement level. For apartments, the rate ranges from $6.5\%$ in Aube to $40.9\%$ in Corse-du-Sud, with a cross-d\'epartement median of $16.5\%$; high values cluster in southern France, Corsica, and parts of the Parisian region (Seine-Saint-Denis at $32.1\%$). For houses, the range is $5.4\%$ in Essonne to $37.4\%$ in Paris, with a median of $10.0\%$; the Paris d\'epartement is an outlier on the high end, while southern and rural d\'epartements lie closer to the metropolitan baseline. The geographic patterns differ across the two segments, reflecting the different local composition of each property type's listing pool, but the underlying message is identical: within-segment cross-d\'epartement dispersion is large in both cases. This heterogeneity rules out a national-level treatment of the information environment in either segment: any credible structural exercise must allow the depth of audit coverage and the salience of EPC information to vary by local market.

\begin{figure}[htbp]
    \centering
    \caption{Non-disclosure rate by d\'epartement, July 2021 to December 2024}
    \label{fig:fact4_map}
    \begin{subfigure}{0.49\textwidth}
        \centering
        \caption{Apartments}
        \includegraphics[width=\textwidth]{fig/stylized_facts/apartment/fig_fact5_nondiscl_rate_map.pdf}
    \end{subfigure}
    \begin{subfigure}{0.49\textwidth}
        \centering
        \caption{Houses}
        \includegraphics[width=\textwidth]{fig/stylized_facts/house/fig_fact5_nondiscl_rate_map.pdf}
    \end{subfigure}
    \begin{minipage}{0.95\textwidth}
        \footnotesize\textit{Notes.} Shares are computed on listings with a valid d\'epartement code; d\'epartements with fewer than 200 listings within each segment are shown in light grey, and overseas d\'epartements are omitted. Classification: quantiles ($k=5$) on metropolitan France, computed separately within each panel so the breakpoints reflect each segment's distribution. Projection: Lambert-93 (EPSG:2154). Source: SeLoger residential listings matched against the ADEME registry of certified energy audits; d\'epartement boundaries from \texttt{france-geojson} (G.~David, 2024).
    \end{minipage}
\end{figure}

\subsection{What the Facts Demand of the Theory}
\label{ss:facts_demand}

Taken together, the four stylized facts impose tight discipline on any candidate theory of EPC disclosure. The documented non-disclosure premium rules out classical full-unraveling models with skeptical buyers, in which silence must trade at or below the worst certified score. The persistence of non-disclosure after the reform rules out frictionless implementations in which enforcement is complete and disclosure is universal at any positive cost of concealment. The composite nature of the non-disclosing pool requires the theory to recognize that silence aggregates two economically distinct subpopulations: informed sellers who choose to conceal unfavorable information, and uninformed sellers who cannot disclose because no certified audit exists. Their joint mass determines the informational content of non-disclosure and cannot be collapsed into a single type. Finally, the substantial cross-market heterogeneity in non-disclosure rates and price gaps rules out a uniform national treatment of these structural parameters: estimation must be conducted market by market. The model developed in the next section is designed to satisfy these four requirements simultaneously.


% =====================================================================
% §4 A STRUCTURAL INTERPRETATION — consolidation of the former
% Model / Equilibrium Analysis / Identification / Results sections.
% 4.1 is new prose; 4.2-4.5 are relocated verbatim, demoted one level.
% Labels sec:model, sec:equilibrium, sec:identification, sec:results are
% carried onto the corresponding subsections (they have live references).
% =====================================================================

\section{A Structural Interpretation}
\label{sec:model}

\subsection{Motivation}
\label{ss:model_motivation}

To recover the forces behind the patterns documented in
Section~\ref{sec:setting}, we specify and estimate a model of the seller's disclosure
decision. The model is a housing-market application of \citet{dye_disclosure_1985}, as
equilibrium-characterised by \citet{jung_kwon_1988}, with one structural addition:
heterogeneity in buyer EPC sensitivity ($\rho$). The binary quality space
($\theta\in\{0,1\}$) and the seller-information heterogeneity ($\pi$) are adapted directly
from those antecedents. What goes beyond them is the closed-form expression for the
non-disclosure premium and its invariance to market composition in the partial-pooling
regime (Proposition~\ref{prop:equilibrium}), the strict cross-regime ordering of the premium
(Eq.~\ref{eq:premium}), and an identification strategy that uses the directly observed audit
status $K$ to recover $c_0$ and $\rho$ market by market.

We allow for, but do not impose, strategic concealment. Nothing in the specification requires
informed sellers to withhold: the hiding rate is an equilibrium object, the model admits full
unravelling as a special case, and whether concealment occurs --- and how much --- is what the
estimation determines.

Two considerations motivate this structure. The first is resemblance. The facts in
Section~\ref{sec:facts} line up with the Dye--Jung-Kwon configuration: a silent pool that is
overwhelmingly silent by constraint rather than by choice, a small but non-trivial mass of
informed sellers who withhold, and a premium attached to silence rather than a penalty. The
second is institutional. The binary quality space has an empirical anchor in the rental-ban
threshold introduced by \emph{loi Climat et R\'esilience}: dwellings rated E, F, or G face a
progressive prohibition on new lettings, which draws a bright line between the bottom of the
grade distribution and the rest and maps onto $\theta\in\{0,1\}$ without an arbitrary cutoff.

A competing explanation must also clear a demanding timing test. Any account that attributes
the post-reform widening to a supply-side shift has to generate an abrupt change at
July~2021 in \emph{both} the quality premium and the non-disclosure premium, in the same
markets and on the same calendar date, while leaving the composition of the listing pool
undisturbed. The share of listings without a disclosed score is, however, essentially flat
across the reform in the apartment segment (Section~\ref{ss:fact2}), which is the segment
where identification is sharpest. A shift in the mix of sellers would generically move the
two gaps by different amounts, and would show up in that share. Neither is observed.

The mechanism itself is the one set out in Section~\ref{sec:unraveling}: sellers who cannot
disclose dilute the adverse inference buyers draw from silence, which makes withholding
attractive to informed sellers of poor quality, whose entry the buyer cannot separate from
the involuntarily silent. France sits in the bottom-right cell of
Figure~\ref{fig:equilibrium_outcomes} --- incomplete registry coverage combined with positive
buyer EPC-sensitivity. Even where the EPC-sensitive buyer share approaches one, as it does in
most apartment markets after the reform, incomplete coverage alone sustains the premium: so
long as a non-trivial share of non-disclosers hold no certificate, buyers cannot fully
penalise silence. Full unravelling requires both conditions; France satisfies only one.

Finally, we are explicit about what the model does not settle. It takes no stance on
\emph{why} an informed seller withholds --- deliberate concealment and an unrecorded but
still-valid certificate are observationally similar in the listing data, and we do not
attempt to separate them. It assumes no falsification, so a displayed score is taken at face
value. And it maintains Corollary~\ref{cor:hq_disclose}, $\sigma_1^*=0$, rather than testing
it: informed non-disclosers are assigned low quality by revealed preference, which is an
assumption and not a result. Appendix~\ref{app:robustness} reports estimates that instead read
the withheld quality from the registry, so the restriction can be assessed. Nor does the
design pin down $\rho_m$ from the interior in most markets: the data are consistent with full
buyer sensitivity but cannot establish how close to one it lies, and we maintain
$\hat\rho_m = 1$ for those markets and treat $\hat c_{0,m}$ as an upper bound accordingly. We
do not purport to estimate the one true model of housing disclosure; we propose one way of
reading these data through a workhorse framework.

\subsection{Setup}
\label{ss:model_setup}

\subsubsection{Environment and Timing}

A seller owns a property with observable attributes $X$ and an unobservable energy performance level $\theta \in \{0,1\}$, where $\theta=1$ denotes high efficiency. Whether the seller knows $\theta$ depends on whether she has obtained a certified Diagnostic de Performance \'{E}nerg\'{e}tique (DPE) audit prior to listing. With probability $\pi\in(0,1)$, the seller is \emph{informed} ($K=1$) and knows $\theta$; with probability $1-\pi$, she is \emph{uninformed} ($K=0$) and does not. The parameter $\pi$ is not a choice variable: it reflects structural features of local audit markets, including the cost and availability of accredited assessors, the time pressure sellers face (relocation, financial distress, inheritance), and the regulatory requirements in force. It is measured directly from the ADEME registry match, as described in Section~\ref{sec:data}.

The seller simultaneously chooses a disclosure decision $d\in\{0,1\}$ and a listing price $p^{\text{list}}$. French property law requires that all mandatory diagnostic documents, including the DPE, be appended to the sales agreement at signing (loi ALUR, Article L271-4 CCH), so the disclosure decision and the listing price are set jointly at the time the property is placed on the market. This institutional timing motivates the simultaneous-choice formulation: the seller does not first decide whether to disclose and then price, but sets a complete strategy that internalizes how the disclosure decision shifts buyer beliefs and thereby affects both the probability of sale and the expected transaction price.

\noindent\textbf{Assumption (No Falsification).} \textit{Disclosed EPC scores are verifiable against the ADEME registry. An informed seller with $K=1$ may choose $d\in\{0,1\}$, while an uninformed seller with $K=0$ is constrained to $d=0$: she cannot disclose what she does not possess. This rules out costless misreporting.}

\medskip
The key friction is that buyers observe $(d, p^{\text{list}}, X)$ but not the seller's information status $K$: they cannot distinguish a seller who chose not to disclose from one who could not. A buyer then arrives and decides whether to purchase at $p^{\text{list}}$. With probability $\rho\in(0,1)$, the buyer is \emph{EPC-sensitive} and incorporates energy performance into her valuation; with probability $1-\rho$, the buyer is \emph{EPC-insensitive} and conditions her bid solely on observable attributes. Sellers cannot observe buyer type ex ante and optimize against expected demand. The parameter $\rho$ is market-specific, reflecting the local composition of the buyer pool, and is identified from the differential price response to the July 2021 reform across markets, as detailed in Section~\ref{sec:identification}.

We provide a visual representation of the model dynamics and timing in Figure~\ref{fig:timing}; a fuller schematic overview appears later in Figure~\ref{fig:vis_model}.

\begin{figure}[H]
    \centering
    \includegraphics[width=\textwidth]{fig/visuals/model_dynamics.png}
    \caption{Model timing}
\label{fig:timing}
\end{figure}

\subsubsection{Payoffs and Hiding Costs}

In French residential markets, listing prices serve as anchors for negotiation rather than posted take-it-or-leave-it offers: the transaction price conditional on a sale typically falls below the listing price by a margin that depends on local market conditions and property characteristics.\footnote{The negotiation margin in France is empirically small: the majority of properties transact at or very near list price, and the average discount rarely exceeds three to five percent. \citet{merlo_bargaining_2004} document similar patterns for England, showing that most residential transactions involve little or no negotiation once a listing price is set. \citet{genesove_loss_2001} provide evidence that list prices act as anchors that constrain final transaction prices. For France specifically, \citet{friggit_marche_2013} documents that the average negotiation margin has remained narrow and declining, with a large share of properties selling without any price reduction.} Buyer arrival and match quality are stochastic, so the seller faces uncertainty over both whether a sale occurs and, conditional on a sale, at what price. The seller's problem therefore separates into a transaction probability $\Gamma$ and an expected transaction price $\phi$, each depending on the listing price and disclosure decision:
\begin{equation}
    \max_{\{d,\,p^{\text{list}}\}} \;\Pi = \underbrace{\mathbb{E}\!\left[p^{\text{trans}}\mid d,p^{\text{list}},X\right]}_{\phi(p^{\text{list}})} \cdot \underbrace{\Pr(T{=}1\mid d,p^{\text{list}},X)}_{\Gamma(p^{\text{list}})} - C(\theta,d)
    \label{eq:payoff}
\end{equation}
where $\phi(p^{\text{list}}) = \gamma_0+\gamma_1 \cdot p^{\text{list}}+\sum_n\gamma_n X_n$ is the expected transaction price conditional on a sale occurring, and $C(\theta,d)=\mathbf{1}[d=0]\cdot(c_0+\theta c_1)$ is the hiding cost, paid only when the seller withholds information.

The hiding cost has two components. The term $c_0\geq 0$ applies to all non-disclosers, regardless of quality. In the French institutional context, $c_0$ captures the regulatory burden imposed by loi ALUR (legal exposure if a buyer subsequently discovers a concealed audit via the public ADEME registry, notaire liability checks, and the informational ambiguity that prolongs negotiation and increases sale failure risk). The term $c_1\geq 0$ is an additional opportunity cost borne exclusively by high-quality sellers ($\theta=1$) who choose to hide: by not disclosing, they forgo the revenue premium that verified high quality would command from EPC-sensitive buyers. This asymmetry, $c_1>0$, is economically decisive: it means that hiding is strictly cheaper for low-quality sellers than for high-quality ones, so the net gain from concealment is larger for $\theta=0$ sellers at any given pool belief. This differential is the sorting pressure that drives full separation of high-quality sellers in equilibrium (Corollary~\ref{cor:hq_disclose}). Hiding costs are paid unconditionally on sale: they represent compliance burdens and reputational exposure that arise at listing time, not at closing. A note on the dual role of $c_0$ is in order. As a structural primitive, $c_0$ is an institutional cost of concealment that we treat as fixed within the prevailing enforcement regime and recover from the equilibrium price gap (Section~\ref{sec:identification}); as a policy object, it is the margin a regulator shifts through stricter enforcement, and we study counterfactual variations in $c_0$ as such in Section~\ref{sec:policy}. The two uses are consistent: estimation reads $c_0$ off the current regime, while the policy exercise perturbs it around that baseline.

\subsubsection{Buyer Valuations and Demand}

The partition of buyers into EPC-sensitive and EPC-insensitive types reflects a well-documented empirical regularity: willingness to pay for energy efficiency varies sharply with how much the buyer will bear the ongoing energy costs. Primary-residence buyers in northern France with long heating seasons, and investor-buyers who face the impending F/G rental ban under the \textit{loi Climat et R\'{e}silience}, have strong incentives to incorporate EPC scores into their bids. Vacation-home buyers who occupy the property a few weeks per year, and buyers in mild climates, may rationally discount energy performance. This heterogeneity determines $\rho$, the share of the local buyer pool that is EPC-sensitive.

EPC-sensitive buyers ($B=1$) condition their valuation on energy performance:
\begin{equation}
    v_1(\theta,d) = \begin{cases}
        \beta_0 + \beta_1\cdot\theta + \textstyle\sum_n\beta_n X_n & \text{if } d=1 \;\text{(verified quality)},\\
        \beta_0 + \beta_1\cdot\mathbb{E}[\theta\mid d{=}0] + \textstyle\sum_n\beta_n \cdot X_n & \text{if } d=0 \;\text{(inferred quality)}.
    \end{cases}
\end{equation}
We adopt the benchmark $\alpha_k=\beta_k$: buyers value quality equally whether it is certified or inferred. This is a conservative assumption in a precise sense: under the alternative $\alpha_k<\beta_k$ (a distrust discount on unverified quality), the non-disclosure premium would be smaller for any given pool composition, so our identified $\rho$ is a lower bound on the true EPC-sensitive buyer share. The non-disclosure premium we document therefore understates what a fully skeptical buyer would generate. EPC-insensitive buyers ($B=0$) value the property solely on observables: $v_0(X)=\delta_0+\sum_n\delta_n X_n$. The aggregate transaction probability mixes the two buyer types:
\begin{equation}
    \Gamma(p^{\text{list}}) = \rho\cdot G\!\left(v_1(d)-p^{\text{list}}\right) + (1-\rho)\cdot G\!\left(v_0-p^{\text{list}}\right)
    \label{eq:demand}
\end{equation}
where $G(\cdot)$ is the logistic CDF governing idiosyncratic taste shocks. The parameter $\rho$ enters the demand function as a weight on the EPC-sensitive buyer valuation: a higher $\rho$ amplifies the revenue difference between a disclosed high-quality listing and a non-disclosing listing, raising the stakes of the disclosure decision for the seller. It is this amplification channel through which $\rho$ governs disclosure incentives. Empirically, $\rho$ is a market-level structural parameter identified from the cross-market variation in the price response to the July 2021 reform, as developed in Section~\ref{sec:identification}.

\subsubsection{Buyer Beliefs and the Non-Disclosing Pool}

When observing a non-disclosing listing, buyers cannot determine whether the seller is strategically concealing an unfavorable audit or simply lacks a certified EPC. They form a posterior belief about expected quality using Bayes' rule. The structure of this belief depends on the composition of the non-disclosing pool, which is itself an equilibrium object.

\paragraph{Who Hides?}

Each seller is characterized by two attributes: information status $K \in \{0,1\}$ and quality $\theta \in \{0,1\}$. We maintain the following assumption throughout.

\begin{assumption}[Independence]\label{ass:independence}
$K \perp \theta$, with $\Pr(K=1) = \pi$ and $\Pr(\theta=1) = \bar{\theta}$.
\end{assumption}

\noindent Independence is most plausible for transactions driven by exogenous life events (succession following a death, court-ordered liquidation, or divorce settlement), in which legal proceedings impose binding timing constraints that preclude any quality-contingent decision to obtain a certified EPC. In such cases, the seller's information status $K$ is determined by administrative capacity and procedural deadlines rather than the energy performance of the dwelling; a seller under a judicial mandate to liquidate an estate has neither the incentive nor the discretion to condition audit acquisition on $\theta$. In France, forced transactions of this type constitute a substantial share of the resale market, providing a natural mechanism through which $K \perp \theta$ obtains in the data.

Table~\ref{tab:who_hides} summarizes disclosure behavior by seller type. Uninformed sellers have no choice: they cannot disclose information they do not possess (No Falsification Assumption). Informed sellers choose strategically, weighing the revenue gain from disclosure against the option of pooling with the non-disclosing mass. We denote by $\sigma_\theta \equiv \Pr(d=0 \mid K=1, \theta)$ the equilibrium hiding probability for an informed seller with quality $\theta$.

\begin{table}[H]
\centering
\caption{Disclosure Behavior by Seller Type}
\label{tab:who_hides}
\small
\renewcommand{\arraystretch}{1.3}
\setlength{\tabcolsep}{20pt}
\begin{tabular}{lcccl}
\toprule
\toprule
\textbf{Type} & \textbf{Mass} & \textbf{Can Disclose?} & \textbf{Hides?} & \textbf{Hiding Mass} \\
\midrule
Informed, $\theta=1$ & $\pi \bar{\theta}$ & Yes & Prob.\ $\sigma_1$ & $\pi \bar{\theta} \sigma_1$ \\
Informed, $\theta=0$ & $\pi (1-\bar{\theta})$ & Yes & Prob.\ $\sigma_0$ & $\pi (1-\bar{\theta}) \sigma_0$ \\
Uninformed, $\theta=1$ & $(1-\pi) \bar{\theta}$ & No & Always & $(1-\pi) \bar{\theta}$ \\
Uninformed, $\theta=0$ & $(1-\pi) (1-\bar{\theta})$ & No & Always & $(1-\pi) (1-\bar{\theta})$ \\
\bottomrule
\bottomrule
\end{tabular}

\vspace{0.5em}
\raggedright
\footnotesize
\textit{Notes:} $\sigma_\theta \equiv \Pr(d=0 \mid K=1, \theta)$ denotes the hiding probability for an informed seller with quality $\theta$. Uninformed sellers cannot disclose and hide with probability one.
\end{table}

\paragraph{The Belief Function}

Applying Bayes' rule to the pool composition in Table~\ref{tab:who_hides}, the expected quality among non-disclosers is:
\begin{equation}
    e(\sigma_0, \sigma_1) \equiv \Pr(\theta=1 \mid d=0) = \frac{\pi \bar{\theta} \sigma_1 + (1-\pi) \bar{\theta}}{\pi \bar{\theta} \sigma_1 + \pi (1-\bar{\theta}) \sigma_0 + (1-\pi)}
    \label{eq:belief_general}
\end{equation}

The numerator counts the high-quality mass in the pool: informed high-quality sellers who choose to hide ($\pi\bar{\theta}\sigma_1$) together with uninformed sellers who happen to have high quality ($(1-\pi)\bar{\theta}$). The denominator is the total non-disclosing mass. The belief $e$ is therefore a weighted average across two economically distinct populations. When uninformed sellers dominate the pool ($\pi$ small or $\sigma_0$ small), non-disclosure looks like a constraint rather than a choice and beliefs remain favorable to the seller. When strategic low-quality hiders dominate ($\pi$ large and $\sigma_0$ large), the pool is contaminated and beliefs deteriorate. This is why unraveling is incomplete: non-disclosure is inherently ambiguous when $\pi < 1$, whereas in classical disclosure theory all sellers are informed and silence is unambiguous bad news.

\begin{lemma}[Belief Monotonicity]\label{lem:belief_monotone}
$e(\sigma_0,\sigma_1)$ is strictly decreasing in $\sigma_0$ and strictly increasing in $\sigma_1$.
\end{lemma}
\begin{proof}[Proof sketch]
Increasing $\sigma_0$ raises the denominator (more low-quality hiders enter the pool) while leaving the numerator unchanged, so $e$ falls. Increasing $\sigma_1$ raises the numerator and the denominator by the same term $\pi\bar{\theta}$, and since $e < 1$ the numerator effect dominates, so $e$ rises. The strict monotonicity in $\sigma_0$ is the self-correcting force that ensures equilibrium uniqueness in Section~\ref{sec:equilibrium}: as more informed low-quality sellers choose to hide, buyer beliefs worsen, attenuating the benefit of hiding and restoring incentive compatibility.
\end{proof}

\subsection{Equilibrium}
\label{sec:equilibrium}

\subsubsection{Disclosure Incentives}

The seller's disclosure decision turns on a single comparison: whether the price attainable by withholding the audit result, $p^*_{d=0}$, priced off the pool belief $e$, exceeds the price attainable by disclosing true quality $\theta$. A low-quality seller ($\theta=0$) gains from non-disclosure by pooling with uninformed sellers and extracting the associated quality premium; a high-quality seller ($\theta=1$) loses from non-disclosure because pooling with low-quality hiders depresses the belief $e$ below 1. This asymmetry drives the central result.

\begin{proposition}[Disclosure Incentives and Seller Sorting]
\label{prop:premium}
Under the benchmark $\alpha_k=\beta_k$, a seller of type $(K,\theta)$ strictly prefers disclosure if and only if $\theta > e(\sigma_0^*)$. In particular:
\begin{enumerate}[label=(\roman*)]
    \item An informed seller with $\theta=1$ strictly prefers disclosure for any $e(\sigma_0^*)<1$.
    \item An informed seller with $\theta=0$ strictly prefers non-disclosure whenever $e(\sigma_0^*)>0$.
\end{enumerate}
\end{proposition}

\noindent\textit{Proof.} See Appendix~\ref{app:proofs}. $\Box$

Proposition~\ref{prop:premium} reflects a standard adverse-selection asymmetry. For a high-quality seller, disclosure reveals quality exactly, whereas non-disclosure pools her with low-quality hiders; the condition $e(\sigma_0^*)<1$ is satisfied for any non-degenerate pool and so is always binding. For a low-quality seller, concealment allows her to price off the pool belief $e(\sigma_0^*)>0$, which is strictly greater than her true quality $\theta=0$; the gain from pooling is therefore strictly positive whenever uninformed sellers are present. The additional hiding cost $c_1>0$ borne by high-quality concealers reinforces but does not drive this separation under the benchmark $\alpha_k=\beta_k$: the pool-inference argument alone is sufficient to induce full separation of high-quality sellers, with $c_1$ becoming binding only when $\alpha_k < \beta_k$ (the distrust discount case).

\begin{corollary}[High-Quality Sellers Always Disclose]
\label{cor:hq_disclose}
In any equilibrium, $\sigma_1^*=0$.
\end{corollary}

\noindent\textit{Proof.} Follows directly from Proposition~\ref{prop:premium}(i). $\Box$

Corollary~\ref{cor:hq_disclose} is a partial unraveling result: the top of the quality distribution fully separates. An informed high-quality seller gains nothing from ambiguity and is strictly harmed by being pooled with low-quality hiders; she therefore always discloses. The non-disclosing pool is consequently free of high-quality informed sellers in any equilibrium, and the strategic question reduces entirely to the behavior of informed low-quality sellers.

Since $\sigma_1^*=0$, the equilibrium belief simplifies. Defining $\sigma_0\equiv\Pr(d=0\mid K=1,\theta=0)$, the non-disclosing pool consists exclusively of uninformed sellers of both quality types and informed low-quality hiders:
\begin{equation}
    e(\sigma_0) = \frac{(1-\pi)\bar{\theta}}{\pi(1-\bar{\theta})\sigma_0 + (1-\pi)}
    \label{eq:belief_simplified}
\end{equation}
The belief $e(\sigma_0)$ is strictly decreasing in $\sigma_0$: each additional low-quality hider enters the denominator with weight $\pi(1-\bar{\theta})$ while leaving the numerator unchanged, diluting pool quality and depressing buyer beliefs.

\subsubsection{Equilibrium Characterization}

Let $\Delta(e)\equiv\rho\cdot\beta_1\cdot e$ denote the revenue gain from non-disclosure for a low-quality seller facing pool belief $e$. Define the net hiding gain:
\begin{equation}
    \Phi(\sigma_0) \;\equiv\; \Delta\!\left(e(\sigma_0)\right) - c_0 \;=\; \frac{\rho\,\beta_1\,(1-\pi)\bar{\theta}}{\pi(1-\bar{\theta})\sigma_0 + (1-\pi)} - c_0.
    \label{eq:Phi}
\end{equation}
The function $\Phi$ is strictly decreasing in $\sigma_0$: as more informed low-quality sellers conceal, pool quality $e$ deteriorates, the gain from hiding falls, and the incentive to conceal dissipates. This monotone feedback delivers a unique equilibrium without reliance on equilibrium refinements. The three equilibrium regimes follow directly from comparing $\Phi$ to zero at the two boundary values: $\Phi(0)=\Delta(\bar{\theta})-c_0$ measures the net gain when the pool is at its cleanest, and $\Phi(1)=\Delta(e(1))-c_0$ when it is maximally contaminated.

\begin{proposition}[Equilibrium Existence and Uniqueness]
\label{prop:equilibrium}
There exists a unique equilibrium, characterized as follows:
\begin{enumerate}[label=(\roman*)]
    \item If $c_0>\Delta(\bar{\theta})$, the unique equilibrium is \emph{Cost-Deterred Hiding}: $\sigma_0^*=0$ and $e^*=\bar{\theta}$.
    \item If $c_0<\Delta(e(1))$, the unique equilibrium is \emph{Complete Adverse Selection}: $\sigma_0^*=1$ and $e^*=(1-\pi)\bar{\theta}/(1-\pi\bar{\theta})$.
    \item If $\Delta(e(1))\leq c_0\leq\Delta(\bar{\theta})$, the unique equilibrium is \emph{Partial Pooling}: $\sigma_0^*\in(0,1)$ is the unique solution to $\Delta(e(\sigma_0))=c_0$.
\end{enumerate}
\end{proposition}

\noindent\textit{Proof.} See Appendix~\ref{app:proofs}. $\Box$

\begin{table}[h]
\centering
\caption{Equilibrium Regimes}
\label{tab:equilibria}
\small\renewcommand{\arraystretch}{1.5}
\begin{tabular}{lcccc}
\toprule\toprule
\textbf{Equilibrium} & \textbf{Condition} & \textbf{Strategy $\sigma_0^*$} & \textbf{Belief $e^*$} & \textbf{Premium $p_0^*-p_1^*$} \\
\midrule
Cost-Deterred Hiding       & $c_0>\Delta(\bar{\theta})$                     & $0$     & $\bar{\theta}$                                          & $\rho\beta_1\bar{\theta}$ \\
Complete Adverse Selection & $c_0<\Delta(e(1))$                             & $1$     & $\dfrac{(1-\pi)\bar{\theta}}{1-\pi\bar{\theta}}$        & $\rho\beta_1 e^*$ \\[6pt]
Partial Pooling            & $\Delta(e(1))\leq c_0\leq\Delta(\bar{\theta})$ & $(0,1)$ & Solves $\Phi(\sigma_0^*)=0$                             & $\rho\beta_1 e(\sigma_0^*)$ \\
\bottomrule\bottomrule
\end{tabular}
\end{table}

Neither corner solution fits the data: universal concealment would require the non-disclosure premium to exceed $c_0$ even against the most contaminated pool, while universal disclosure would require concealment to be unprofitable even against the cleanest. What we observe instead is an interior equilibrium in which some low-quality sellers disclose and the remainder pool with uninformed non-disclosers. Substituting~\eqref{eq:Phi} into $\Phi(\sigma_0^*)=0$ yields the closed-form hiding rate:
\begin{equation}
    \sigma_0^* = \frac{1}{\pi(1-\bar{\theta})}\left[\frac{\rho\,\beta_1\,(1-\pi)\bar{\theta}}{c_0} - (1-\pi)\right].
    \label{eq:pp_sigma}
\end{equation}

We derive the non-disclosure premium as follows:
\begin{equation}
    p_0^*-p_1^* \;=\;
    \begin{cases}
        \rho\beta_1\bar{\theta}
            & \text{Cost-Deterred Hiding }(\sigma_0^*=0), \\[4pt]
        c_0
            & \text{Partial Pooling }(\sigma_0^*\in(0,1)), \\[4pt]
        \rho\beta_1\dfrac{(1-\pi)\bar{\theta}}{1-\pi\bar{\theta}}
            & \text{Complete Adverse Selection }(\sigma_0^*=1),
    \end{cases}
    \quad \text{with}\quad \rho\beta_1\frac{(1-\pi)\bar{\theta}}{1-\pi\bar{\theta}} < c_0 < \rho\beta_1\bar{\theta}.
    \label{eq:premium}
\end{equation}

The three equilibria characterized in Proposition~\ref{prop:equilibrium} arise from a single comparison: the hiding gain $\Delta(e)=\rho\beta_1 e$ evaluated at the two extremes of pool contamination against the structural hiding cost $c_0$. When the gain is dominated by $c_0$ even at the cleanest possible pool ($\sigma_0=0$, $e=\bar{\theta}$), no informed seller hides; when it survives the most contaminated pool ($\sigma_0=1$, $e=(1-\pi)\bar{\theta}/(1-\pi\bar{\theta})$), every informed seller hides; in between, an interior fraction hides and the system self-equilibrates. The non-disclosure premium inherits a strict ordering across regimes,
\begin{equation*}
    \underbrace{\rho\beta_1 \bar{\theta}}_{\text{Cost-Deterred Hiding}}
    \;>\;
    \underbrace{c_0}_{\text{Partial Pooling}}
    \;>\;
    \underbrace{\rho\beta_1\,\dfrac{(1-\pi)\bar{\theta}}{1-\pi\bar{\theta}}}_{\text{Complete Adverse Selection}},
\end{equation*}
with the interior value depending neither on $\pi$ nor on $\rho$.

In the \emph{Cost-Deterred Hiding} regime, the hiding cost is prohibitive even against the cleanest conceivable non-disclosing pool, the one populated entirely by uninformed sellers drawn from the unconditional quality distribution. No informed low-quality seller finds it profitable to suppress her type, so $\sigma_0^*=0$ and the non-disclosing pool is uncontaminated by strategic behavior: $e^*=\bar{\theta}$. The premium $\rho\beta_1\bar{\theta}$ reflects the full unconditional quality belief, scaled by the share of EPC-sensitive buyers and their marginal valuation. Relative to \citet{milgrom_good_1981} and \citet{grossman_informational_1981}, this regime delivers a form of deterrence (informed bad types refrain from concealment) but through a fundamentally different mechanism: unraveling in those models is sustained by a belief-penalty channel, in which silence triggers the worst possible inference, whereas here deterrence runs through the cost channel $c_0$ because uninformed sellers prevent silence from carrying the worst-type inference. This is the only regime in which the non-disclosing pool is informationally clean.

In the \emph{Complete Adverse Selection} regime, the hiding gain survives the most contaminated pool the model can sustain. Every informed low-quality seller hides, $\sigma_0^*=1$, and buyer beliefs about the non-disclosing pool collapse to the floor $e^*=(1-\pi)\bar{\theta}/(1-\pi\bar{\theta})$, strictly below $\bar{\theta}$. The non-disclosure premium accordingly hits its lower bound. This lower bound is nonetheless strictly positive: the uninformed seller mass dilutes the adverse inference and prevents the complete belief collapse that Milgrom--Grossman unraveling would imply. The mechanism is exactly the one identified by \citet{dye_disclosure_1985}, in which a residual silent population that is not strategically motivated keeps silence partially informative, and it operates here even when strategic hiders saturate.

In the \emph{Partial Pooling} regime, $\Phi(\sigma_0)$ is positive at $\sigma_0=0$ and negative at $\sigma_0=1$. Strict monotonicity yields a unique interior crossing $\sigma_0^*\in(0,1)$ at which the marginal hider is exactly indifferent: $\Delta(e(\sigma_0^*))=c_0$. The underlying mechanism is a feedback loop. As more informed low-quality sellers join the non-disclosing pool, $e(\sigma_0)$ falls, perceived pool quality deteriorates, and the gain $\Delta(e)$ from hiding shrinks. Entry into the pool stops precisely when this gain has been driven down to the hiding cost. The fraction $\sigma_0^*$ of low-quality types hides and pools with the uninformed; the remaining $1-\sigma_0^*$ disclose, accepting the low-quality price rather than continuing to contaminate the pool.

The structural payoff of this regime is that the equilibrium premium equals exactly $c_0$, invariant to $\pi$ and $\rho$ throughout the interior. The reason is the indifference condition itself: in equilibrium, the price difference between the non-disclosing pool and the disclosed low-quality outcome must equal the hiding cost, otherwise the marginal seller would strictly prefer one action over the other. To see the invariance, consider a rise in $\pi$. A larger informed share would, holding $\sigma_0^*$ fixed, mechanically contaminate the pool and depress $e$. Equilibrium restores $\Phi=0$ by reducing $\sigma_0^*$: fewer hiders produce a cleaner pool, $e$ rebounds, and the premium returns to $c_0$. An analogous adjustment absorbs changes in $\rho$. The non-disclosure premium therefore functions as a revealed hiding cost rather than a free function of market composition, and the observable price gap between non-disclosing and disclosing low-quality sellers directly identifies $c_0$. This is the only regime in which $\pi$, $\rho$, and $c_0$ are jointly identified.

\subsubsection{Comparative Statics}

\begin{proposition}[Comparative Statics]
\label{prop:comparative_statics}
The following comparative statics hold:
\begin{enumerate}[label=(\roman*)]
    \item In the \emph{Complete Adverse Selection} equilibrium: $\partial(p_0^*-p_1^*)/\partial\pi < 0$ and $\partial(p_0^*-p_1^*)/\partial\rho > 0$.
    \item In the \emph{Partial Pooling} equilibrium: $\partial\sigma_0^*/\partial c_0 < 0$, $\partial\sigma_0^*/\partial\rho > 0$, and $\partial\sigma_0^*/\partial\pi < 0$.
\end{enumerate}
\end{proposition}

\noindent\textit{Proof.} See Appendix~\ref{app:proofs}. $\Box$

\begin{itemize}
    \item Part~(i) reflects two distinct transmission channels. The effect of $\pi$ operates through the \textit{belief channel}: as the share of informed sellers rises, non-disclosure becomes more informative of low quality, $e^*$ falls, and the gain from concealment dissipates. The effect of $\rho$ operates through the \textit{revenue channel}: a larger share of EPC-sensitive buyers widens the price spread between quality levels, raising the value of concealment.
    \item Part~(ii) follows from the implicit function theorem applied to $\Phi(\sigma_0^*)=0$. Since $\Phi'(\sigma_0^*)<0$, one obtains $\partial\sigma_0^*/\partial c_0 < 0$ directly. The effect $\partial\sigma_0^*/\partial\rho > 0$ follows because a higher $\rho$ increases $\Delta$ at any fixed $e$, shifting $\Phi$ upward and raising the interior equilibrium. The effect of $\pi$ operates through the belief function: a higher $\pi$ increases the contamination weight $\pi(1-\bar{\theta})$ in the denominator of $e(\sigma_0)$, reducing $\Delta(e(\sigma_0))$ and thereby shifting $\Phi$ downward, so $\partial\sigma_0^*/\partial\pi < 0$. The parameters $\pi$ and $c_0$ are therefore \textit{policy substitutes}: each reduces strategic hiding through a distinct channel ($\pi$ through beliefs, $c_0$ through costs), and their marginal effects on $\sigma_0^*$ are both negative and quantifiable from the estimated model.
\end{itemize}

\begin{corollary}[Effect of $\rho$ on the non-disclosure premium]
\label{cor:rho}
\begin{enumerate}[label=(\roman*)]
    \item The non-disclosure premium $\Delta = p_0^* - p_1^*$ is weakly increasing in $\rho$: strictly so in the Cost-Deterred Hiding and Complete Adverse Selection regimes ($\partial\Delta/\partial\rho > 0$), and zero in the Partial Pooling regime ($\partial\Delta/\partial\rho = 0$, with the premium pinned at $c_0$).
    \item The parameter interval sustaining the Partial Pooling regime, $[\rho\,\beta_1\,e(1),\,\rho\,\beta_1\,\bar{\theta}]$, is weakly expanding in $\rho$: its width $\rho\,\beta_1[\bar{\theta}-e(1)]$ is strictly increasing in $\rho$ whenever $\pi>0$.
\end{enumerate}
\end{corollary}

\noindent\textit{Proof.} See Appendix~\ref{app:proofs}. $\Box$

\noindent Economically, a larger share of buyers who price the certificate raises what silence
is worth wherever the hiding rate is at a corner, but inside the partial-pooling regime it does
not: there the premium is pinned to the concealment cost, and an increase in $\rho$ is absorbed
by sellers hiding more rather than by the price gap widening. Part~(ii) adds that higher $\rho$
widens the parameter window over which partial pooling survives. Empirically this predicts that
across partial-pooling markets the non-disclosure premium should be uncorrelated with
$\hat\rho_m$ while the hiding rate rises with it, and that markets with higher $\hat\rho_m$
should be less likely to appear in a corner regime --- a pattern we can check directly against
the regime classification reported in Section~\ref{sec:results}.


\begin{proposition}[Equilibrium effect of $\pi$]
\label{prop:pi}
As $\pi$ increases, the lower boundary $\Delta(e(1))$ of the Partial Pooling interval falls ($\partial e(1)/\partial\pi < 0$), so the condition for Complete Adverse Selection ($c_0 < \Delta(e(1))$) becomes harder to satisfy and the equilibrium moves toward the Cost-Deterred Hiding regime. For any fixed hiding cost $c_0$, a sufficiently large increase in $\pi$ drives the equilibrium out of the Partial Pooling region toward full unraveling.
\end{proposition}

\noindent\textit{Proof.} See Appendix~\ref{app:proofs}. $\Box$

\noindent In words: better registry coverage makes silence more informative, because a larger
fraction of non-disclosers are hiding by choice rather than by constraint. Buyers therefore
discount silence more heavily, concealment becomes less attractive, and the market moves toward
disclosure. The testable content is a cross-market comparative static: markets with higher
$\hat\pi_m$ should display lower hiding rates $\hat\sigma_{0,m}$ and should be
disproportionately represented among the markets classified as fully unravelling in
Section~\ref{sec:results}.


\begin{proposition}[Equilibrium effect of $c_0$]
\label{prop:c0}
An increase in $c_0$ shifts $\Phi(\sigma_0) = \Delta(e(\sigma_0)) - c_0$ uniformly downward. The upper boundary condition for Partial Pooling ($c_0 \leq \Delta(\bar{\theta})$) becomes harder to satisfy and the condition for Cost-Deterred Hiding ($c_0 > \Delta(\bar{\theta})$) becomes easier, so higher hiding costs push the equilibrium toward the Cost-Deterred Hiding regime. In the interior of the Partial Pooling regime, $\partial\sigma_0^*/\partial c_0 < 0$: higher costs reduce the equilibrium hiding rate.
\end{proposition}

\noindent\textit{Proof.} See Appendix~\ref{app:proofs}. $\Box$

The premium in~\eqref{eq:premium} generates a transfer from EPC-sensitive buyers to the non-disclosing pool, but it does not accrue symmetrically within that pool. The informed low-quality hider receives the premium but nets zero: the indifference condition $\Delta(e(\sigma_0^*))=c_0$ ensures that the gross revenue gain from concealment is exactly consumed by the hiding cost. The uninformed seller, who bears no hiding cost, is the residual claimant and pockets the full premium $c_0$ as a pure rent. The welfare cost of the equilibrium therefore has two distinct components. First, an allocative distortion: EPC-sensitive buyers pay a positive premium for assets whose true energy performance falls below the pooled belief $e(\sigma_0^*)$, a misallocation that would be eliminated if the concealed information were transmitted. Second, a deadweight loss: the hiding cost $c_0$ is incurred by informed sellers at listing time without any counterpart in buyer valuations. This is a social loss only to the extent that $c_0$ represents real resources destroyed (legal fees, prolonged negotiation, administrative burden); any portion that is a transfer, such as the expected value of a fine accruing to the state, is a redistribution rather than a loss. We do not identify this split, so the deadweight-loss component is an upper bound. A formal efficiency characterization of these two components and the optimal subsidy that addresses them is developed in Section~\ref{sec:policy}.

\subsection{Identification}
\label{sec:identification}

The two structural objects to identify are $\rho_m$, the market-specific share of EPC-sensitive buyers, and $c_{0,m}$, the seller hiding cost. The structural equation $\Delta_m = \rho_m \cdot \beta_{1,m} \cdot e^*_m$ relates these to observables: $\Delta_m$ and $\beta_{1,m}$ are equilibrium price gaps, and $e^*_m$ is directly observed from the ADEME registry. This might appear to make $\rho_m$ point-identified from a single cross-section via $\rho_m = \Delta_m / (\beta_{1,m} \cdot e^*_m)$. The problem is that cross-sectional price gaps reflect \emph{level} differences across markets, not the demand-induced component we need. Any market characteristic that simultaneously shapes the composition of the undisclosed pool and is correlated with buyer EPC-sensitivity would bias a cross-sectional estimate of $\rho_m$. Purging this confound requires variation that isolates the \emph{change} in price gaps attributable to buyers becoming EPC-sensitive, not their levels.

The July 2021 DPE reform provides this variation, and the two-step approach uses it as follows. In the first step, a difference-in-differences estimator compares price gaps in each market before and after the reform, stripping out market-specific level differences and isolating the reform-induced shift in buyer demand locally. This yields $\hat\rho_m = \hat\Delta_m / (\hat\beta_{1,m} \cdot \hat{e}^*_m)$, where the hat quantities are DiD estimates (post minus pre period means). In the second step, the hiding cost $c_{0,m}$ is recovered from the equilibrium hiding rate $\hat\sigma_{0,m}$ via a closed-form inversion of the seller indifference condition (given the equilibrium characterized in Section~\ref{sec:model}). No additional external variation is required: $\hat\sigma_{0,m}$ is directly observed as the fraction of informed low-quality sellers who withhold their EPC, and the inversion is exact. Together, three observables $(\hat\beta_{1,m},\,\hat\Delta_m,\,\hat\sigma_{0,m})$ and one registry-based input $(\hat{e}^*_m)$ map to two structural parameters $(\hat\rho_m,\,\hat{c}_{0,m})$, market by market.

\subsubsection{The July 2021 Reform as a Natural Experiment}

The change in the legal status of the DPE described in Section~\ref{ss:institutions} is the source of identification: it created a precise, legally documented change in the value of a disclosed EPC score to buyers, with a known implementation date of July 1, 2021. Crucially, the reform operated on the \emph{demand} side --- the weight a disclosed rating carries in buyers' valuations and legal recourse --- and not on the seller's cost of withholding. The obligation to provide a DPE at sale, and the compliance and liability exposure attached to coming to market without the required diagnostic, were already in force under loi ALUR (Article L271-4 CCH) well before 2021. These are precisely the components of the concealment cost $c_{0,m}$. What July 2021 changed was not whether a seller must disclose, nor the cost she bears for withholding, but how much a \emph{disclosed} rating now matters to the buyer. The reform therefore shifts $\rho_m$ while leaving $c_{0,m}$, governed by the pre-existing disclosure mandate, unchanged across the reform date. This is what licenses reading the DiD change in the premium as the demand response that identifies $\rho_m$, holding the hiding cost fixed. Whether the supply side of the market responded concurrently is a maintained assumption we cannot directly test with the data at hand; we return to it below.\footnote{An alternative strategy based on a Bartik instrument (interacting wholesale gas price variation with the pre-reform share of low-efficiency properties per market) fails at the first stage ($F$-statistic: 1.0 to 1.6). The French \emph{bouclier tarifaire} (energy price shield, October 2021 to 2023) decoupled household energy costs from the TTF gas price spike driven by the Russian invasion of Ukraine, removing the demand-side variation the instrument was designed to exploit.}

The identifying assumption is the standard parallel-trends condition: in the absence of the reform, the cross-group price gaps $\beta_{1,m,t}$ (energy-efficient vs.\ energy-intensive disclosed properties) and $\Delta_{m,t}$ (undisclosed vs.\ energy-intensive) would have continued on their pre-reform trajectory. The DiD does not require that buyers were entirely inattentive to EPC scores before July 2021; it requires only that their relative pricing of disclosure groups was not already trending. Any post-reform shift in the gaps, relative to that pre-reform baseline, is then attributable to the change in buyer incentives induced by legal opposability.

The pre-trend evidence supports this assumption directly, and more strongly than a single test would. Both $\beta_{1,m,t}$ and $\Delta_{m,t}$ are statistically flat across all markets in the 33 months before the reform (Figure~\ref{fig:ATT}). These two gaps measure distinct economic objects: $\beta_{1,m,t}$ captures buyers' willingness to pay for verified high quality, while $\Delta_{m,t}$ reflects the informational wedge created by strategic non-disclosure. They respond to different primitives of the model. A supply-side confound (a shift in $\pi_m$, $\bar\theta_m$, or seller composition at the reform date) would generically move these two gaps by different amounts and in potentially different directions, since the mechanisms through which supply composition enters each gap differ. The fact that both gaps are jointly flat before July 2021 and both shift upward simultaneously after the reform, at the same markets and on the same calendar date, is difficult to reconcile with a supply-side story. It is exactly what the demand-side interpretation predicts: legal opposability increases the return to EPC quality for sensitive buyers, raising both the price of verified high quality ($\beta_{1,m,t}$) and the information premium extracted by non-disclosers ($\Delta_{m,t}$) at the same moment.

We acknowledge that supply-side stability (the assumption that $\pi_m$, $\bar\theta_m$, and $c_{0,m}$ did not shift discontinuously at the reform) is maintained rather than directly tested. Our listing data begin in 2021Q3, so pre-reform ADEME match rates are not observable. The timing of audit decisions provides some reassurance: certified DPE audits are commissioned and completed before listing, with a gap of weeks to months, so the audit stock at any listing date reflects decisions made prior to that date. But this is a maintained assumption, and we treat it as such. The consequence of a supply-side shift at the reform is bounded: any resulting bias in $\hat\rho_m$ is upward, reinforcing the interpretation of $\hat\rho_m$ as an upper bound rather than a point estimate.\footnote{Because the disclosure obligation and its associated liability predate the reform (loi ALUR, Article L271-4 CCH), we treat $c_{0,m}$ as governed by that long-standing mandate and unchanged across July 2021. In the conservative case that opposability nonetheless nudged the perceived cost of withholding upward, $\hat\Delta_m$ would absorb both the demand response and the cost change, causing $\hat\rho_m = \hat\Delta_m / (\hat\beta_{1,m}\cdot\hat{e}^*_m)$ to overstate true buyer sensitivity. This is the same direction of bias as from supply-side instability: both reinforce the upper-bound interpretation.}

\subsubsection{DiD Identification of $\rho_m$}

The identification uses a panel of monthly hedonic price indices constructed by Chapter~1 from listing-level data covering all three disclosure groups. For each market $m$, group $g\in\{H,\,L,\,\emptyset\}$, and month $t$, the index $\bar{p}_{m,t}^g$ is the predicted log price per m$^2$ for a reference property conditional on hedonic controls $X_i$. The $H$ index covers energy-efficient disclosed properties (EPC grades A to D, $d=1$, $K=1$); the $L$ index covers energy-intensive disclosed properties (grades E to G, $d=1$, $K=1$); and the $\emptyset$ index covers all non-disclosing listings ($d=0$), a pooled group that blends strategic hiders ($K=1$) and genuinely uninformed sellers ($K=0$). By conditioning on observables, the cross-group gap $\bar{p}_{m,t}^g - \bar{p}_{m,t}^{g'}$ isolates the EPC-attributable price wedge net of compositional differences between groups.

From these indices, two time-varying log price gaps are defined:
\begin{align}
\beta_{1,m,t} &= \bar p_{m,t}^{H} - \bar p_{m,t}^{L} \quad \text{(energy-efficiency premium)} \\
\Delta_{m,t}  &= \bar p_{m,t}^{\emptyset} - \bar p_{m,t}^{L} \quad \text{(non-disclosure premium)}
\end{align}
The reform-induced shifts in each gap are estimated by differencing period means:
\begin{align}
\hat\beta_{1,m} &= \overline{\beta_{1,m,t}}^{\,\text{post}} - \overline{\beta_{1,m,t}}^{\,\text{pre}}
  \label{eq:did_beta1} \\
\hat\Delta_m    &= \overline{\Delta_{m,t}}^{\,\text{post}} - \overline{\Delta_{m,t}}^{\,\text{pre}}
  \label{eq:did_delta}
\end{align}
where the pre-period spans October~2018 through June~2021 (33~months) and the post-period spans July~2021 through July~2024 (37~months). A volume filter requiring at least 20 average monthly transactions per disclosure group per market retains 40 apartment markets and 100 house markets; markets with $\hat\beta_{1,m}\leq 0$ are excluded from structural inversion.

To validate the parallel pre-trends assumption and characterise the dynamic profile of the reform, the period-mean estimates are complemented with a month-by-month event study. For each market $m$ and relative month $k$ (where $k=0$ is July 2021), the normalised gaps
\begin{align}
\tilde\beta_{1,m,k} &= \beta_{1,m,t(k)} - \beta_{1,m,t(-1)}, \qquad
\tilde\Delta_{m,k}  = \Delta_{m,t(k)}  - \Delta_{m,t(-1)}
\end{align}
are set to zero at $k=-1$ (the last pre-reform month) by construction. Figure~\ref{fig:ATT} plots the cross-market means with 95\% confidence intervals over $k\in[-33,\,+36]$. The pre-reform coefficients ($k<0$) are statistically indistinguishable from zero for both gaps, confirming parallel pre-trends. The post-reform coefficients increase gradually from $k=0$, consistent with progressive capitalisation of the legally opposable DPE. The period-mean estimates in Equations~\eqref{eq:did_beta1} and \eqref{eq:did_delta} average these relative-month profiles over the respective windows.

Given these estimates, the structural equation $\Delta_m = \rho_m\cdot\beta_{1,m}\cdot e^*_m$ (Equation~\eqref{eq:premium}) implies a closed-form estimator for buyer sensitivity. Since the reform shifts $\rho_m$ while $e^*_m$ is a supply-side object measured directly from the ADEME registry (Section~\ref{sec:beliefs} below), the DiD isolates the post-reform demand response:
\begin{equation}
    \hat\rho_m = \frac{\hat\Delta_m}{\hat\beta_{1,m} \cdot \hat{e}^*_m}
    \label{eq:rho_estimator}
\end{equation}
The constraint $\hat\rho_m \leq 1$ binds in the majority of markets. For apartments, the unconstrained ratio exceeds unity in all 40 markets (cross-market mean of approximately 7.7); for houses, 92 of 107 markets with $\hat\beta_{1,m}>0$ are clipped (cross-market mean ratio of approximately 0.58). Two compounding sources drive the excess: composition bias in $\hat\Delta_m$, because the non-disclosure price index pools strategic hiders with genuinely uninformed sellers and inflates the numerator; and thin-cell attenuation of $\hat\beta_{1,m}$, because sparse disclosure-group cells compress the denominator. These effects are quantitatively consistent with the magnitudes observed, but they preclude point identification of $\rho_m$ in most markets. In clipped markets, $\rho_m = 1$ is maintained as a working assumption, and $\hat{c}_{0,m}$ is interpreted as an upper bound: since the indifference condition $c_{0,m} = \rho_m\cdot\beta_{1,m}\cdot\hat{e}^*_m(\hat\sigma_{0,m})$ is strictly increasing in $\rho_m$, setting $\rho_m = 1$ overstates $c_{0,m}$ relative to any interior value. Point identification is achieved in 8 house markets where the inversion yields an interior solution; these are characterised separately in the Results.

\subsubsection{Measuring Equilibrium Beliefs}
\label{sec:beliefs}

Computing $\hat\rho_m$ and $\hat{c}_{0,m}$ requires $\hat{e}^*_m$, the expected quality buyers assign to a non-disclosing listing. In standard signalling models this belief must be jointly estimated with the structural parameters. Here it is directly observable: the ADEME registry provides the researcher with the same information structure that $e^*_m$ encodes for buyers.

For each market $m$, the four-group classification of Section~\ref{sec:data} yields three sample frequencies: $\hat\pi_m$ (the share of listings with $K=1$, the ADEME match rate); $\hat{\bar\theta}_m$ (the share of informed sellers with high quality, $\theta=1$, EPC $\leq$ D); and $\hat\sigma_{0,m}$ (the share of informed low-quality sellers, $K=1$, $\theta=0$, who hide by setting $d=0$). Substituting into the belief function (Equation~\eqref{eq:belief_simplified}, at $\sigma_1^*=0$ per Corollary~\ref{cor:hq_disclose}):
\begin{equation}
    \hat{e}^*_m = \frac{(1-\hat\pi_m)\,\hat{\bar\theta}_m}{\hat\pi_m\,(1-\hat{\bar\theta}_m)\,\hat\sigma_{0,m} + (1-\hat\pi_m)}
    \label{eq:estar}
\end{equation}
The numerator is the uninformed high-quality mass in the non-disclosing pool; the denominator is the total non-disclosing mass. No model estimation is involved: $\hat{e}^*_m$ is a function of directly observed sample proportions.

\subsubsection{Closed-Form Identification of $c_{0,m}$}

One might expect $\hat{c}_{0,m} = \hat\Delta_m$, since the equilibrium indifference condition gives $c_{0,m} = \Delta_m$ for a marginal hider. This logic fails because the ``undisclosed'' category in the price index pools \emph{two heterogeneous groups}: strategic hiders ($K=1$, $\theta=0$, $d=0$), who are the economically relevant group, and genuinely uninformed sellers ($K=0$, $d=0$), who constitute approximately 85\% of all non-disclosers. $\hat\Delta_m$ is therefore a weighted average of the strategic hider premium and the uninformed-seller premium, an upper bound on $c_{0,m}$, not a point estimate.\footnote{A listing-level DiD targeting $K=1$ non-disclosers directly is infeasible: the listing data begin in 2021, providing no pre-reform baseline at the listing level.}

The route to $c_{0,m}$ uses the equilibrium hiding rate $\hat\sigma_{0,m}$ among informed low-quality sellers ($K=1$, $\theta=0$). This group is identified precisely by the ADEME match and is immune to contamination from uninformed sellers. In the Partial Pooling regime, the equilibrium indifference condition
\begin{equation}
    \Delta\!\left(e^*(\hat\sigma_{0,m})\,;\,\hat\rho_m\right) = c_{0,m}
    \label{eq:indifference}
\end{equation}
has a unique solution in $c_{0,m}$ for any observed $\hat\sigma_{0,m}\in(0,1)$. Given the first-stage estimate $\hat\rho_m$, this yields the closed-form estimator:
\begin{equation}
    \hat{c}_{0,m} = \hat\rho_m \cdot \hat\beta_{1,m} \cdot \hat{e}^*_m(\hat\sigma_{0,m})
    \label{eq:c0_estimator}
\end{equation}
where $\hat{e}^*_m(\hat\sigma_{0,m})$ is the belief evaluated at the observed mixing rate. The estimator is exact: one free parameter ($c_{0,m}$) is matched to one moment ($\hat\sigma_{0,m}$) via the equilibrium condition. No numerical optimisation is required.

For house markets, the data provide two distinct upper bounds on $c_{0,m}$. The first is the equilibrium inversion under the maintained $\hat\rho_m = 1$ assumption:
\begin{equation}
    \hat{c}_{0,m}^{\,\text{inv}} = \hat\beta_{1,m} \cdot \hat{e}^*_m(\hat\sigma_{0,m})
    \label{eq:c0_interval}
\end{equation}
The second is $\hat\Delta_m$, the DiD non-disclosure premium, which is an upper bound on $c_{0,m}$ because uninformed sellers (approximately 75\% of non-disclosers for houses) inflate the undisclosed pool's average price above what strategic hiders alone would generate. The cross-market means are $\hat{c}_{0,m}^{\,\text{inv}} = 0.018$ log points and $\hat\Delta_m = 0.034$ log points. Internal consistency requires $\hat{c}_{0,m}^{\,\text{inv}} \leq \hat\Delta_m$, which holds in all but a handful of markets. The tighter bound $\hat{c}_{0,m}^{\,\text{inv}} = 0.018$ is the preferred estimate and is used in the policy counterfactuals of Section~\ref{sec:policy}.

The structural estimates use listing-level data from 2022 to 2024.\footnote{The year 2021 is excluded from the listing-level analysis because the market had not yet fully repriced the reform: the pooled hedonic estimate of $\hat\beta_{1}$ for apartments in H2~2021 is approximately $-0.012$, indicating the EPC quality premium had not yet turned positive, which renders the structural equation $\Delta_m = \rho_m\cdot\beta_{1,m}\cdot e^*_m$ unreliable for that period.} The estimation sample comprises 40 apartment markets and 100 house markets.

Bootstrap standard errors are computed with $B=500$ within-market resamples. In each draw, listings in market $m$ are resampled with replacement and all market-level moments $(\hat\pi_m, \hat{\bar\theta}_m, \hat\sigma_{0,m})$ are recomputed. The first-stage estimate $\hat\rho_m$ is held fixed at its DiD value across bootstrap iterations; only $\hat\sigma_{0,m}$ is resampled to propagate sampling uncertainty into $\hat{c}_{0,m}$ via Equation~\eqref{eq:c0_estimator}. The 500 replications yield percentile-based 95\% confidence intervals.

\subsection{Estimates}
\label{sec:results}

\subsubsection{Main Estimates}

Table~\ref{tab:did_results} reports market-level estimates for both property types. Panel A covers the DiD estimates from the price index panel; Panel B covers the two-step structural estimates; Panel C covers supply-side observables from the ADEME match.

\begin{table}[H]
\centering
\caption{Market-level estimates by property type}
\label{tab:did_results}
\small\renewcommand{\arraystretch}{1.3}
\begin{tabular*}{\textwidth}{@{\extracolsep{\fill}}l*{6}{c}@{}}
\toprule\toprule
 & \multicolumn{3}{c}{\textbf{Apartments}} & \multicolumn{3}{c}{\textbf{Houses}} \\
\cmidrule(lr){2-4}\cmidrule(lr){5-7}
\textbf{Parameter} & \textbf{Mean} & \textbf{Std} & \textbf{N} & \textbf{Mean} & \textbf{Std} & \textbf{N} \\
\midrule
\multicolumn{7}{l}{\textit{Panel A: DiD estimates from price index panel}} \\[2pt]
$\hat\beta_{1,m}$ (energy-efficiency premium) & 1.7\% & 3.5 & 49  & 5.1\% & 4.0 & 104 \\
$\hat\Delta_m$ (non-disclosure premium, UB)   & 6.9\% & 6.3 & 49  & 3.4\% & 4.8 & 104 \\
\midrule
\multicolumn{7}{l}{\textit{Panel B: Two-step structural estimates, pooled 2022 to 2024}} \\[2pt]
$\hat\rho_m$ (buyer EPC-sensitivity)          & 90.9\% & 28.6 & 45 & 68.1\% & 42.0 & 97 \\
$\hat c_{0,m}$ (hiding cost)                  &  3.4\% &  2.9 & 45 &  4.5\% &  3.7 & 97 \\
\midrule
\multicolumn{7}{l}{\textit{Panel C: Supply-side observables from ADEME match}} \\[2pt]
$\hat\pi_m$ (registry-matched share)          & 50.7\% &  7.4 & 45 & 41.9\% &  4.4 & 97 \\
$\hat{\bar\theta}_m$ (mean quality $K{=}1$)   & 62.5\% &  8.8 & 45 & 60.5\% & 10.4 & 97 \\
$\hat\sigma_{0,m}$ (informed hiding rate)     & 14.0\% & 10.5 & 45 &  5.7\% &  6.5 & 97 \\
\bottomrule\bottomrule
\end{tabular*}
\vspace{0.5em}
\begin{minipage}{\linewidth}
\footnotesize\textit{Notes.} Panel A: market-level DiD estimates from the price index panel (pre: Oct 2018 to Jun 2021; post: Jul 2021 to Jul 2024), volume-filtered at $\geq 20$ transactions/month/EPC category. $N=43$ apartment markets after excluding 25 with $\hat\beta_{1,m}\leq 0$; $N=107$ house markets after excluding 6. $\hat\Delta_m$ is an upper bound for $c_{0,m}$ because the undisclosed category pools strategic hiders and uninformed sellers. Panel B: two-step structural estimates using 2022 to 2024 listing data. $\hat\rho_m = \hat\Delta_m / (\hat\beta_{1,m} \cdot \hat{e}^*_m)$ from closed-form DiD inversion, clipped to $[0,1]$; $\hat{c}_{0,m}$ is recovered by equilibrium inversion from $\hat\sigma_{0,m}$ with $\hat\rho_m$ maintained at the clipped value (upper bound when $\hat\rho_m = 1$, since $c_{0,m} = \rho_m\cdot\hat\beta_{1,m}\cdot\hat{e}^*_m$ is increasing in $\rho_m$). Panel C: supply-side moments computed as sample frequencies from the ADEME match. Means in Panel~A and the $\hat c_{0,m}$ row of Panel~B are percentage changes (computed as $e^\delta - 1$, expressed in percent). Share parameters in Panels~B and~C ($\hat\rho_m$, $\hat\pi_m$, $\hat{\bar\theta}_m$, $\hat\sigma_{0,m}$) are in percentage points. Standard-deviation columns are in the same units as the mean (percentage points).
\end{minipage}
\end{table}

\vspace{-1em}

Table~\ref{tab:did_results} yields a coherent and interpretable picture across all three panels.
In Panel~A, the energy-efficiency premium $\hat\beta_{1,m}$ is approximately 1.7\% for
apartments and 5.1\% for houses: disclosing an energy-efficient EPC score commands that price advantage
over disclosing an energy-intensive score, after the July~2021 reform. The non-disclosure premium
$\hat\Delta_m$ (6.9\% for apartments and 3.4\% for houses) measures how much more
non-disclosers earn over low-quality disclosed properties compared with the pre-reform
baseline; because the undisclosed pool commingles strategic hiders with genuinely uninformed
sellers, this figure is an upper bound on the true concealment cost rather than a direct
estimate of it. Panel~B translates these reduced-form inputs into seller-side parameters: $\hat\rho_m = 90.9\%$ for apartments and 68.1\% for
houses indicate that nearly all buyers in the average market price energy performance into
their offers, while the inverted hiding cost $\hat{c}_{0,m}$ settles at 3.4\% of property value for
apartments and 4.5\% for houses --- an economically substantial margin, and one of the same
order in both segments despite how differently they are traded. Panel~C
illuminates the supply side: just over half of apartment sellers (50.7\%) and 41.9\% of house
sellers carry a post-2021 registry record at listing ($\hat\pi_m$), and among those, roughly
six in ten hold a property rated A--D ($\hat{\bar\theta}_m = 62.5\%$ and $60.5\%$).
The most striking contrast across property types lies in the strategic-hiding rate
$\hat\sigma_{0,m}$: 14.0\% of informed apartment sellers who hold an E/F/G rating choose to
withhold it, versus 5.7\% for houses, a difference that mirrors the three-fold gap in
$\hat\beta_{1,m}$, since a larger forgone quality premium makes concealment a costlier choice
for house sellers.

\subsubsection{Structural Estimates and Economic Interpretation}

The first-stage estimator $\hat\rho_m = \hat\Delta_m / (\hat\beta_{1,m} \cdot \hat{e}^*_m)$, clipped to $[0,1]$, yields means of 0.909 for apartments and 0.681 for houses. These averages, however, conceal a pattern of boundary saturation. For apartments, 40 of the 42 partial-pooling markets reach the $\hat\rho_m = 1$ upper boundary; only the 3 full-unraveling markets pull the mean below unity. For houses, 54 of 76 partial-pooling markets reach the boundary, with 21 full-unraveling markets contributing $\hat\rho_m = 0$ and 22 markets achieving interior identification. The boundary reflects a structural feature of the data: the non-disclosure premium $\hat\Delta_m$ typically exceeds $\hat\beta_{1,m} \cdot \hat{e}^*_m$, so the data are consistent with full buyer EPC-sensitivity but cannot pin down exactly how close to one $\rho_m$ truly is. We maintain $\hat\rho_m = 1$ as the working assumption for these markets and treat $\hat{c}_{0,m}$ as an upper bound on the true hiding cost accordingly.

Under this convention, the concealment cost is economically material: 0.034 log points (approximately 3.4\% of property value) for apartments and 0.045 log points (4.5\%) for houses. That the two segments land within roughly one percentage point of each other suggests the cost is governed primarily by the legal disclosure regime and the state of registry coverage rather than by property-specific characteristics. The reduced-form premium provides an independent check: $\hat c_{0,m} \leq \hat\Delta_m$ holds market by market by construction, with equality exactly in the interior markets, so the inversion cannot exceed the bound the DiD imposes on it.

The contrast between property types is itself informative. Houses have an energy-efficiency premium three times larger than apartments ($\hat\beta_{1,m} \approx 0.051$ vs.\ 0.017 log points), reflecting their larger surface areas and older building stock, which make heating costs directly material to buyers' running expenses. A higher forgone premium makes concealment more expensive for low-quality sellers, and the data reflect this exactly: the hiding rate $\hat\sigma_{0,m}$ is well under half as large for houses (5.7\% vs.\ 14.0\%), and the share of strategic hiders is correspondingly lower (1.1\% vs.\ 3.5\% of all listings). The non-disclosure premium $\hat\Delta_m$ is also lower for houses (0.034 vs.\ 0.069), since a thinner pool of strategic hiders means buyers' beliefs about withheld quality are less pessimistic. The entire pattern (higher quality premium, lower hiding rate, lower non-disclosure premium) is the qualitative prediction of the model under equilibrium comparative statics across property types.

One further check concerns the gap between the non-disclosure premium $\hat\Delta_m$ and the structural hiding cost $\hat{c}_{0,m}$. In the apartment markets where both are available, the raw DiD estimate exceeds the inverted hiding cost by a factor of roughly two. This systematic gap is not a contradiction: it is what the model predicts. Uninformed sellers (approximately 83\% of non-disclosers for apartments and 89\% for houses) earn a non-disclosure premium simply by being in the market, pooling with strategic hiders and inflating the undisclosed category's average price above what strategic concealment alone would generate. $\hat\Delta_m$ is the market-level average of this pooled premium; $\hat{c}_{0,m}$ is the cost borne by the marginal informed hider, recovered by targeting only the informed-seller hiding rate $\hat\sigma_{0,m}$ in the equilibrium inversion. The two objects measure different things, and the fact that $\hat\Delta_m$ greatly exceeds $\hat{c}_{0,m}$ is precisely the compositional mechanism the model is designed to capture.


% ============================================================
% Section 4.6 -- Limited Attention to Silence (extension)
% DRAFT inserted 2026-10-08 from playground/notes.md Part IV (plan harmonic-munching-lecun.md).
% Every number is a macro generated by scripts/compute_attention.py and
% scripts/mc_attention_estimator.py; values are provisional until the certificate-date (G2)
% and certified-consumption (G3) re-pulls. Paragraphs tagged [AUTHOR REWRITE] are
% load-bearing and should be rewritten in the author's voice (writing-with-ai.md).
% ============================================================

\subsection{Limited Attention to Silence}
\label{ss:attention}

% [AUTHOR REWRITE] opening
The estimates in Section~\ref{sec:results} leave one result unexplained, and its shape is specific enough to point toward a mechanism. In apartment markets the reform raised the price of a listing that shows no rating, relative to one rated E to G, by more than it raised the price of a verified A-to-D listing. We call a listing that shows no rating \emph{silent}; in the notation of Section~\ref{ss:model_setup} it is a non-disclosing listing, $d=0$. Averaged over the \attNMarketsApt{} apartment markets of Table~\ref{tab:did_results}, the two shifts are \attDthreeApt{} and \attBthreeApt{} log points, and the first exceeds the second in \attNDgtBApt{} of them. In house markets the order is reversed on average, \attDthreeHouse{} against \attBthreeHouse{}. This subsection asks what would have to be true of buyers for silence to be worth more than a verified efficient rating.

One answer is that buyers attend to energy when a listing puts it in front of them, and much less when it does not. We allow for this mechanism without imposing it: full attention remains a special case, and the data decide between the two. The extension changes the buyer side of the model and nothing else. It nests the model of Sections~\ref{ss:model_setup}--\ref{sec:equilibrium} as a special case and preserves its regimes. It leaves $\hat\rho_m$ unchanged as a number but reads it as effective attention to displayed ratings, and it reads the concealment cost off the non-disclosure premium rather than off $\hat\beta_{1,m}\hat e^*_m$ (Proposition~\ref{prop:att_identification}). It also turns the over-identifying wedge $\hat\lambda_m$ of Section~\ref{sec:identification} into an estimate of how much attention silent listings receive.

\subsubsection{Motivation}
\label{ss:att_motivation}

\paragraph{Why the model above cannot reach the apartment wedge.}
The difficulty is not one of calibration. In the model of Section~\ref{ss:model_setup}, a buyer who sees no rating prices the listing at her belief about its quality, and a belief about a binary rating lies between the worst rating and the best. A silent listing therefore sells at most at the price of a verified efficient listing: the ratio of the non-disclosure premium to the energy-efficiency premium is the pool belief, which is bounded by $\bar\theta\le1$ (Lemma~\ref{lem:model_bounds}(i)). The same bound applies to the differences-in-differences whenever the pool belief is stable across the reform.\footnote{The ordering survives exact pricing under the monotonicity condition discussed in Appendix~\ref{app:att_exact}. The levels make the same point: silent apartments sold above verified A-to-D listings by \attLevelGapPreApt{} log points before the reform and by \attLevelGapPostApt{} after, whereas the model prices silence below verified efficiency at any date and for any belief in $[0,1]$. Under Assumption~\ref{assum:no_pre_pricing} the pre-reform gap is non-energy, and the energy component of the post-reform gap is the DiD, \attExcessApt{} log points, still positive.} The apartment wedge of \attLamApt{} log points decomposes accordingly. Of it, \attExcessApt{} points are the amount by which $\hat\Delta$ exceeds $\hat\beta_1$, a region the model cannot enter. Only the remaining \attReachApt{} points, $(1-\hat e^*)\hat\beta_1$, are within reach of any change in how buyers form beliefs about silent listings.

Figure~\ref{fig:att_belief} shows that most apartment markets lie above the ceiling that the core model imposes, whatever belief buyers hold about silent listings.

\begin{figure}[htbp]
    \centering
    \caption{The core model cannot produce the non-disclosure premium}
    \label{fig:att_belief}
    \includegraphics[width=\linewidth]{fig/visuals/fig_att_F2a_belief.pdf}
    \begin{minipage}{0.95\linewidth}
        \vspace{0.5em}
        \footnotesize\textit{Notes.} Each point is a market; both axes are log points. In the core model a silent listing sells at the buyers' belief about its quality times the energy-efficiency premium $\hat\beta_{1,m}$, and a belief about a binary rating lies between zero and one, so the non-disclosure premium $\hat\Delta_m$ cannot exceed $\hat\beta_{1,m}$ (the line). Markets in the shaded region, drawn as filled points, are outside what the core model can produce for any belief about silent listings: \attNDgtBApt{} of \attNMarketsApt{} apartment markets and \attNDgtBHouse{} of \attNMarketsHouse{} house markets. Sampling error in $\hat\beta_{1,m}$ and $\hat\Delta_m$ is not shown.
    \end{minipage}
\end{figure}

\paragraph{Belief-based amendments cannot reach the wedge.}
This rules out the amendments that come to mind first. First, buyers may draw too little inference from silence, as in cursed equilibrium and in experimental evidence on disclosure games \citep[e.g.,][]{eyster_cursed_2005, jin_is_2021, hagenbach_motivated_2025}. Any such rule replaces the Bayesian belief $e^*$ by another number between zero and one, so it stays below the ceiling. It is also quantitatively small here. Sellers without a certificate make up about \attShareUninfApt{} percent of the silent apartment pool, so the Bayesian belief already sits within a few points of the prior, and full cursedness would account for about \attCursedApt{} of the \attLamApt{} points. Second, buyers may over-react to a displayed poor rating. That lowers the price of E-to-G listings, which is the base of both gaps, and it cancels from the comparison between silent and verified efficient listings. Third, the distrust discount on inferred quality allowed in Section~\ref{ss:model_setup} ($\alpha_k<\beta_k$) moves the price of silence down, not up. Observable composition does not resolve the puzzle either, because the indices hold the hedonic reference property fixed within each market and the comparison is a difference-in-differences. What these arguments do not exclude is a change in the \emph{unobserved} composition of the silent pool, most of which consists of sellers without a certificate: if the dwellings that were silent after the reform were of better unobserved quality, the wedge would rise without any change in attention.


% [AUTHOR REWRITE] mechanism paragraphs
What these amendments share is that the buyer facing a silent listing still thinks about its energy performance; she only reasons about it imperfectly. A large body of evidence suggests that consumers often do not think about costs that are not put in front of them, and that when they do not, they act as if those costs were small rather than average. Buyers under-react to sales taxes that are not included in posted prices \citep{chetty_salience_2009}, to road tolls collected electronically \citep{finkelstein_eztax_2009} and to shrouded shipping fees \citep{brown_shrouded_2010}. Energy operating costs are a documented instance in durable-goods markets \citep[e.g.,][]{allcott_consumers_2011, sallee_rational_2014, allcott_evaluating_2015}. Models of focusing, salience and sparsity formalise the same idea, that the attributes a choice set makes salient receive more weight than those it leaves in the background \citep[e.g.,][]{koszegi_model_2013, bordalo_salience_2013, gabaix_sparsity_2014}. In financial markets, limited attention changes how disclosed and undisclosed information is priced \citep[e.g.,][]{hirshleifer_limited_2003, dellavigna_investor_2009}.

Three features of the French setting make this behaviour plausible for apartment buyers in particular. First, a listing that discloses shows the rating as a coloured label and, from January 2022, an estimate of the annual energy bill in euros, so disclosure itself puts energy into the buyer's comparison of dwellings.\footnote{Since 1 January 2022, advertisements for the sale of a dwelling must state the estimated annual energy expenditure for standard use, with the reference year of energy prices, and dwellings rated F or G must carry the mention ``logement \`a consommation \'energ\'etique excessive''. Omitting these mentions exposes the advertiser to fines under Article L126-33 of the Code de la construction et de l'habitation. % [CHECK] cite the implementing decree (D\'ecret n\textsuperscript{o} 2020-1609, 17 Dec.\ 2020) against L\'egifrance before submission.
The obligation also gives a dated test of the mechanism: under limited attention the gap between silent and rated listings should widen again at that date.} A listing that shows nothing does not. Second, heating in many apartment buildings is collective and paid through service charges, so apartment buyers are likely to see energy as part of a monthly charge rather than as an attribute of the dwelling. Third, a missing rating is common: about \attShareUninfApt{} percent of silent apartment listings match no post-reform certificate in the registry. Silence is therefore unremarkable rather than conspicuous.

The mechanism is easiest to see in three steps. Suppose first that buyers attend to energy as often for a silent listing as for a rated one. They then price silence at the expected cost of the silent pool, and the model is exactly the one estimated above. Suppose next that they attend half as often. Half of them still price silence at the expected cost of the pool, while the other half price it at whatever they implicitly assume when energy never comes up. If that default is ``nothing out of the ordinary'', the silent listing gains on every rated listing, efficient ones included. Suppose finally that buyers never attend to energy on a silent listing. It is then priced as if it carried the default cost, which can lie below the cost of an A-rated dwelling, and silence beats every rating.

\subsubsection{Setup}
\label{ss:att_setup}

We embed display-dependent attention in the buyer side of Section~\ref{ss:model_setup}. This lets us read the over-identifying wedge as a measure of the attention silent listings receive. The estimators of Section~\ref{sec:identification} keep their numerical values, but under the extension they measure different objects (Proposition~\ref{prop:att_identification}). Everything not stated below is as in Section~\ref{ss:model_setup}: seller types, the No Falsification assumption, Assumption~\ref{ass:independence}, the payoff~\eqref{eq:payoff} and the demand system~\eqref{eq:demand}.

\paragraph{Environment and timing.}
The environment is that of Section~\ref{ss:model_setup}, with one addition on the buyer side: after seeing the listing, an EPC-sensitive buyer either attends to its energy performance or does not, with a probability that depends on what the listing shows. The timing is as follows.
\begin{enumerate}[label=(\arabic*)]
    \item Nature draws the seller's quality $\theta$ and information status $K$, independently, with $\Pr(\theta=1)=\bar\theta$ and $\Pr(K=1)=\pi$.
    \item The seller lists the property, choosing the display $d$ (only if $K=1$) and the listing price.
    \item A buyer arrives. She is EPC-sensitive with probability $\rho$. If sensitive, she attends to energy with a probability $\omega_d$ that depends on the display.
    \item She values the property given what she attends to, and buys or not.
\end{enumerate}

\paragraph{Payoffs and hiding costs.}
The seller's problem is~\eqref{eq:payoff}, with hiding cost $c(\theta,d)=\mathbf 1[d=0](c_0+\theta c_1)$, interpreted as in Section~\ref{ss:model_setup}. The extension leaves the seller side of the model unchanged. It changes only how buyers price what a listing shows, and therefore what a seller gains by showing nothing.

\paragraph{Buyer valuations and demand.}
Section~\ref{ss:model_setup} writes the valuation of an EPC-sensitive buyer in what we call \emph{bonus form}: she adds $\beta_1$ times her belief that the dwelling is efficient. This form records the difference in energy value between efficient and inefficient dwellings, $\beta_1$, but never the level of the energy cost of either, and it has no place for a buyer who forms no belief at all, because every sensitive buyer enters through some number $\mathbb E[\theta\mid\cdot\,]\in[0,1]$. The ceiling documented in Section~\ref{ss:att_motivation} is thus a property of the representation as much as of the economics. To ask what a buyer pays for a listing whose energy performance she does not consider, we rewrite the valuation in \emph{cost form}. The rewriting changes nothing in the model of Section~\ref{ss:model_setup} (Lemma~\ref{lem:cost_form}); it changes what the model can express. Appendix~\ref{app:attention} explains why the normalisation implicit in the bonus form becomes binding once attention depends on what a listing shows.

A dwelling of quality $\theta$ in market $m$ carries a capitalised energy cost $C_{\theta,m}$, with $C_{0,m}>C_{1,m}\ge0$, measured in the units of the price gaps; we drop the market index when no confusion arises. We take these costs from the construction of the anchor in Assumption~\ref{assum:energy_anchor}, $C_{\theta,m}=\kappa\,p^E a(r,T)\,\bar E_{\theta,m}/P_m$, so that $\beta^{\rm anchor}_{1,m}=C_{0,m}-C_{1,m}$. The extension also uses the cost \emph{level}. Because the conversion factor $\kappa p^Ea(r,T)/P_m$ is common to both qualities, any energy cost below can equivalently be expressed in kWh$_{\rm ep}$/m$^2$/yr. Define the \emph{energy-cost ratio}
\begin{equation}
    \eta_m\;\equiv\;\frac{C_{0,m}}{C_{0,m}-C_{1,m}}\;=\;\frac{\bar E_{0,m}}{\bar E_{0,m}-\bar E_{1,m}}\;>\;1,
    \label{eq:eta}
\end{equation}
which is computed from the registry alone and does not depend on the energy price, the discount rate or the horizon.

\begin{definition}[Gross Value and Cost-Form Valuation]
\label{def:cost_form}
The \emph{gross value} of a dwelling with attributes $X$ is
\begin{equation}
    \bar V(X)\;\equiv\;\beta_0+\sum_n\beta_nX_n+C_0,
    \label{eq:gross_value}
\end{equation}
its value before any energy cost is subtracted. A buyer's \emph{cost-form valuation} of a listing is
\begin{equation}
    u\;=\;\bar V(X)-\mathcal C,
    \label{eq:cost_valuation}
\end{equation}
where $\mathcal C$ is the energy cost she subtracts, given what the listing shows and what she attends to.
\end{definition}

The gross value adds back $C_0$ because, in the bonus form, $\beta_0+\sum_n\beta_nX_n$ is not the value of a dwelling without an energy bill. It is the value of an inefficient dwelling, whose energy cost is already netted out of the intercept. Adding $C_0$ moves the zero point of the valuation from the worst dwelling to a dwelling with no energy cost, so that every dwelling, inefficient ones included, carries its energy cost explicitly. The level of $\bar V(X)$ never matters for a buyer who treats every listing alike, because it cancels from every price gap; what the cost form adds is the ability to compare what a buyer subtracts on one listing with what she subtracts on another.

\begin{lemma}[Bonus and Cost Forms]
\label{lem:cost_form}
The cost-form valuation~\eqref{eq:cost_valuation} reproduces the valuations of Section~\ref{ss:model_setup} exactly, with $\beta_1=C_0-C_1$. An EPC-sensitive buyer subtracts $\mathcal C=C_\theta$ for a disclosed rating of quality $\theta$ and $\mathcal C=\bar C^{\rm sil}$ for a silent listing, where
\begin{equation}
    \bar C^{\rm sil}\;\equiv\;\mathbb E[C\mid d=0]\;=\;e\,C_1+(1-e)\,C_0
    \label{eq:csil}
\end{equation}
is the expected energy cost of the silent pool under the belief $e$ of~\eqref{eq:belief_general}. An EPC-insensitive buyer subtracts $\mathcal C=C^{\rm ins}\equiv\bar V(X)-v_0(X)$, the same amount for every listing with attributes $X$. Moreover, $C_0$ is the only constant which, added to $\beta_0+\sum_n\beta_nX_n$, makes the two forms agree for both qualities.
\end{lemma}

\noindent\textit{Proof.} See Appendix~\ref{app:attention}. $\Box$

\begin{definition}[Default Energy Cost]
\label{def:default_cost}
The \emph{default energy cost} $C^{\rm def}\ge0$ is the energy cost an EPC-sensitive buyer subtracts from a listing's value when she does not attend to energy. It is measured in the same units as $C_\theta$.
\end{definition}

The term follows the limited-attention literature, in which an agent who does not attend to an attribute replaces it with a default value rather than with an inference \citep{gabaix_sparsity_2014, gabaix_behavioral_2019}. The default is therefore not a belief about the dwelling's rating. It is what the buyer implicitly assumes when the question of energy never arises. Three benchmark values fix ideas. With $C^{\rm def}=0$ no energy cost is priced at all, which does not mean the buyer believes the dwelling consumes nothing, only that energy plays no part in her valuation. With $C^{\rm def}$ equal to the prior mean cost $\bar\theta C_1+(1-\bar\theta)C_0$, the inattentive buyer behaves like one who assumes an average dwelling, which is cursedness by another name. With $C^{\rm def}=\bar C^{\rm sil}$ she happens to subtract exactly what an attentive buyer would infer, and attention becomes irrelevant. The same default applies to rated and silent listings, because a buyer who does not attend to energy, by definition, does not use what the listing shows about it. Which value is relevant is an empirical question: Proposition~\ref{prop:att_incentives}(iii) shows that the extension can explain the apartment wedge only if $C^{\rm def}<C_1$, and Proposition~\ref{prop:att_identification}(iii) shows what the data can say about its level.

\begin{assumption}[Display-Dependent Attention]
\label{assum:attention}
An EPC-sensitive buyer attends to energy with probability $\omega_1$ when the listing displays a rating and with probability $\omega_0\le\omega_1$ when it does not. An attentive buyer subtracts the expected energy cost given what the listing shows: $C_\theta$ if a rating is displayed, $\bar C^{\rm sil}$ otherwise. An inattentive buyer subtracts the default energy cost $C^{\rm def}$ of Definition~\ref{def:default_cost}, whatever the listing shows.
\end{assumption}

We call $\psi\equiv\omega_0/\omega_1\in[0,1]$ the \emph{relative attention to silence}: the attention a listing without a rating receives, per unit of the attention a rated listing receives. We call $\rho\,\omega_d$ the \emph{effective attention} to listings with display $d$, the share of all buyers who price energy when they look at such a listing. Table~\ref{tab:att_buyers} records what each buyer subtracts.

\begin{table}[htbp]
\centering
\caption{Energy Cost Subtracted, by Buyer Type and Listing Display}
\label{tab:att_buyers}
\small\renewcommand{\arraystretch}{1.5}
\begin{tabular*}{\textwidth}{@{\extracolsep{\fill}}lccc@{}}
\toprule\toprule
\textbf{Buyer} & \textbf{Probability} & \textbf{Displayed rating $\theta$} & \textbf{Silent listing} \\
\midrule
EPC-insensitive & $1-\rho$ & $C^{\rm ins}$ & $C^{\rm ins}$ \\
EPC-sensitive, attentive & $\rho\,\omega_d$ & $C_\theta$ & $\bar C^{\rm sil}$ \\
EPC-sensitive, inattentive & $\rho\,(1-\omega_d)$ & $C^{\rm def}$ & $C^{\rm def}$ \\
\bottomrule\bottomrule
\end{tabular*}

\vspace{0.5em}
\raggedright
\footnotesize
\textit{Notes:} $\omega_d=\omega_1$ for a displayed rating and $\omega_0$ for a silent listing. $C^{\rm ins}\equiv\bar V(X)-v_0(X)$ (Lemma~\ref{lem:cost_form}) does not depend on the rating or on whether one is displayed, so it cancels from every price gap.
\end{table}

Each buyer values a listing by~\eqref{eq:cost_valuation}, with $\mathcal C$ given by Table~\ref{tab:att_buyers}. The transaction probability generalises~\eqref{eq:demand} to three buyer segments,
\begin{equation}
    \Gamma(p^{\rm list}\mid d)=\rho\Big[\omega_d\,G\big(\bar V-\mathcal C_d-p^{\rm list}\big)+(1-\omega_d)\,G\big(\bar V-C^{\rm def}-p^{\rm list}\big)\Big]+(1-\rho)\,G\big(v_0-p^{\rm list}\big),
    \label{eq:demand_att}
\end{equation}
where $\mathcal C_d$ is $C_\theta$ for a displayed rating and $\bar C^{\rm sil}$ for a silent listing. All results below are stated under the linear pricing approximation used to derive~\eqref{eq:premium_derived}, in which the equilibrium price equals the buyer-share-weighted mean valuation. Whether their orderings survive the exact demand~\eqref{eq:demand_att} is discussed, not proven, in Appendix~\ref{app:attention}.

Write $p_{\rm H}$, $p_{\rm L}$ and $p_{\rm sil}$ for the equilibrium log prices, at the reference property $X$, of a listing displaying an efficient rating (A to D), a listing displaying an inefficient rating (E to G), and a silent listing. In the notation of Section~\ref{sec:equilibrium}, $p_{\rm sil}-p_{\rm L}$ is the premium $p_0^*-p_1^*$.

\begin{lemma}[Prices and the Perceived Cost of Silence]
\label{lem:att_prices}
Under Assumption~\ref{assum:attention}, the energy-efficiency premium $b$ and the non-disclosure premium $\Delta$ are
\begin{align}
    b&\;\equiv\;p_{\rm H}-p_{\rm L}\;=\;\rho\,\omega_1\,\beta_1, \label{eq:b_att}\\
    \Delta&\;\equiv\;p_{\rm sil}-p_{\rm L}\;=\;\rho\,\omega_1\,(C_0-C^*),\qquad C^*\;\equiv\;C^{\rm def}+\psi\,\big(\bar C^{\rm sil}-C^{\rm def}\big). \label{eq:premium_att}
\end{align}
More generally, a listing displaying a rating with energy cost $C$ sells at $p_{\rm sil}-\rho\,\omega_1\,(C-C^*)$.
\end{lemma}

\noindent\textit{Proof.} See Appendix~\ref{app:attention}. $\Box$

$C^*$ is the energy cost the market behaves as if a silent listing carries: a weighted average of the default and the rational cost of the silent pool, with weight $\psi$ on the latter. The default cancels from the energy-efficiency premium because both listings in that comparison display a rating; it does not cancel from the non-disclosure premium, because a silent listing draws attention $\omega_0$ while a rated listing draws $\omega_1$. We keep $\beta_1=C_0-C_1$ for the valuation coefficient of Section~\ref{ss:model_setup} and write $b$ for the premium. With $\omega_1=1$, $b=\rho\beta_1$, the premium of Section~\ref{sec:equilibrium}, whose reform-induced change $\hat\beta_{1,m}$ estimates.

\paragraph{Beliefs and the effective belief.}
Of the three buyer types in Table~\ref{tab:att_buyers}, only the attentive EPC-sensitive buyer forms a belief about a silent listing. Her belief is the Bayesian $e(\sigma_0,\sigma_1)$ of~\eqref{eq:belief_general}, unchanged: the share of efficient dwellings among silent listings. The insensitive buyer does not care about energy, and the inattentive sensitive buyer would care but does not think about it on this listing, so she subtracts the default instead. The Bayesian belief therefore reaches prices only through attentive buyers.

In Section~\ref{sec:equilibrium} everything the market thinks about silence is summarised by one number, because the ratio of the non-disclosure premium to the energy-efficiency premium \emph{is} the belief. Under limited attention prices aggregate a belief held by some buyers and a default used by others, and no single buyer's belief summarises them. We define the number that plays the role $e$ played before.

\begin{definition}[Effective Belief]
\label{def:effective_belief}
The \emph{effective belief} is the ratio of the non-disclosure premium to the energy-efficiency premium, $\tilde e\equiv\Delta/b$, so that $\Delta=b\,\tilde e$ as $\Delta=\rho\beta_1 e$ in Section~\ref{sec:equilibrium}.
\end{definition}

Equivalently, $\tilde e$ is the belief that rationalises observed prices within the model of Section~\ref{sec:equilibrium}: if every EPC-sensitive buyer attended to energy and held belief $\tilde e$ about every silent listing, the non-disclosure premium would be exactly $\Delta$. That model confines beliefs to $[0,1]$. An effective belief above one is therefore a direct statement that no belief, Bayesian or behavioural, can account for the premium in that model, and the gap between $\tilde e$ and the registry-measured Bayesian belief $e^*$ is precisely what limited attention has to explain.

The effective belief follows in four steps. By Lemma~\ref{lem:att_prices}, the factor $\rho\omega_1$ cancels from the ratio, so $\tilde e=(C_0-C^*)/(C_0-C_1)$. This places the perceived cost of silence on a \emph{quality scale}, $q(C)\equiv(C_0-C)/(C_0-C_1)$, on which an inefficient dwelling's cost maps to 0 and an efficient one's to 1. Because $q$ is linear and $C^*=\psi\,\bar C^{\rm sil}+(1-\psi)\,C^{\rm def}$, the effective belief is the same weighted average of the two points, $\tilde e=\psi\,q(\bar C^{\rm sil})+(1-\psi)\,q(C^{\rm def})$. By~\eqref{eq:csil}, $q(\bar C^{\rm sil})=e$: the cost an attentive buyer subtracts maps back exactly to her Bayesian belief. The point at which an inattentive buyer places every listing is the \emph{default-implied belief}
\begin{equation*}
    e^{\rm def}\;\equiv\;q\big(C^{\rm def}\big)\;=\;\frac{C_0-C^{\rm def}}{C_0-C_1}\;=\;\eta-\frac{C^{\rm def}}{C_0-C_1},
\end{equation*}
which equals $\eta$ when $C^{\rm def}=0$. Combining the steps,
\begin{equation}
    \tilde e\;=\;\psi\,e+(1-\psi)\,e^{\rm def}.
    \label{eq:etilde}
\end{equation}

The effective belief is a weighted average of the belief of buyers who think about silence and the default of those who do not, with weight $\psi$ on the belief. The two pieces live on different ranges. The Bayesian belief is a probability, so $e\le1$. The default-implied belief is not: it exceeds one whenever $C^{\rm def}<C_1$, that is, whenever a buyer who ignores energy implicitly assumes a lower cost than that of an average efficient dwelling. A buyer who prices no energy cost at all treats a silent listing as better than a verified efficient one, because even a verified efficient listing shows a positive cost. With $\psi=1$ the model is that of Section~\ref{sec:equilibrium}; with $\psi=0$ the composition of the silent pool plays no role in prices; and $\tilde e>1$, hence $\Delta>b$, is possible as soon as $\psi<1$ and $e^{\rm def}>1$.

\begin{lemma}[Effective-Belief Monotonicity]
\label{lem:etilde_monotone}
$\tilde e(\sigma_0,\sigma_1)$ is weakly decreasing in $\sigma_0$ and weakly increasing in $\sigma_1$, strictly so if and only if $\psi>0$, with $\partial\tilde e/\partial\sigma_j=\psi\,\partial e/\partial\sigma_j$.
\end{lemma}

\noindent\textit{Proof.} In~\eqref{eq:etilde} only $e$ depends on $(\sigma_0,\sigma_1)$. Differentiate and apply Lemma~\ref{lem:belief_monotone}. $\Box$

In Section~\ref{sec:equilibrium} the hiding rate is pinned down by a feedback loop: more concealment worsens the pool, the worse pool lowers the belief, and the lower belief lowers the gain from further concealment. Lemma~\ref{lem:etilde_monotone} says that under limited attention this loop runs at a fraction $\psi$ of its full strength. Inattentive buyers keep subtracting the default whatever the pool contains.

\subsubsection{Equilibrium}
\label{ss:att_equilibrium}

\paragraph{Disclosure incentives.}
An informed seller compares the price of silence with the price of her own displayed rating. By Lemma~\ref{lem:att_prices}, the gross gain from hiding for an informed seller of quality $\theta$ is
\begin{equation}
    g_0\;=\;b\,\tilde e,\qquad g_1\;=\;b\,(\tilde e-1),\qquad g_0-g_1=b.
    \label{eq:gains}
\end{equation}
The low type's gain is the non-disclosure premium itself, $g_0=\Delta$. Section~\ref{sec:equilibrium} writes this gain as $\Delta(e)=\rho\beta_1e$; under limited attention it becomes
\begin{equation}
    \Delta_\psi(e)\;\equiv\;b\,\big[\psi\,e+(1-\psi)\,e^{\rm def}\big]\;=\;\omega_1\Big[\psi\,\Delta(e)+(1-\psi)\,\rho\beta_1\,e^{\rm def}\Big],
    \label{eq:Delta_psi}
\end{equation}
which equals $\Delta(e)$ when $\omega_1=\psi=1$. This one substitution carries the equilibrium analysis of Section~\ref{sec:equilibrium} over to limited attention: the hiding gain is a weighted average of the gain in Section~\ref{sec:equilibrium}, which depends on the pool, and a gain against the default, which does not.

\begin{proposition}[Disclosure Incentives under Limited Attention]
\label{prop:att_incentives}
Under Assumption~\ref{assum:attention}:
\begin{enumerate}[label=(\roman*)]
    \item an informed seller with $\theta=0$ strictly prefers non-disclosure if and only if $\tilde e>c_0/b$;
    \item an informed seller with $\theta=1$ strictly prefers non-disclosure if and only if $\tilde e>1+(c_0+c_1)/b$, which requires that silence sell above a verified efficient rating;
    \item if $C^{\rm def}<\bar C^{\rm sil}$, silence sells above a verified efficient rating ($\tilde e>1$, equivalently $C^*<C_1$) if and only if $\psi<(C_1-C^{\rm def})/(\bar C^{\rm sil}-C^{\rm def})$, which requires $C^{\rm def}<C_1$;
    \item if $\psi=1$, then $\tilde e=e$ for any $\omega_1$ and $C^{\rm def}$; conditions (i)--(ii) are then the net-of-cost comparisons underlying Proposition~\ref{prop:equilibrium}, and with $c_0=c_1=0$ they reduce to Proposition~\ref{prop:premium}.
\end{enumerate}
\end{proposition}

\noindent\textit{Proof.} See Appendix~\ref{app:attention}. $\Box$

Inattention that is the same for rated and silent listings is equivalent to having fewer EPC-sensitive buyers: it scales both premia together and leaves the model of Section~\ref{sec:equilibrium} intact (part iv). Differential attention matters only if the default is optimistic (part iii); a buyer who, without thinking about energy, implicitly assumes an average dwelling is a cursed buyer by another name and stays below the ceiling.\footnote{If the default equals the prior mean cost $\bar\theta C_1+(1-\bar\theta)C_0$, it is at least $C_1$, and silence never outprices a verified efficient rating. If it equals the rational cost of the silent pool, attention is irrelevant. The evidence on shrouded costs cited above finds under-weighting toward zero rather than toward the mean, which supports an optimistic default. Housing-specific evidence on this point does not yet exist, and it is the assumption on which the extension most depends.} And the high-quality seller loses her automatic preference for disclosure once silence outprices a verified efficient rating, so that only the hiding cost stops her from hiding (part ii). Applied to the seven letters instead of the two quality levels, Lemma~\ref{lem:att_prices} also predicts where silence sits among ratings: the gap between a rated listing and a silent one should be linear in the letter's energy cost and cross zero at $C^*$. The model of Section~\ref{sec:equilibrium} places that crossing between D and E; limited attention with an optimistic default places it at or above A.

\begin{corollary}[High-Quality Sellers Disclose outside Complete Adverse Selection]
\label{cor:att_hq}
In any equilibrium with $\sigma_0^*<1$, $\sigma_1^*=0$. In Complete Adverse Selection ($\sigma_0^*=1$), high-quality hiding is possible only if $\tilde e(1,\sigma_1)>1+(c_0+c_1)/b$ for some $\sigma_1$, which requires that silence sell above a verified efficient rating.
\end{corollary}

\noindent\textit{Proof.} See Appendix~\ref{app:attention}. $\Box$

Corollary~\ref{cor:hq_disclose} survives, but for a different reason. In Section~\ref{sec:equilibrium} high-quality sellers disclose because silence costs them money. Here silence may pay them too, but by~\eqref{eq:gains} it always pays them $b$ less than it pays low-quality sellers, so with a common hiding cost the low type is the first to want to hide. Whenever $\sigma_1^*=0$, the belief reduces to~\eqref{eq:belief_simplified} and the effective belief to $\tilde e(\sigma_0)=\psi\,e(\sigma_0)+(1-\psi)\,e^{\rm def}$.

\paragraph{Equilibrium characterization.}
By analogy with~\eqref{eq:Phi}, define the net hiding gain under limited attention,
\begin{equation}
    \Phi_\psi(\sigma_0)\;\equiv\;\Delta_\psi\big(e(\sigma_0)\big)-c_0\;=\;b\Big[\psi\,\frac{(1-\pi)\bar\theta}{\pi(1-\bar\theta)\sigma_0+(1-\pi)}+(1-\psi)\,e^{\rm def}\Big]-c_0,
    \label{eq:Phi_psi}
\end{equation}
and write $e(1)\equiv(1-\pi)\bar\theta/(1-\pi\bar\theta)$ for the belief when every informed low-quality seller hides.

\begin{proposition}[Equilibrium under Limited Attention]
\label{prop:att_equilibrium}
Under Assumption~\ref{assum:attention}, an equilibrium exists, and it is characterised as follows.
\begin{enumerate}[label=(\roman*)]
    \item If $c_0>\Delta_\psi(\bar\theta)$, the unique equilibrium is Cost-Deterred Hiding, with $\sigma_0^*=0$ and $e^*=\bar\theta$.
    \item If $c_0<\Delta_\psi(e(1))$, every informed low-quality seller hides, $\sigma_0^*=1$ (Complete Adverse Selection). Let $\bar g_1(\sigma_1)\equiv b\,[\tilde e(1,\sigma_1)-1]$ be the high type's gain. Then $\sigma_1^*=0$, with $e^*=e(1)$, is an equilibrium if and only if $\bar g_1(0)\le c_0+c_1$, and $\sigma_1^*=1$ is an equilibrium if and only if $\bar g_1(1)\ge c_0+c_1$. Both can hold, so equilibrium need not be unique in this regime. If $\psi=1$, only $\sigma_1^*=0$ survives.
    \item If $\Delta_\psi(e(1))\le c_0\le\Delta_\psi(\bar\theta)$, the equilibrium is Partial Pooling. If $\psi>0$ it is unique, with
    \begin{equation}
        e^*\;=\;\frac{c_0/b-(1-\psi)\,e^{\rm def}}{\psi},\qquad \sigma_0^*\;=\;\frac{1-\pi}{\pi(1-\bar\theta)}\Big(\frac{\bar\theta}{e^*}-1\Big).
        \label{eq:att_pp}
    \end{equation}
    If $\psi=0$, Partial Pooling requires $c_0=b\,e^{\rm def}$ exactly, and then every $\sigma_0\in[0,1]$ is an equilibrium.
\end{enumerate}
\end{proposition}

\noindent\textit{Proof.} See Appendix~\ref{app:attention}. $\Box$

\begin{table}[htbp]
\centering
\caption{Equilibrium Regimes under Limited Attention}
\label{tab:att_equilibria}
\small\renewcommand{\arraystretch}{1.5}
\begin{tabular*}{\textwidth}{@{\extracolsep{\fill}}lccc@{}}
\toprule\toprule
\textbf{Equilibrium} & \textbf{Condition} & \textbf{Strategy $\sigma_0^*$} & \textbf{Non-disclosure premium $\Delta^*$} \\
\midrule
Cost-Deterred Hiding & $c_0>\Delta_\psi(\bar\theta)$ & $0$ & $b\,[\psi\bar\theta+(1-\psi)e^{\rm def}]$ \\
Partial Pooling & $\Delta_\psi(e(1))\le c_0\le\Delta_\psi(\bar\theta)$ & \eqref{eq:att_pp} & $c_0$ \\
Complete Adverse Selection & $c_0<\Delta_\psi(e(1))$ & $1$ & $b\,[\psi e(1)+(1-\psi)e^{\rm def}]$ \\
\bottomrule\bottomrule
\end{tabular*}

\vspace{0.5em}
\raggedright
\footnotesize
\textit{Notes:} The premium in Complete Adverse Selection is evaluated at $\sigma_1^*=0$. At $\psi=1$ and $\omega_1=1$ the table reduces to Table~\ref{tab:equilibria}.
\end{table}

% [AUTHOR REWRITE] lead-in to the proposition on seller types
The three regimes keep their logic, and the central result of Section~\ref{sec:equilibrium} survives intact: in Partial Pooling the non-disclosure premium equals the hiding cost, because the marginal hider is indifferent. That statement concerns the marginal seller, not buyers, so it holds whatever buyers attend to. What limited attention changes is where silence sits among the ratings, and with it which seller gains most from silence. The next proposition derives the price of each type of listing and the payoff of each type of seller. Uninformed sellers cannot hide a rating they do not have, so they bear no hiding cost.

\begin{proposition}[Prices and Payoffs by Seller Type]
\label{prop:att_seller_types}
Under Assumption~\ref{assum:attention}, with $b=\rho\omega_1\beta_1>0$, measure prices and payoffs relative to a listing displaying an inefficient rating, and consider an equilibrium with $\sigma_1^*=0$. A listing displaying an efficient rating sells at $b$, and a silent listing at
\begin{equation}
    \Delta^*\;=\;b\,\tilde e^*\;=\;\min\Big\{\max\big\{c_0,\,\Delta_\psi(e(1))\big\},\;\Delta_\psi(\bar\theta)\Big\},
    \label{eq:att_premium_trunc}
\end{equation}
the hiding cost truncated to the Partial Pooling band. An informed high-quality seller earns $b$, an informed low-quality seller earns $\max\{\Delta^*-c_0,\,0\}$, and an uninformed seller of either quality earns $\Delta^*$ (Table~\ref{tab:att_seller_types}). By regime:
\begin{itemize}
    \item in \textbf{Cost-Deterred Hiding}, $\Delta^*=\Delta_\psi(\bar\theta)<c_0$, and every informed low-quality seller displays her rating and earns zero;
    \item in \textbf{Partial Pooling}, $\Delta^*=c_0$ and $\tilde e^*=c_0/b$. Every informed low-quality seller earns zero, whether she hides or displays, and the uninformed seller earns $c_0$. A silent listing sells above a verified efficient rating if and only if $c_0>b$. Payoffs are then ordered uninformed, informed high quality, informed low quality; if $0<c_0<b$ they are ordered informed high quality, uninformed, informed low quality;
    \item in \textbf{Complete Adverse Selection}, $\Delta^*=\Delta_\psi(e(1))>c_0$, and every informed low-quality seller hides and earns $\Delta^*-c_0>0$, which is $c_0$ less than the uninformed seller.
\end{itemize}
In every regime, a silent listing sells above a verified efficient rating, $\Delta^*>b$, if and only if $e^{\rm def}>1$ and
\begin{equation}
    \psi\;<\;\bar\psi(e^*)\;\equiv\;\frac{e^{\rm def}-1}{e^{\rm def}-e^*},
    \label{eq:att_psibar}
\end{equation}
the threshold of Proposition~\ref{prop:att_incentives}(iii) evaluated at the equilibrium belief. The uninformed seller then earns more than an informed high-quality seller and at least as much as any informed seller. In an equilibrium with $\sigma_1^*>0$, which exists only in Complete Adverse Selection, $\Delta^*\ge b+c_0+c_1$.
\end{proposition}

\noindent\textit{Proof.} See Appendix~\ref{app:attention}. $\Box$

\begin{table}[htbp]
\centering
\caption{Prices and Payoffs by Listing and Seller Type under Limited Attention}
\label{tab:att_seller_types}
\small\renewcommand{\arraystretch}{1.5}
\begin{tabular*}{\textwidth}{@{\extracolsep{\fill}}lccc@{}}
\toprule\toprule
 & \textbf{Cost-Deterred Hiding} & \textbf{Partial Pooling} & \textbf{Complete Adverse Selection} \\
\midrule
\multicolumn{4}{l}{\textit{A. Price, relative to a displayed inefficient rating}} \\
Displayed inefficient rating (E to G) & $0$ & $0$ & $0$ \\
Displayed efficient rating (A to D) & $b$ & $b$ & $b$ \\
Silent listing, $\Delta^*$ & $\Delta_\psi(\bar\theta)$ & $c_0$ & $\Delta_\psi(e(1))$ \\
\midrule
\multicolumn{4}{l}{\textit{B. Net payoff, relative to a displayed inefficient rating}} \\
Informed, high quality & $b$ & $b$ & $b$ \\
Informed, low quality, displays & $0$ & $0$ & -- \\
Informed, low quality, hides & -- & $0$ & $\Delta_\psi(e(1))-c_0$ \\
Uninformed, either quality & $\Delta_\psi(\bar\theta)$ & $c_0$ & $\Delta_\psi(e(1))$ \\
\midrule
Silence above an efficient rating iff & $\psi<\bar\psi(\bar\theta)$ & $c_0>b$ & $\psi<\bar\psi(e(1))$ \\
\bottomrule\bottomrule
\end{tabular*}

\vspace{0.5em}
\raggedright
\footnotesize
\textit{Notes:} Equilibria with $\sigma_1^*=0$. A dash marks a seller type that is not on the equilibrium path; in Cost-Deterred Hiding a low-quality seller who deviated to silence would earn $\Delta_\psi(\bar\theta)-c_0<0$. Displayed prices do not depend on beliefs (Lemma~\ref{lem:att_prices}) and are defined in every regime. $\bar\psi(e)$ is defined in~\eqref{eq:att_psibar} for $e^{\rm def}>1$; if $e^{\rm def}\le1$ silence never sells above an efficient rating. Uninformed sellers bear no hiding cost; relative to her own verified price, an uninformed high-quality seller earns $\Delta^*-b$. In Complete Adverse Selection with $\sigma_1^*>0$ (Proposition~\ref{prop:att_equilibrium}), $\Delta^*\ge b+c_0+c_1$ and an informed high-quality seller hides and earns $\Delta^*-c_0-c_1\ge b$. At $\psi=\omega_1=1$ the silent price is $\rho\beta_1e^*<b$ in every column (Table~\ref{tab:equilibria}).
\end{table}

\begin{corollary}[Where Silence Sits among Ratings]
\label{cor:att_ordering}
Suppose $c_0+c_1>0$. If $e^{\rm def}>1$ and $\psi<\bar\psi(e(1))$, a silent listing sells above a verified efficient rating in every equilibrium, whatever the hiding cost. If $e^{\rm def}\le1$ or $\psi\ge\bar\psi(\bar\theta)$, it sells at most at the price of a verified efficient rating in every equilibrium, whatever the hiding cost. Otherwise, in any equilibrium with $\sigma_1^*=0$, it sells above a verified efficient rating if and only if $c_0>b$. In the model of Section~\ref{sec:equilibrium}, $\psi=1$ and $0<\Delta^*<b$ in every equilibrium.
\end{corollary}

\noindent\textit{Proof.} See Appendix~\ref{app:attention}. $\Box$

% [AUTHOR REWRITE] empirical reading of the proposition
The proposition gives the two segments one reading. In apartments the non-disclosure premium, \attDthreeApt{} log points, is about twice the energy-efficiency premium, \attBthreeApt{}: the effective belief is \attRatioApt{}, and silent listings sell above verified efficient ones, which sell above inefficient ones. In houses the effective belief is \attRatioHouse{} and silence sells between the two ratings. These orderings are read directly from the two estimated premia. Which case of Corollary~\ref{cor:att_ordering} each segment falls in also depends on the normalisation of the default ($C^{\rm def}=0$): at the estimated attention apartments are in the first case ($\hat\psi_\delta=\attPsiApt{}$ against a threshold of \attPsiBarOneApt{}), so silence outprices a verified efficient rating whatever the hiding cost, and houses are in the second ($\attPsiHouse{}$ against \attPsiBarThetaHouse{}). The model with one $\psi$ per segment therefore predicts the first ordering in every apartment market and the second in every house market. The market-level estimates are noisy and $\psi$ is common within a segment, so the counts differ: silence sells above a verified efficient rating in \attNDgtBApt{} of \attNMarketsApt{} apartment markets and \attNDgtBHouse{} of \attNMarketsHouse{} house markets, between the two displayed ratings in \attNBetweenApt{} and \attNBetweenHouse{}, and below a displayed inefficient rating in \attRawNNegApt{} and \attRawNNegHouse{}. The last cannot arise in Partial Pooling and reflects sampling error or heterogeneity that we do not model. The payoff ranking among informed sellers is stated at the hiding cost that reproduces the observed hiding rate, which at the aggregate and in Partial Pooling is the premium itself, $c_0=\hat\Delta$ (Proposition~\ref{prop:att_identification}(iv)); the price ordering needs only $\hat\Delta$ against $\hat\beta_1$. Read through Partial Pooling, an uninformed apartment seller earns \attDthreeApt{} log points over a displayed inefficient rating, an informed high-quality seller \attBthreeApt{}, and an informed low-quality seller nothing, whether she hides or not. Silence paying more than a verified rating does not make high-quality sellers hide, which is an implication of Partial Pooling and not a separate finding: their gain, \attExcessApt{} log points, falls short of the hiding cost \attDthreeApt{} that keeps low-quality sellers indifferent. The uninformed seller's rent is a property of the maintained assumptions, that uninformed sellers have no strategy, bear no hiding cost and are exogenous in number, and not a finding.

Limited attention also changes the location and the width of the band of hiding costs that sustains Partial Pooling. Both bounds shift up by $(1-\psi)\,b\,e^{\rm def}$, the part of the gain that does not depend on the pool, and the width shrinks by the factor $\psi$.

\begin{corollary}[The Partial Pooling Band]
\label{cor:att_band}
The interval of hiding costs that supports Partial Pooling is $[\Delta_\psi(e(1)),\,\Delta_\psi(\bar\theta)]$, of width $\psi\,b\,[\bar\theta-e(1)]$. Its width goes to zero as $\psi\to0$.
\end{corollary}

\noindent\textit{Proof.} Direct from~\eqref{eq:Delta_psi} and Proposition~\ref{prop:att_equilibrium}, as in Corollary~\ref{cor:rho}(ii). $\Box$

The corollary has a sharp empirical counterpart. At the apartment estimates of Section~\ref{ss:att_estimates}, taking the reform-period estimate $\hat\psi_\delta$ as a proxy for $\psi$ (exact under the maintained assumption of Proposition~\ref{prop:att_identification}), the band spans \attBandLoApt{} to \attBandHiApt{} log points, a width of \attBandWidthApt{}, against \attBandWidthBaseApt{} in the model of Section~\ref{sec:equilibrium}. The mean premium lies inside the aggregate band by construction of the aggregation, which is not evidence. Market by market, with a common $\psi$, the premium falls inside the market's own band in \attShareInBandApt{} percent of apartment markets: a single hiding cost per market would have to fall inside a window of about \attBandWidthApt{} log points for an apartment market to show interior hiding, yet the registry shows interior hiding rates in essentially every market. In houses the band is wider, \attBandWidthHouse{} log points, but only \attShareInBandHouse{} percent of house markets fall inside their own band. A single hiding cost per market therefore generates interior hiding only within a narrow window of costs, particularly in apartments; we do not model heterogeneity in hiding costs across sellers, and we read the band as a limitation of the single-cost specification.

\paragraph{Comparative statics.}
The next proposition collects how the equilibrium responds to attention, to the default, and to the parameters of Section~\ref{sec:equilibrium}.

\begin{proposition}[Comparative Statics under Limited Attention]
\label{prop:att_cs}
Suppose $C^{\rm def}<\bar C^{\rm sil}$.
\begin{enumerate}[label=(\roman*)]
    \item In Partial Pooling, $\partial\sigma_0^*/\partial\psi<0$, $\partial\sigma_0^*/\partial C^{\rm def}<0$, $\partial\sigma_0^*/\partial c_0<0$, $\partial\sigma_0^*/\partial(\rho\omega_1)>0$ and $\partial\sigma_0^*/\partial\pi<0$. The effect of $c_0$ on the equilibrium belief, $\partial e^*/\partial c_0=1/(\psi b)$, grows without bound as $\psi\to0$.
    \item In the corner regimes, the premium is decreasing in $\psi$ whenever $C^{\rm def}<\bar\theta C_1+(1-\bar\theta)C_0$ (which implies the corresponding condition at $e(1)$). In Complete Adverse Selection, $\partial\Delta^*/\partial\pi=\psi\,b\,\partial e(1)/\partial\pi<0$: the belief channel of Proposition~\ref{prop:comparative_statics}(i), scaled by $\psi$.
\end{enumerate}
\end{proposition}

\noindent\textit{Proof.} See Appendix~\ref{app:attention}. $\Box$

The comparative statics separate three channels. The \emph{revenue channel} ($\rho\omega_1$) works as in Section~\ref{sec:equilibrium}: a larger effective attention to displayed ratings raises the gain from concealment and the hiding rate. The \emph{belief channel} ($\pi$) keeps its sign but is scaled by $\psi$, because a higher share of informed sellers makes silence more informative only to buyers who think about what silence means. The \emph{attention channel} ($\psi$, $C^{\rm def}$) is new: more attention to silence, or a less optimistic default, lowers the gain from hiding and the hiding rate. These are the margins on which a salience policy acts (Section~\ref{ss:att_policy}). The last statement in part~(i) is the formal reason the single-cost model becomes fragile in apartments: when silence is barely read, a small change in the hiding cost produces a large swing in the hiding rate, because nothing on the buyer side pushes back.

\paragraph{High-quality hiding and non-strategic silence.}
Proposition~\ref{prop:att_incentives}(ii) and Corollary~\ref{cor:att_hq} say when high-quality sellers may hide: only if silence sells above a verified efficient rating. The registry shows similar hiding rates for the two grades, \attSigmaOneApt{} against \attSigmaZeroApt{} percent in apartments and \attSigmaOneHouse{} against \attSigmaZeroHouse{} percent in houses. In apartments silence pays a high-quality seller \attExcessApt{} log points more than her rating, yet she displays, because in Partial Pooling her slack is $b+c_1$ (Proposition~\ref{prop:att_seller_types}). The pattern in the registry is not evidence for strategic hiding by high-quality sellers: house sellers of high quality stay silent although silence costs them about \attHighLossHouse{} log points relative to their rating, so some silence is not strategic at all, for instance listings whose certificate is still pending. The same non-strategic silence is in the hiding rate of low-quality sellers, so the interior hiding rates we use may overstate strategic hiding.

\paragraph{Welfare.}
The welfare accounting of Section~\ref{sec:equilibrium} acquires a third component. As before, the hiding cost $c_0$ is a deadweight loss to the extent it destroys real resources and a transfer to the extent it reflects expected sanctions, informed hiders net zero in Partial Pooling and $\Delta^*-c_0>0$ in Complete Adverse Selection (Proposition~\ref{prop:att_seller_types}), and uninformed sellers earn the non-disclosure premium as a rent. The premium itself now contains a behavioural transfer. Relative to what a fully attentive market would pay, a silent listing sells at a mark-up of
\begin{equation}
    p_{\rm sil}-\big(\bar V-(1-\rho)\,C^{\rm ins}-\rho\,\bar C^{\rm sil}\big)\;=\;\rho\,(1-\omega_0)\,\big(\bar C^{\rm sil}-C^{\rm def}\big),
    \label{eq:att_markup}
\end{equation}
paid by inattentive EPC-sensitive buyers to every silent seller, uninformed ones included. Unlike the allocative distortion of Section~\ref{sec:equilibrium}, which arises because concealed information is not transmitted, this transfer arises because transmitted absence is not read. The model identifies its level only under additional assumptions (Section~\ref{ss:att_identification}); with displayed ratings fully attended and no pre-reform energy pricing it equals $\hat\lambda$, about \attLamApt{} log points per silent apartment sale, and if displayed ratings are not fully attended $\hat\lambda$ is a lower bound.

\subsubsection{The Reform and Identification}
\label{ss:att_identification}

The estimands of Section~\ref{sec:identification} are reform-induced changes, so attention enters through how the reform moved it. For any object $x$, write $\delta x\equiv x^{\rm post}-x^{\rm pre}$ for its change across the July 2021 reform. In this notation $\hat\beta_{1,m}$ estimates $\delta b_m$ and $\hat\Delta_m$ estimates $\delta\Delta_m$. Let the reform raise effective attention to rated listings by $\delta(\rho\omega_1)$ and to silent listings by $\delta(\rho\omega_0)$, and define the \emph{spillover} $\psi_\delta\equiv\delta(\rho\omega_0)/\delta(\rho\omega_1)$, the share of the reform's new attention to energy that reached listings showing no rating. Under the maintained assumption of Section~\ref{sec:identification} that buyers did not price energy before July 2021 (Assumption~\ref{assum:no_pre_pricing}), $\psi_\delta$ equals post-reform relative attention. Normalising the default to zero, equations~\eqref{eq:b_att}--\eqref{eq:premium_att} in differences give
\begin{equation}
    \hat\lambda_m\;\equiv\;\hat\Delta_m-\hat\beta_{1,m}\,\hat e^*_m\;=\;(1-\psi_\delta)\,\hat\beta_{1,m}\,\big(\eta_m-\hat e^*_m\big).
    \label{eq:att_lambda}
\end{equation}

\begin{proposition}[Identification under Limited Attention]
\label{prop:att_identification}
Under Assumptions~\ref{assum:attention} and~\ref{assum:no_pre_pricing} and the maintained assumptions of Section~\ref{sec:identification}:
\begin{enumerate}[label=(\roman*)]
    \item $\hat\rho_m=\hat\beta_{1,m}/\hat\beta^{\rm anchor}_{1,m}$ identifies $\delta(\rho\omega_1)$, the reform-induced increase in effective attention to displayed ratings; with $\omega_1=1$ it is the $\rho_m$ of Section~\ref{sec:identification};
    \item $\hat\Delta_m/\hat\beta_{1,m}$ identifies the perceived cost of silence $C^*_\delta\equiv C^{\rm def}+\psi_\delta(\bar C^{\rm sil}-C^{\rm def})$, in units of certified consumption;
    \item given $C^{\rm def}=0$,~\eqref{eq:att_lambda} point-identifies $\psi_\delta$; without that normalisation, $\psi_\delta\in[0,\,C^*_\delta/\bar C^{\rm sil}]$ and $C^{\rm def}\in[0,C^*_\delta]$;
    \item in Partial Pooling, the hiding cost of the informed low-quality seller equals the post-reform non-disclosure premium, $\delta\Delta$, which $\hat\Delta_m$ estimates.
\end{enumerate}
\end{proposition}

\noindent\textit{Proof.} See Appendix~\ref{app:attention}. $\Box$

The over-identifying wedge is the part of the reform's price effect that silence escaped. Had silent listings received the same increase in attention as rated ones, their price relative to E-to-G listings would have moved by $\hat\beta_1\hat e^*$, as Section~\ref{sec:equilibrium} predicts; they received less, and the shortfall, priced at the cost of the silent pool, is $\hat\lambda$. A spillover of one means the reform made energy something buyers ask about whether or not a rating is shown; a spillover near zero means it made the rating something buyers read but not something they ask for. Part~(iii) states what the design cannot do. The two gaps pin down one number per market, the perceived cost of silence, and low attention with a modest default fits them as well as slightly higher attention with an even more optimistic one. Normalising the default to zero delivers the largest attention compatible with the data.\footnote{The same algebra in levels identifies relative attention at a point in time rather than the spillover. We do not estimate it, because across apartment markets the level gaps carry a large market-specific component that is unrelated to energy and that the differences-in-differences remove: the ratio of mean level gaps and the cross-market slope of one level gap on the other imply very different values, $\attPsiLevROMApt$ and $\attPsiLevSlopeApt$ post-reform. We use levels only for signs.}

Three testable implications follow. By~\eqref{eq:etilde}, $\Delta/b\le\eta$ whenever $C^{\rm def}\ge0$, and since $\eta_m$ comes from the registry alone this is a market-by-market bound on the two estimated gaps. Equation~\eqref{eq:att_lambda} has no constant: a market where the reform did not move the energy-efficiency premium should show no wedge. And the registry belief $\hat e^*_m$ enters~\eqref{eq:att_lambda} with a weight proportional to $\psi_\delta$ (Lemma~\ref{lem:etilde_monotone}), so where attention to silence is low, mismeasurement of the silent pool's composition has little bearing on the estimate.\footnote{At the apartment estimates an error of 0.1 in $\hat e^*$ moves $\hat\psi_\delta$ by about $0.1\times\attSensEApt$. In house markets the sensitivity is about \attSensRatio{} times larger, as in Section~\ref{sec:equilibrium}, and the measurement of the silent pool matters correspondingly more there. For the same reason, the restriction $\sigma_1=0$ imposed in $\hat e^*$ (Section~\ref{sec:beliefs}) matters little for the apartment estimate.}

\subsubsection{Estimates}
\label{ss:att_estimates}

Table~\ref{tab:attention} compares what the core model and the limited-attention model predict for the non-disclosure premium, and reports the attention estimates. We estimate $1-\psi_\delta$ as the ratio of the cross-market sums of the two sides of~\eqref{eq:att_lambda}. Absent selection, and unlike a regression through the origin, this estimator is free of the attenuation and correlated-noise biases that arise when $\hat\beta_{1,m}$ is noisy and shares the E-to-G index with $\hat\Delta_m$. The estimation sample is, however, selected on $\hat\beta_{1,m}>0$, and selecting on a noisy $\hat\beta_{1,m}$ biases both estimators toward full attention: a Monte Carlo calibrated to our markets puts the ratio of sums between \attMCromSelLo{} and \attMCromSelHi{} for a true value of \attMCtrue{} (Appendix~\ref{app:att_estimator}). The apartment estimate is therefore, if anything, an upper bound on attention to silence, and the house estimate may overstate it. We estimate one $\psi_\delta$ per segment rather than per market. A market-specific value is identified by the market's own ratio $\hat\Delta_m/\hat\beta_{1,m}$, but it would fit each premium exactly, place every market in Partial Pooling by construction, and leave no restriction to test, and the market ratios are too noisy to be informative on their own.\footnote{The estimators rank by how much weight they place on markets with a large energy-efficiency premium: the median of market ratios gives $\attPsiMedApt$, the ratio of sums $\attPsiApt$, the regression through the origin $\attPsiOLSApt$, and a regression with a constant, which the model excludes, $\attPsiConstApt$ (intercept $\attIcptApt$, s.e.\ $\attIcptSeApt$).}

\begin{table}[htbp]
\centering
\caption{The Non-Disclosure Premium under the Core Model and under Limited Attention}
\label{tab:attention}
\small\renewcommand{\arraystretch}{1.3}
\begin{tabular*}{\textwidth}{@{\extracolsep{\fill}}lcc@{}}
\toprule\toprule
 & \textbf{Apartments} ($N=45$) & \textbf{Houses} ($N=97$) \\
\midrule
\multicolumn{3}{l}{\textit{A. What the data show (log points)}} \\
Energy-efficiency premium $\hat\beta_1$ & 4.1 & 5.8 \\
Non-disclosure premium $\hat\Delta$ & 8.5 & 3.9 \\
Ratio $\hat\Delta/\hat\beta_1$ (at most one in the core model) & 2.10 & 0.67 \\
Markets above the core-model ceiling $\hat\Delta=\hat\beta_1$ & 37 of 45 & 33 of 97 \\
\midrule
\multicolumn{3}{l}{\textit{B. Predicted non-disclosure premium (log points)}} \\
Core model (full attention): $\hat\beta_1\hat e^*$ & 2.4 & 3.5 \\
\quad Observed minus predicted ($t$) & 6.1 (6.8) & 0.4 (1.1) \\
\cmidrule[1pt]{1-3}
$\triangleright$ \textbf{Limited attention: $\hat\beta_1[\hat\psi\hat e^*+(1-\hat\psi)\eta]$} & \textbf{8.5} & \textbf{3.9} \\
 \quad \textbf{Observed minus predicted} & \textbf{0.0} & \textbf{0.0} \\
\cmidrule[1pt]{1-3}
\multicolumn{3}{l}{\textit{C. Attention to silence}} \\
$\hat\psi_\delta$, ratio of sums [95\% CI] & 0.13 [\ensuremath{-}0.12, 0.37] & 0.95 [0.86, 1.04] \\
$\hat\psi_\delta$, regression through the origin (s.e.) & 0.21 (0.11) & 0.96 (0.04) \\
$t$-statistic, full attention ($\psi_\delta=1$) & 6.9 & 1.1 \\
Identified interval for $\psi_\delta$ ($C^{\rm def}\ge0$) & $[0,\,0.13]$ & $[0,\,0.95]$ \\
Perceived cost of silence $C^*_\delta$ (kWh$_{\rm ep}$/m$^2$/yr) [95\% CI] & 31 [\ensuremath{-}31, 89] & 233 [210, 256] \\
Rational cost of the silent pool $\bar C^{\rm sil}$ (kWh$_{\rm ep}$/m$^2$/yr) & 244 & 245 \\
Energy-cost ratio $\eta$ (weighted) & 2.32 & 2.11 \\
Mean slack in the attention ceiling $\eta\hat\beta_1-\hat\Delta$ ($t$) & 0.90 (1.0) & 8.37 (13.0) \\
Markets above their own attention ceiling & 21 of 45 & 9 of 97 \\
\bottomrule\bottomrule
\end{tabular*}


\vspace{0.5em}
\raggedright
\footnotesize
\textit{Notes:} Difference-in-differences estimates on the estimation sample of Table~\ref{tab:did_results}, log points unless stated; means across markets. Panel~B is the comparison the table is built around. The core model predicts a non-disclosure premium of $\hat\beta_1\hat e^*$, which falls short of the observed premium by \attLamApt{} log points in apartments and by \attLamHouse{} in houses. With limited attention the predicted premium is $\hat\beta_1[\hat\psi\hat e^*+(1-\hat\psi)\eta]$, which matches the observed premium in both segments. In apartments the match uses one parameter, $\hat\psi$, estimated from the ratio of the same two aggregates, so it shows that a single attention parameter is enough and is not an out-of-sample test. The evidence for limited attention is in panel~C: attention to silence is a small fraction of attention to ratings in apartments (\attPsiApt{}) and close to full in houses (\attPsiHouse{}), and full attention is rejected for apartments and not for houses. Aggregates weight markets by $\hat\beta_{1,m}$, the weights implied by the ratio-of-sums estimator. $\eta_m$ and $\bar C^{\rm sil}$ value registry letters at DPE band midpoints; values are provisional pending certified consumption. Confidence intervals from a bootstrap over markets (\attNBoot{} draws); percentile intervals are not truncated to the admissible range $[0,1]$ of $\psi_\delta$. The identified interval allows any default $C^{\rm def}\in[0,C^*_\delta]$. The attention ceiling is $\hat\Delta=\eta_m\hat\beta_{1,m}$ (Section~\ref{ss:att_estimates}).
\end{table}


Figure~\ref{fig:att_contrast} plots the quantity that does not depend on the normalisation of the default, the perceived cost of silence, together with the estimates of $\psi_\delta$.

\begin{figure}[htbp]
    \centering
    \caption{Perceived cost of silence and relative attention, apartments and houses}
    \label{fig:att_contrast}
    \includegraphics[width=\linewidth]{fig/visuals/fig_att_F3_contrast.pdf}
    \begin{minipage}{0.95\linewidth}
        \vspace{0.5em}
        \footnotesize\textit{Notes.} Left: perceived cost of silence $C^*_\delta$ in kWh$_{\rm ep}$/m$^2$/yr with 95\% bootstrap intervals (apartments \attCstarApt{} [\attCstarLoApt, \attCstarHiApt]; houses \attCstarHouse{} [\attCstarLoHouse, \attCstarHiHouse]), against the rational cost of the silent pool (\attCsilApt{} and \attCsilHouse{}). Right: estimates of relative attention $\psi_\delta$ by estimator, and the ratio of sums with the G-band energy midpoint set to 420 or 600 kWh$_{\rm ep}$/m$^2$/yr. The shaded band is the admissible range $[0,1]$. Intervals are not truncated. Apartment intervals are wide, and selection of markets on $\hat\beta_{1,m}>0$ biases the estimates toward full attention.
    \end{minipage}
\end{figure}

% [AUTHOR REWRITE] interpretation of the estimates
Full attention is rejected in apartments and not in houses. The evidence is the contrast between segments, which does not depend on the normalisation of the default: the perceived cost of silence, identified without it, is \attCstarApt{} kWh per square metre and year in apartments (95\% CI \attCstarLoApt{} to \attCstarHiApt{}) against \attCstarHouse{} (\attCstarLoHouse{} to \attCstarHiHouse{}) in houses ($t=\attDiffCstarT$ for the difference). In apartments it is a small fraction of the \attCsilApt{} kWh a fully attentive buyer would assign to the silent pool, so silent apartments are priced roughly as though energy were not part of the transaction. In houses it is close to the rational cost of the pool, which is why the model of Section~\ref{sec:equilibrium} fits that segment. One attention parameter per segment thus accounts for both the rejection in apartments and the fit in houses, without any change to the seller side; its level, unlike the contrast, depends on the normalisation of the default. By Proposition~\ref{prop:att_identification}(iii), the apartment estimate is best read as an upper bound. House prices are consistent with silence being priced at the rational cost of the pool, which is all Section~\ref{sec:equilibrium} requires; they do not show that house buyers are attentive in any absolute sense.

% [AUTHOR REWRITE] fit claim
The extension also changes how well the model predicts the apartment premium, though less than the aggregate suggests (Figure~\ref{fig:att_fit}). It removes the level bias in the predicted premia, the mean residual falling from \attResidCoreApt{} to \attResidExtApt{} log points, but that zero holds by construction, because $\hat\psi_\delta$ is estimated from the same ratio. The root mean squared error falls from \attRmseCoreApt{} to \attRmseExtApt{}, yet net of the level the residual dispersion is unchanged (\attSdResidCoreApt{} against \attSdResidExtApt{}) and the correlation between predicted and observed premia moves only from \attCorrCoreApt{} to \attCorrExtApt{}. The extension therefore explains why silent apartments sell above verified efficient ones, not which markets carry the largest premia. In houses the two models coincide. The evidence for the extension rests on the ceilings of Figures~\ref{fig:att_belief} and~\ref{fig:att_ceilings} (the latter in Appendix~\ref{app:attention}) and the contrast of Figure~\ref{fig:att_contrast}, not on the aggregate fit.

\begin{figure}[htbp]
    \centering
    \caption{Observed and predicted non-disclosure premium, core model and extension}
    \label{fig:att_fit}
    \includegraphics[width=\linewidth]{fig/visuals/fig_att_F1_fit.pdf}
    \begin{minipage}{0.95\linewidth}
        \vspace{0.5em}
        \footnotesize\textit{Notes.} Each market appears twice at the same observed premium $\hat\Delta_m$ (vertical axis, log points): at the core-model prediction $\hat\beta_{1,m}\hat e^*_m$ (hollow) and at the limited-attention prediction $\hat\beta_{1,m}[\hat\psi_\delta\hat e^*_m+(1-\hat\psi_\delta)\eta_m]$ (filled), joined by a thin line. The dashed line marks observed equal to predicted. The box reports the mean residual, the root mean squared error and the correlation of each prediction with $\hat\Delta_m$. The zero mean residual of the extension holds by construction, because $\hat\psi_\delta$ is estimated from the same aggregate ratio. The apartment sample has \attNMarketsApt{} markets and the house sample \attNMarketsHouse{}.
    \end{minipage}
\end{figure}

The attention reading also implies a different concealment cost for apartments: the hiding cost of the marginal hider is the non-disclosure premium itself, \attDthreeApt{} log points, about \attCzeroRatio{} times the \attCzeroBaseApt{} implied by Section~\ref{sec:identification}. A cost of that size is well above the transaction-cost benchmarks of Section~\ref{sec:results}, one reason we keep the baseline estimate as the paper's measure of concealment costs and report this as the alternative reading the extension implies; in houses the two coincide closely (\attCzeroBaseHouse{} and \attDthreeHouse{}).

The model does not take a stance on why apartment buyers fail to attend to energy when no rating is shown. It may be that energy is bundled into service charges. It may be that apartment buyers, often moving from rented housing in which heating was included, have no habit of pricing it. Or it may be that decisions in fast urban markets leave little room to ask. Irrespective of which applies, the implication for prices is the same, and so is the estimate.

Three qualifications bound these results. First, apartment markets cannot be distinguished from the bound of~\eqref{eq:etilde} on average (mean slack \attSlackApt{} log points, $t=\attSlackTApt$): the extension fits them only with a default close to ``no energy cost''. Second, \attNViolApt{} of the \attNMarketsApt{} apartment markets exceed the bound individually; violations are more frequent where $\hat\beta_{1,m}$ is small and the ratio is noisy (\attShareViolBottomApt{} percent of markets in the bottom tercile of $\hat\beta_{1,m}$), but \attShareViolTopApt{} percent of markets in the top tercile also exceed it, so limited attention is the leading force on average but probably not the only one in every market. Third, the level of $\hat\psi_\delta$ depends on the estimator, while the contrast between segments does not.\footnote{The bound depends on how certified consumption is valued. Our $\eta_m$ values registry letters at band midpoints and the open-ended G band at 500 kWh. Valuing G at 420 or 600 kWh moves the apartment estimate of $\psi_\delta$ to $\attPsiGfourtwentyApt$ or $\attPsiGsixhundredApt$, and the number of apartment markets violating the bound to \attNViolGfourtwentyApt{} or \attNViolGsixhundredApt. % [AUTHOR] replace with certified consumption (gate G3) and delete this footnote.
}

Splitting the post-reform window by year, $1-\hat\psi_\delta$ equals \attOneMinusPsiYoneApt, \attOneMinusPsiYtwoApt{} and \attOneMinusPsiYthreeApt{} in apartments and \attOneMinusPsiYoneHouse, \attOneMinusPsiYtwoHouse{} and \attOneMinusPsiYthreeHouse{} in houses. These are imprecise point estimates, computed with the pooled 2022--2024 registry moments; we do not interpret their time pattern without year-specific moments. % [AUTHOR] revisit after G2/G3.

Section~\ref{sec:policy} builds the policy extension on this model.


% ============================================================
\section{Policy Extension: Renovation Subsidies}
\label{sec:policy}
% ============================================================

\subsection{What the Model Adds}
\label{ss:policy_scope}

% [AUTHOR REWRITE] section opening
The core model treats $\theta$ as exogenous, but the policy question is what happens once quality becomes a choice. An informed seller who holds a low-quality certificate can hide it, disclose it, or renovate and disclose a verified high rating, and a subsidy is the natural instrument when renovation costs more than the market pays for it. This section works through that choice in the model of Section~\ref{ss:attention}, in which a buyer's attention to energy depends on whether the listing displays a rating. The model of Sections~\ref{ss:model_setup}--\ref{sec:equilibrium} is the case $\psi=\omega_1=1$, so every statement below reduces to its full-attention counterpart at that point. We use the attention model throughout for reasons of fit: it reproduces the apartment premium that the full-attention model cannot reach, and in houses, where $\hat\psi$ is \attPsiHouse{}, the two models nearly coincide (Section~\ref{ss:att_estimates}). The hiding cost is a single $c_0$ per market, held at the value that reproduces the observed hiding rate at the estimated attention, which is \attCzeroFitApt{} log points in apartments and \attCzeroFitHouse{} in houses (\attCzeroCoreMeanApt{} and \attCzeroCoreMeanHouse{} in the full-attention model). Uninformed sellers have no strategy set, and the thresholds are computed holding beliefs fixed, except in Section~\ref{ss:policy_ge}, which lets the pool respond.

Two caveats frame what the model adds. First, wherever the renovation threshold is the price gap $b$ between a verified efficient and a verified inefficient rating (Proposition~\ref{prop:threshold}), that gap is the energy-efficiency premium $\hat\beta_1$ that Section~\ref{sec:identification} estimates directly, so the model does not produce a new number there. What it adds is the identification of which price gap governs renovation (not the non-disclosure premium $\hat\Delta$, and not $\hat\Delta+c_0$, which is what the definition of the threshold suggests on first reading), the conditions under which that reading fails, and the counterfactuals of Sections~\ref{ss:policy_ge} and~\ref{ss:att_policy}, which need a model to say what happens to hiding and to the population of eligible sellers. Second, the instruments in Table~\ref{tab:instruments} act on different margins of the seller's problem, and only one of them moves the threshold outside Complete Adverse Selection.

\begin{table}[htbp]
\centering
\caption{Instruments, Parameters and Margins}
\label{tab:instruments}
\footnotesize\renewcommand{\arraystretch}{1.3}
\begin{tabular*}{\textwidth}{@{\extracolsep{\fill}}lccl@{}}
\toprule\toprule
\textbf{Instrument} & \textbf{Parameter} & \textbf{Margin it acts on} & \textbf{Moves $\bar R$?} \\
\midrule
Renovation subsidy & $S$ & renovation & acts on $\hat R-\bar R$ \\
Enforcement against hiding & $c_0$ & who hides & only in Complete Adverse Selection \\
Share of informed sellers & $\pi$ & who can renovate & only in Complete Adverse Selection \\
Attention to silent listings & $\psi$ & who hides; price of silence & only in Complete Adverse Selection \\
Full attention to every listing & $\omega_0=\omega_1=1$ & rating premium; price of silence & yes, through $b$ \\
\bottomrule\bottomrule
\end{tabular*}
\vspace{0.5em}

\raggedright
\footnotesize
\textit{Notes:} Complete Adverse Selection is the regime in which every informed low-quality seller hides. In the other regimes $\bar R=b$ (Proposition~\ref{prop:threshold}).
\end{table}

\subsection{The Renovation Problem}
\label{ss:renov_problem}

Formalizing the choice requires one addition to the model: a renovation decision $r\in\{0,1\}$ that turns an informed low-quality seller into an informed high-quality seller with a new certificate.\footnote{Uninformed sellers have no strategy set: information status is a friction rather than a choice (Section~\ref{ss:model_setup}), so they cannot display, hide or renovate, since a renovation decision presupposes a diagnostic. They enter the analysis through the composition of the silent pool and, through $\pi$, the size of the population of potential renovators.} By Lemma~\ref{lem:att_prices}, a disclosed low rating sells for $p_{\rm L}$, a disclosed high rating for $p_{\rm L}+b$ and a silent listing for $p_{\rm L}+b\tilde e$, where $b=\rho\omega_1\beta_1$ and $\tilde e$ is the effective belief of Definition~\ref{def:effective_belief}. Measure payoffs relative to a disclosed low rating and write
\[
    h\;\equiv\;b\,\tilde e-c_0
\]
for the net gain from hiding. A subsidy $S\ge0$ lowers the renovation cost $\hat R$ that the seller faces to $\hat R-S$. Table~\ref{tab:strategies} lists the four strategies. We set $c_1=0$, so a renovated seller who hides bears the same cost as before.

\begin{table}[htbp]
\centering
\caption{Strategies for an Informed Low-Quality Seller ($\theta=0$)}
\label{tab:strategies}
\small
\renewcommand{\arraystretch}{2.1}
\begin{tabular*}{\textwidth}{@{\extracolsep{\fill}}lcccc@{}}
\toprule\toprule
\textbf{Strategy} & \textbf{Renovate} & \textbf{Disclose} & \textbf{Buyer observes} & \textbf{Payoff} \\
\midrule
1.\ Disclose low & No & Yes & low rating & $0$ \\
2.\ Hide & No & No & silence & $\underbrace{b\,\tilde e-c_0}_{h}$ \\
3.\ Renovate and disclose & Yes & Yes & high rating & $\underbrace{\rho\omega_1\beta_1}_{b}-(\hat R-S)$ \\
4.\ Renovate and hide & Yes & No & silence & $\underbrace{b\,\tilde e-c_0}_{h}-(\hat R-S)$ \\
\bottomrule\bottomrule
\end{tabular*}
\vspace{0.5em}

\raggedright
\footnotesize
\textit{Notes:} Payoffs are measured relative to a disclosed low rating. $b$ is the rating premium, $h$ the net gain from hiding and $\hat R-S$ the renovation cost net of the subsidy. An informed high-quality seller earns $b$ by disclosing and $h$ by hiding. An uninformed seller has no strategy.
\end{table}

\begin{lemma}[A renovated seller behaves like an informed high type]
\label{lem:renov_behaviour}
Strategy~4 is weakly dominated by Strategy~3 if and only if $b(\tilde e-1)\le c_0$, the condition under which an informed high-quality seller discloses. This holds in Cost-Deterred Hiding and in Partial Pooling (Corollary~\ref{cor:att_hq}) and in Complete Adverse Selection when $\sigma_1^*=0$. In Complete Adverse Selection with $\sigma_1^*>0$ a renovated seller also hides.
\end{lemma}

\noindent\textit{Proof.} See Appendix~\ref{app:attention}. $\Box$

\begin{definition}[Renovation Threshold]
\label{def:threshold}
The renovation threshold is the highest renovation cost, net of subsidy, at which an informed low-quality seller prefers to renovate to her best option without renovating:
\begin{equation}
    \bar{R}\;\equiv\;\max\{b,\,h\}\;-\;\max\{0,\,h\}\;=\;b-\min\big\{\max\{h,0\},\,b\big\}\;\in\;[0,\,b].
    \label{eq:threshold}
\end{equation}
\end{definition}

\noindent Her payoff from renovating is $\max\{b,h\}-(\hat R-S)$ and her payoff from not renovating is $\max\{0,h\}$, so she renovates if and only if $\hat R-S\le\bar R$. When hiding is her best unrenovated option ($h\ge0$) and a renovated seller displays ($h\le b$), the threshold is $b-h=b-b\tilde e+c_0$, the form in which the problem first suggests itself: the market premium for verified quality, net of the gross price of silence, plus the hiding cost she no longer pays.

\begin{definition}[Minimum incentive-compatible subsidy]
\label{def:minimum_subsidy}
The minimum incentive-compatible subsidy is the smallest transfer that makes renovation weakly preferred for a seller with renovation cost $\hat R$, holding equilibrium beliefs fixed:
\begin{equation}
    S^*\;\equiv\;\max\!\left\{\hat{R}-\bar{R},\;0\right\}.
    \label{eq:subsidy}
\end{equation}
\end{definition}

\medskip\noindent\textbf{Remark ($S^*$ is incentive-compatibility, not welfare optimum).}\label{rem:welfare}
The object $S^*$ is the minimum transfer that flips the marginal informed low-quality seller's decision from not renovating to renovating, given everyone else's behaviour. It is \emph{not} a Ramsey-optimal subsidy in the sense of \citet{ramsey_taxation_1927}. A welfare problem would maximise
\[
    \max_S \; \int_X \mathbf{1}\{\hat R(X) \le \bar R + S\} \cdot \big[\Delta W(X) - \lambda S\big]\, dF(X),
\]
where $\Delta W(X)$ is the social gain per renovation (private surplus change, any energy-use externality, and the change in the transfer that buyers make to sellers who stay silent), $\lambda\ge 1$ is the marginal cost of public funds and $F$ is the within-market distribution of renovation costs. The corner solution $S^* = \hat R - \bar R$ emerges only in the limit of degenerate $F$, $\lambda=1$ and zero externality. We treat $S^*$ as the operative positive object and leave the Ramsey characterisation to future work.

\subsection{The Renovation Threshold}
\label{ss:threshold}

\begin{proposition}[Renovation Threshold under Limited Attention]
\label{prop:threshold}
Let $e(1)=(1-\pi)\bar\theta/(1-\pi\bar\theta)$ and take the equilibrium regime of Proposition~\ref{prop:att_equilibrium}. In every regime $\bar R\le b$, so the minimum subsidy satisfies $S^*\ge\max\{\hat R-b,0\}$. By regime:
\begin{enumerate}[label=(\roman*)]
    \item \textbf{Cost-Deterred Hiding and Partial Pooling.} $\bar R=b=\rho\omega_1\beta_1$. It does not depend on $c_0$, $\pi$, $\psi$ or $C^{\rm def}$, and $\partial\bar R/\partial(\rho\omega_1)=\beta_1>0$.
    \item \textbf{Complete Adverse Selection} (every informed low-quality seller hides, so $h>0$).
    \begin{enumerate}[label=(\alph*)]
        \item \emph{High-quality sellers disclose} ($\sigma_1^*=0$): $\bar R=b-h=b\,[1-\tilde e(1,0)]+c_0\in[0,b)$, with $\tilde e(1,0)=\psi e(1)+(1-\psi)e^{\rm def}$. Table~\ref{tab:threshold_cs} gives the derivatives.
        \item \emph{High-quality sellers also hide} ($\sigma_1^*>0$): $\bar R=0$.
    \end{enumerate}
\end{enumerate}
At $\psi=\omega_1=1$ the proposition is the full-attention statement: $b=\rho\beta_1$, $\tilde e=e$, and only the case in which high-quality sellers disclose exists in Complete Adverse Selection.
\end{proposition}

\begin{table}[htbp]
\centering
\caption{Derivatives of $\bar R$ in Complete Adverse Selection with Disclosing High-Quality Sellers}
\label{tab:threshold_cs}
\small\renewcommand{\arraystretch}{1.4}
\begin{tabular*}{\textwidth}{@{\extracolsep{\fill}}lcc@{}}
\toprule\toprule
\textbf{Parameter} & \textbf{Derivative of $\bar R$} & \textbf{Sign} \\
\midrule
Hiding cost $c_0$ & $1$ & $+$ \\
Relative attention $\psi$ & $b\,[e^{\rm def}-e(1)]$ & $+$ if $e^{\rm def}>e(1)$ \\
Share of informed sellers $\pi$ & $-\,b\,\psi\,e'(1)$ & $+$ \\
Default cost $C^{\rm def}$ & $\rho\omega_1(1-\psi)$ & $+$ \\
Rating weight $\rho\omega_1$ & $\beta_1\,(1-\tilde e)$ & $-$ if $\tilde e>1$ \\
\bottomrule\bottomrule
\end{tabular*}
\vspace{0.5em}

\raggedright
\footnotesize
\textit{Notes:} $\tilde e>1$ means that silence sells above a verified efficient rating, as it does in apartments. The sign of the $\psi$ derivative requires $e^{\rm def}>e(1)$.
\end{table}

\noindent\textit{Proof.} See Appendix~\ref{app:attention}. $\Box$

\medskip

The Cost-Deterred Hiding and Partial Pooling case is the central result of this section, and the mechanism is the one that Section~\ref{ss:renov_problem} makes visible. In Partial Pooling the marginal informed low-quality seller is exactly indifferent between hiding and disclosing a low rating, so $h=0$ and her best option without renovating is worth a disclosed low rating. Renovating gains her $b$ and nothing else. The hiding cost has cancelled, because the equilibrium hiding rate adjusts until the net gain from hiding is zero. In Cost-Deterred Hiding she does not hide and the same comparison applies. Raising the hiding cost therefore moves the hiding rate along the indifference condition and leaves the threshold where it was: tightening enforcement against concealment cannot buy renovation in these regimes, and renovation requires a subsidy that improves the renovate-and-disclose payoff. Complete Adverse Selection is the textbook case. When every informed low-quality seller hides and hiding pays strictly, renovating means giving up the rent $h$, the threshold falls to $b-h$, and enforcement, by lowering $h$, raises the threshold one for one. The sign of $\partial\bar R/\partial(\rho\omega_1)$ in Complete Adverse Selection is the sign of $1-\tilde e$, which is negative whenever silence sells above a verified efficient rating, as it does in apartments.

The bound $S^*\ge\max\{\hat R-b,0\}$ is the claim we rest on. The result for Cost-Deterred Hiding and Partial Pooling holds only where the market is outside Complete Adverse Selection, and which regime a market is in depends on the hiding cost. With a single cost per market, calibrated to the observed hiding rate, every market is in Partial Pooling at the estimated attention by construction (\attScBaseNPpApt{} of \attNMarketsApt{} apartment markets and \attScBaseNPpHouse{} of \attNMarketsHouse{} house markets), which follows from the calibration; the registry shows interior hiding rates in essentially every market, which in this model is possible only in Partial Pooling. Reading the hiding cost instead off each market's premium, $c_0=\hat\Delta_m$, places \attRawNCasApt{} apartment markets and \attRawNCasHouse{} house markets in Complete Adverse Selection, with mean threshold \attRawRbarRatioApt{} and \attRawRbarRatioHouse{} of $b$ (Appendix~\ref{app:attention}); we do not take that reading, because it returns negative costs in \attRawNNegApt{} apartment and \attRawNNegHouse{} house markets. The bound is what survives both readings: whatever the hiding cost, the subsidy required is at least the shortfall of $\hat R$ below the rating premium $\hat b$. Because $\bar R=\hat b$ in the regimes the data support, the minimum subsidy does not depend on the observed hiding rate or on $\hat\Delta$, which is an implication for how subsidy needs vary across markets; testing it needs programme take-up by commune, which we do not have.

\subsection{Subsidies and Disclosure in Equilibrium}
\label{ss:policy_ge}

The threshold holds beliefs fixed. A subsidy that induces renovation also removes disclosed low ratings from the market, and the silent pool responds. The next proposition lets it. Let $F$ be the share of informed low-quality sellers who renovate and $e(1;F)=(1-\pi)\bar\theta/\big(1-\pi+\pi(1-\bar\theta)(1-F)\big)$ the pool belief when every one of the remaining informed low-quality sellers hides.

\begin{proposition}[Subsidies and Disclosure]
\label{prop:subsidy_ge}
Take a market in Partial Pooling with $\psi>0$, and write $e^*$ for its pool belief~\eqref{eq:att_pp} and $\sigma_0^*$ for its hiding rate. The cost distribution $F_R$ is continuous.
\begin{enumerate}[label=(\roman*)]
    \item \emph{Phase 1.} For $F\le1-\sigma_0^*$ the mass of silent informed low-quality sellers, $(1-\pi)(\bar\theta-e^*)/e^*$ per listing, does not depend on $F$. The premium stays at $c_0$, the threshold stays at $b$, and the hiding rate among informed low-quality sellers who have not renovated rises to $\sigma_0^*/(1-F)$.
    \item \emph{Phase 2.} For $F>1-\sigma_0^*$ every remaining informed low-quality seller hides, the market is in Complete Adverse Selection, and $\bar R(F)=b-\psi b\,[e(1;F)-e^*]$, which falls as $F$ rises.
    \item \emph{Equilibrium.} Let $F_R$ be the cost distribution. The renovating share $F^*$ solves $F=F_R(\bar R(F)+S)$ and is unique. If instead all sellers share one cost $\hat R$ (a point mass, where equilibrium involves mixing at the atom), every informed low-quality seller renovates in a pure-strategy equilibrium if and only if $S\ge\hat R-b+\psi b\,(\bar\theta-e^*)$.
\end{enumerate}
\end{proposition}

\noindent\textit{Proof.} See Appendix~\ref{app:attention}. $\Box$

\medskip

Inside Partial Pooling the indifference of the marginal hider pins down the belief $e^*$, and with it the number of silent informed sellers. Renovators therefore come out of the stock of sellers who would have disclosed a low rating. A subsidy buys renovation but not disclosure until that stock is exhausted: the premium does not move, and the share of the remaining informed low-quality sellers who hide rises. The stock is large. Disclosed poor ratings are \attGeDisclosedShareApt{} percent of informed low-quality apartment listings and \attGeDisclosedShareHouse{} percent of house listings, so up to that share can renovate at the partial-equilibrium subsidy $S^*$ without any effect on hiding. Beyond it the remaining sellers hide and the threshold falls, so the subsidy needed for full take-up exceeds $S^*$ by $\psi b(\bar\theta-e^*)$, which is at most the width of the Partial Pooling band (Corollary~\ref{cor:att_band}). At the estimates this correction averages \attGeCorrPctBApt{} percent of $b$ in apartments (\euro{}\attGeCorrEurApt{} per m\textsuperscript{2}) and \attGeCorrPctBHouse{} percent in houses (\euro{}\attGeCorrEurHouse{}), so the partial-equilibrium subsidy is the general-equilibrium one to within about one percent of the rating premium. Uniqueness follows because renovation decisions are strategic substitutes through the silent pool: each renovation lowers the number of disclosed ratings and raises the gain from hiding for those left, which lowers the threshold for the next renovator. The proposition has an implication that does not rely on any calibrated quantity: in markets where a subsidy reaches more informed low-quality sellers, the registry hiding rate among those who remain should rise while the non-disclosure premium stays flat. Testing it requires programme take-up by commune and registry data after the subsidy.

% ============================================================
\subsection{Renovation Subsidies in French Housing Markets: Empirical Magnitudes and Spatial Heterogeneity}
\label{sec:policy_empirical}
% ============================================================

% [AUTHOR REWRITE] magnitudes
In Partial Pooling and Cost-Deterred Hiding the threshold equals the rating premium $b$. Because $\hat\beta_1$ measures the quality premium in log points against a local price per m\textsuperscript{2}, the nominal threshold $\bar R_m=\hat b_m\times P_m$ scales with the local price level, and high-price urban markets have the largest verified-quality premia in euros even when the proportional premium is similar across locations. Across the \attNMarketsApt{} estimable apartment markets, $\bar R_m$ averages \euro{}\attScBaseRbarApt{}/m\textsuperscript{2} (median \euro{}\attBaseRbarMedApt{}, range \euro{}\attBaseRbarMinApt{} to \euro{}\attBaseRbarMaxApt{}); across the \attNMarketsHouse{} house markets the mean is \euro{}\attScBaseRbarHouse{} and the median \euro{}\attBaseRbarMedHouse{} (Figure~\ref{fig:policy_rbar_dist}).\footnote{These figures use the rating premium itself, $\bar R=\hat b$. The exact-pricing solver of Section~\ref{sec:equilibrium}, which takes the threshold from the full-attention equilibrium, gives means of \euro{}\attRbarCoreEurTwoFiftyApt{} and \euro{}\attRbarCoreEurTwoFiftyHouse{} per m\textsuperscript{2}. The difference is consistent with linear versus exact pricing and not with attention, since it is of the same size in the two segments (Appendix~\ref{app:attention}).} Cross-market dispersion is substantial, so the markets at the upper end are intrinsically more amenable to privately funded renovation than those at the lower end.

\begin{figure}[t]
    \centering
    \caption{Distribution of the renovation threshold $\bar{R}_m$ across markets}
    \label{fig:policy_rbar_dist}
    \includegraphics[width=\linewidth]{fig/visuals/fig_att_policy_rbar_dist.pdf}
    \vspace{0.4em}
    \begin{minipage}{0.97\linewidth}
        \footnotesize\textit{Notes.} Each panel shows the cross-market distribution of $\bar{R}_m=\hat b_m P_m$ (euros per m\textsuperscript{2}) for one property type, with $P_m$ the market's median price per m\textsuperscript{2}. Vertical lines mark the renovation-cost benchmarks $\hat R=250$ and $500$. Every market is in Partial Pooling at the estimated attention, so the threshold is the rating premium throughout (Proposition~\ref{prop:threshold}).
    \end{minipage}
\end{figure}

The minimum subsidy is $S^*_m=\max\{\hat R-\bar R_m,0\}$, evaluated at two renovation-cost benchmarks: \euro{}250/m\textsuperscript{2} (a light, ``geste par geste'' intervention such as roof or wall insulation) and \euro{}500/m\textsuperscript{2} (a medium, two-measure upgrade), drawing on MaPrimeR\'enov' programme documentation.\footnote{MaPrimeR\'enov' distinguishes ``geste par geste'' interventions (insulation of walls or roof, costing roughly \euro{}100--400/m\textsuperscript{2} of treated area) from comprehensive ``r\'enovation d'ampleur'' operations (\euro{}700--1{,}200/m\textsuperscript{2} or more); the \euro{}500/m\textsuperscript{2} benchmark corresponds broadly to a two-measure upgrade.} The benchmarks are illustrative. The light benchmark corresponds to a partial upgrade and not to the full move from an E--G to an A--D rating that the premium $b$ prices, and the medium benchmark only approaches it; a comprehensive renovation at the cost range of the footnote would exceed the threshold in every market. A single $\hat R$ per market predicts that no informed low-quality seller exists where $\hat R<\hat b$, yet low ratings are disclosed in every market, so effective renovation costs exceed the benchmark in the \attBaseNoSubsidyTwoFiftyApt{} percent of apartment and \attBaseNoSubsidyTwoFiftyHouse{} percent of house markets where it falls below the premium. At $\hat R=250$, \attScBaseShareTwoFiftyApt{} percent of apartment markets require a positive subsidy; conditional on requiring one, the median $S^*_m$ is \euro{}\attBaseSMedTwoFiftyApt{}/m\textsuperscript{2}, or about \euro{}\attBaseUnitTwoFiftyApt{} per unit at the median apartment area of \attBaseAreaMedApt{}~m\textsuperscript{2}. For houses, \attScBaseShareTwoFiftyHouse{} percent of markets require support, with a conditional median of \euro{}\attBaseSMedTwoFiftyHouse{}/m\textsuperscript{2} (about \euro{}\attBaseUnitTwoFiftyHouse{} per unit at \attBaseAreaMedHouse{}~m\textsuperscript{2}). At $\hat R=500$ the shares are \attScBaseShareFiveHundredApt{} and \attScBaseShareFiveHundredHouse{} percent and the conditional medians \euro{}\attBaseSMedFiveHundredApt{} and \euro{}\attBaseSMedFiveHundredHouse{}/m\textsuperscript{2}. The cost-sharing ratio $S^*_m/\hat R$ is the natural decision variable for the regulator. Among markets requiring support its mean is \attBaseCostShareTwoFiftyApt{} percent for apartments and \attBaseCostShareTwoFiftyHouse{} percent for houses at the light benchmark, and \attBaseCostShareFiveHundredApt{} and \attBaseCostShareFiveHundredHouse{} percent at the medium one. A commitment to a fixed rate is therefore more defensible than a commitment to a fixed euro amount, because it adjusts to the depth of the intervention and to local conditions.

No flat subsidy fits every market. Covering every apartment market at the light benchmark takes \euro{}\attBaseSMaxTwoFiftyApt{} per m\textsuperscript{2}, against a conditional median of \euro{}\attBaseSMedTwoFiftyApt{}, and for houses \euro{}\attBaseSMaxTwoFiftyHouse{} against \euro{}\attBaseSMedTwoFiftyHouse{}; at the medium benchmark the figures are \euro{}\attBaseSMaxFiveHundredApt{} against \euro{}\attBaseSMedFiveHundredApt{} and \euro{}\attBaseSMaxFiveHundredHouse{} against \euro{}\attBaseSMedFiveHundredHouse{}. A uniform national subsidy is therefore poorly targeted because the cost-sharing requirement varies across markets, although $S^*$ is an incentive-compatibility object and not a welfare optimum (Remark~\ref{rem:welfare}). Figure~\ref{fig:policy_map} maps $S^*_m/\hat R$ at the light benchmark, with bubble size encoding the local threshold $\bar R_m$. A flat per-m\textsuperscript{2} grant set at the cross-market mean would over-pay markets where the private incentive already suffices and under-pay those where the gap is largest, the first through fiscal waste and the second through forgone renovations. The dispersion is systematic: the cost-sharing ratio is low where $\bar R_m$ is high, and $\bar R_m$ is high where local prices are high and the energy-efficiency premium $\hat\beta_{1,m}$ is large. The premium reflects the capitalised value of verified quality to sensitive buyers, not a freestanding measure of how many buyers are sensitive (Appendix~\ref{app:rho_identification} shows the two are not identified apart in price data alone). An efficient programme would differentiate the subsidy along these axes and co-target audit-infrastructure investment on markets where the threshold is large but verification is scarce, since $\pi$ does not change the threshold but scales the population it applies to (Section~\ref{ss:att_policy}).

\begin{figure}[htbp]
    \centering
    \caption{Spatial distribution of the minimum incentive-compatible subsidy, $\hat R=250$}
    \label{fig:policy_map}
    \includegraphics[width=\linewidth]{fig/visuals/fig_att_policy_map_R250.pdf}
    \vspace{0.4em}
    \begin{minipage}{0.97\linewidth}
        \footnotesize\textit{Notes.} $\hat R=250$~\euro{}/m\textsuperscript{2}. Left panel: apartment markets; right panel: house markets. Projection: Lambert-93 (EPSG:2154). Bubble colour encodes the cost-sharing ratio $S^*_m/\hat R$, the share of the renovation cost the regulator must cover, in fixed bins comparable across the two benchmarks (Appendix~\ref{app:attention} shows $\hat R=500$). Bubble size encodes the threshold $\bar R_m$ (euros per m\textsuperscript{2}). Markets with $S^*_m=0$ are dashed hollow circles. Near-coincident markets are offset by about 15~km for legibility.
    \end{minipage}
\end{figure}

We aggregate over the informed low-quality listings, the only ones that can renovate: a share $\pi(1-\bar\theta)$ of listings in each market, counted over the \attSampleYears{} years of listings in the estimation markets (July 2021 to December 2024) and expressed per year.\footnote{A dwelling listed more than once in the sample is counted each time, so the figure is an upper bound on the number of dwellings. Counting every listing, as an earlier version of this section did, would multiply the bill several-fold by treating verified efficient dwellings as renovators.} The largest market accounts for \attTopMarketShareApt{} percent of the apartment base, so the apartment total reflects that market's threshold more than the typical market's. At the light benchmark the bill is \euro{}\attScBaseCostTwoFiftyApt{} billion per year for apartments and \euro{}\attScBaseCostTwoFiftyHouse{} billion for houses, \euro{}\attScCostBaseTotTwoFifty{} billion in all, which, for scale, is \attScCostBaseTotTwoFiftyLo{} to \attScCostBaseTotTwoFiftyHi{} percent of MaPrimeR\'enov's annual budget of roughly \euro{}\attBudgetLo{}--\attBudgetHi{} billion. % [CHECK] cite the budget figure (Cour des comptes or ANAH) before submission.
At the medium benchmark it is \euro{}\attScCostBaseTotFiveHundred{} billion per year, \attScCostBaseTotFiveHundredLo{} to \attScCostBaseTotFiveHundredHi{} percent of the budget. The comparison is one of scale and not a share of the programme: the bill covers listings on one platform in the estimation markets, counted per listing, whereas the programme reaches all owners. It does not say the programme is cheap. The bill covers only the sellers who are in the model's renovation problem, and an extension that lets uninformed sellers acquire a certificate and then renovate would enlarge the eligible population. These are upper bounds on a per-listing basis and they ignore the extensive-margin participation decision. By Remark~\ref{rem:welfare} they are incentive-compatibility objects and not optimal transfers: with a shadow cost of public funds $\lambda>1$ the regulator would set the operative subsidy below $S^*_m$.

\subsection{Other Instruments and Counterfactuals}
\label{ss:att_policy}

% [AUTHOR REWRITE] section motivation
The subsidy is not the regulator's only instrument, and the others act on different margins of the seller's problem (Table~\ref{tab:instruments}). Enforcement against hiding raises the cost of concealment and changes who hides. Raising the share of informed sellers $\pi$ changes how many sellers can renovate at all. Two further instruments act on buyers: making silence as salient as a rating raises the attention buyers pay to a listing that shows no rating, and full attention raises the attention they pay to every listing. Three questions follow. Does each instrument move the renovation threshold $\bar R$, and with it the subsidy a seller needs? Does it change who hides? And would a simple measure such as making silence salient reduce what the state has to spend on subsidies? The next proposition answers the first two questions for each instrument, with the hiding cost held at its calibrated value, and the closing subsection returns to the third.

\begin{proposition}[Instruments]
\label{prop:instruments}
Take a market in Partial Pooling at the observed hiding rate, with beliefs held fixed in the renovation decision, and assume $e^{\rm def}\ge\bar\theta$, as in our estimates, where $C^{\rm def}=0$ gives $e^{\rm def}=\eta>1$.
\begin{enumerate}[label=(\roman*)]
    \item \emph{Raising the share of informed sellers $\pi$ to one,} with the off-path belief after silence equal to zero, leaves $\bar R=b$ and the minimum subsidy unchanged. It raises the population who can renovate from the informed to every low-quality dwelling, and removes hiding whenever $c_0\ge b(1-\psi)e^{\rm def}$, which holds in any market that starts in Partial Pooling, since with no uninformed sellers a silent listing is priced at the default.
    \item \emph{Making silence as salient as a rating} (raising $\omega_0$ to $\omega_1$, so $\psi\to1$, with $\rho$ and $\omega_1$ fixed). What a displayed rating is worth, $b$, does not change; only the price of silent listings falls, by $\lambda=(1-\hat\psi)\,b\,(e^{\rm def}-e^*)$. The threshold stays at $b$ and the minimum subsidy $S^*$ is unchanged, and with beliefs held fixed so is the share of informed low-quality sellers who renovate at any subsidy. The hiding rate falls weakly, to zero if $c_0\ge b\bar\theta$. In the equilibrium of Proposition~\ref{prop:subsidy_ge} the subsidy needed for full take-up falls by $\hat\psi b(\bar\theta-e^*)-(b\bar\theta-c_0)^+$, which is $\hat\psi b(\bar\theta-e^*)$ when the market moves to Cost-Deterred Hiding.
    \item \emph{Full attention} (every EPC-sensitive buyer attends to energy on every listing, rated or not: $\omega_0=\omega_1=1$, so also $\psi\to1$). The rating premium rises to $\rho\beta_1=b/\omega_1$. The data identify only $\hat\rho=\rho\omega_1$, so $\omega_1\in[\hat\rho,1]$ and the change in the threshold satisfies
    \[
        0\;\le\;\Delta\bar R\;\le\;\max\{\beta_1^{\rm anchor}-\hat b,\,0\}.
    \]
    The lower bound is the maintained reading $\omega_1=1$ and the upper bound corresponds to $\rho=1$. Where $\hat\rho\ge1$ the threshold is unchanged. Outside Complete Adverse Selection $\Delta\bar R=b(1/\omega_1-1)$; in a market that moves into it the threshold is $b(1/\omega_1)[1-e(1)]+c_0$, which is at least $b$.
    \item \emph{Enforcement.} A rise in $c_0$ leaves $\bar R$ unchanged outside Complete Adverse Selection and moves the market toward Cost-Deterred Hiding.
\end{enumerate}
\end{proposition}

\noindent\textit{Proof.} See Appendix~\ref{app:attention}. $\Box$

\medskip

The two attention scenarios are easiest to read side by side. Making silence as salient as a rating changes what buyers think about only when they see no rating: the price of a displayed rating is unchanged and the price of silence falls, so the sellers who hid have no reason to. Renovation pays the same $b$ as before, so the threshold does not move. Full attention also changes what buyers think about when they do see a rating: a displayed good rating is worth $\rho\beta_1$ instead of $\rho\omega_1\beta_1$. The data identify only the product $\hat\rho=\rho\omega_1$, so they cannot say how much of a low $\hat\rho$ is few EPC-sensitive buyers (then nothing changes) and how much is inattention to displayed ratings (then the premium rises a great deal). The result is a bracket between the maintained reading and an upper bound set by the energy-cost anchor. Table~\ref{tab:att_scenarios} reports the scenarios for both segments.

\begin{table}[htbp]
\centering
\caption{Policy Scenarios under Limited Attention}
\label{tab:att_scenarios}
\small\renewcommand{\arraystretch}{1.3}
\begin{tabular*}{\textwidth}{@{\extracolsep{\fill}}lcccc@{}}
\toprule\toprule
 & \textbf{$\hat\psi$ (baseline)} & \textbf{$\psi\to1$} & \textbf{$\pi\to1$} & \textbf{Full attention (bound)} \\
\midrule
\multicolumn{5}{l}{\textit{Panel A: apartments}} \\
Mean threshold $\bar R$ (\euro/m$^2$) & 122 & 122 & 122 & 239 \\
Markets in Cost-Deterred Hiding & 0 & 45 & 45 & 33 \\
Markets in Partial Pooling & 45 & 0 & 0 & 4 \\
Markets in Complete Adverse Selection & 0 & 0 & 0 & 8 \\
Mean hiding rate (\%) & 15.2 & 0.0 & 0.0 & 20.9 \\
Markets with $S^*>0$, $\hat R=250$ (\%) & 98 & 98 & 98 & 60 \\
Annual cost, $\hat R=250$ (\euro bn) & 0.12 & 0.12 & 0.23 & 0.01 \\
Annual cost, $\hat R=500$ (\euro bn) & 0.76 & 0.76 & 1.49 & 0.60 \\
\midrule
\multicolumn{5}{l}{\textit{Panel B: houses}} \\
Mean threshold $\bar R$ (\euro/m$^2$) & 128 & 128 & 128 & 227 \\
Markets in Cost-Deterred Hiding & 0 & 97 & 97 & 9 \\
Markets in Partial Pooling & 97 & 0 & 0 & 8 \\
Markets in Complete Adverse Selection & 0 & 0 & 0 & 80 \\
Mean hiding rate (\%) & 6.8 & 0.0 & 0.0 & 86.2 \\
Markets with $S^*>0$, $\hat R=250$ (\%) & 90 & 90 & 90 & 67 \\
Annual cost, $\hat R=250$ (\euro bn) & 0.33 & 0.33 & 0.82 & 0.08 \\
Annual cost, $\hat R=500$ (\euro bn) & 0.95 & 0.95 & 2.37 & 0.70 \\
\bottomrule\bottomrule
\end{tabular*}


\vspace{0.5em}
\raggedright
\footnotesize
\textit{Notes:} Hiding cost held at its calibrated value in every column. $\hat\psi$: relative attention at the estimate. $\psi\to1$: attention to silent listings raised to that of rated listings. $\pi\to1$: every dwelling's owner is informed. Full attention: all buyers attend to energy on every listing ($\omega_0=\omega_1=1$) and $\rho=1$ (the estimated $\hat\rho$ is capped at one market by market). Annual costs aggregate over informed low-quality listings (all low-quality listings under $\pi\to1$) over the \attSampleYears{} years of listings in the estimation markets. In the last column the rating premium rises while the single hiding cost is held fixed, so in the markets that move into Complete Adverse Selection every informed seller switches to hiding at once. Values are provisional pending certified energy consumption and certificate dates.
\end{table}

\paragraph{Salience changes who hides, not who renovates.}
A mandatory ``rating not provided'' field in advertisements, or a platform default that treats a missing rating as a rating, would raise attention to silent listings. At a fixed hiding rate the policy lowers the non-disclosure premium by $\hat\lambda_m$, by~\eqref{eq:att_lambda}: \attSalienceApt{} log points in apartments and \attSalienceHouse{} in houses. With the equilibrium response every market moves to Cost-Deterred Hiding. The mean hiding rate falls from \attSigmaObsApt{} to \attSigmaFullApt{} percent in apartments and from \attSigmaObsHouse{} to \attSigmaFullHouse{} percent in houses, and \attCfRevealedApt{} hidden unfavourable ratings of informed apartment sellers become visible, with no additional renovation (Table~\ref{tab:att_scenarios}). The threshold does not move, and neither does the subsidy bill. The apartment response is a knife-edge, because a single hiding cost makes the Partial Pooling band only \attBandWidthApt{} log points wide (Corollary~\ref{cor:att_band}), so we read it as a direction and not as an elasticity.

\paragraph{Raising the share of informed sellers reaches more sellers, not each seller differently.}
Doing so leaves the per-seller subsidy where it was and changes only the population: the annual bill at the light benchmark rises from \euro{}\attScBaseCostTwoFiftyApt{} to \euro{}\attScPiOneCostTwoFiftyApt{} billion in apartments and from \euro{}\attScBaseCostTwoFiftyHouse{} to \euro{}\attScPiOneCostTwoFiftyHouse{} billion in houses, and hiding disappears. The renovation flow~\eqref{eq:renovation_flow} below expresses the same point as an accounting identity.
\begin{equation}
    \text{Renovation flow} \;=\; \underbrace{\pi}_{\text{informed share}} \cdot \underbrace{(1-\bar\theta)}_{\text{low among informed}} \cdot \underbrace{F_R\big(\bar R + S\big)}_{\text{share with }\hat R\le\bar R + S}.
    \label{eq:renovation_flow}
\end{equation}
There is also a stock effect that operates only along this margin. As audits accumulate the housing stock becomes increasingly characterised by verified energy labels, the uninformed pool that sustains the non-disclosure premium shrinks, and the equilibrium drifts toward Cost-Deterred Hiding. Renovations undertaken today become permanent reclassifications from low to high quality in the verified stock.

\paragraph{Full attention.}
The full-attention result bounds how far the threshold could move if every EPC-sensitive buyer attended to energy on every listing. At the upper bound the apartment threshold rises from \euro{}\attScBaseRbarApt{} to \euro{}\attScHiRbarApt{} per m\textsuperscript{2} on average and the annual bill at the light benchmark falls from \euro{}\attScBaseCostTwoFiftyApt{} to \euro{}\attScHiCostTwoFiftyApt{} billion, and at the medium benchmark from \euro{}\attScBaseCostFiveHundredApt{} to \euro{}\attScHiCostFiveHundredApt{} billion (Figure~\ref{fig:policy_bound}). A higher rating premium also raises the gain from hiding, so \attScHiNCasApt{} of \attNMarketsApt{} apartment markets move into Complete Adverse Selection and the mean hiding rate rises to \attScHiSigmaApt{} percent. The lower bound is zero. These are bounds, not estimates: the data do not separate the share of EPC-sensitive buyers from the attention they pay to displayed ratings, and the upper bound rests on the energy-cost anchor and its ten-year horizon. In houses the same exercise moves \attScHiNCasHouse{} of \attNMarketsHouse{} markets into Complete Adverse Selection at once, with a mean hiding rate of \attScHiSigmaHouse{} percent. With one hiding cost per market every informed seller faces the same comparison, so the rise in the rating premium by a factor of about $1/\hat\rho$ pushes the gain from hiding past that cost for all of them together. With costs that differ across sellers only those whose cost lies below the new gain would start hiding. We regard the jump as a feature of the single cost and report the bound for apartments only.

\begin{figure}[htbp]
    \centering
    \caption{Bounds on the effect of forcing attention to displayed ratings, apartments}
    \label{fig:policy_bound}
    \includegraphics[width=\linewidth]{fig/visuals/fig_att_policy_bound.pdf}
    \vspace{0.4em}
    \begin{minipage}{0.97\linewidth}
        \footnotesize\textit{Notes.} Each segment is a market, sorted by the rating premium. The left end is the threshold at the maintained reading $\omega_1=1$ and the right end the upper bound $\rho=1$, in euros per m\textsuperscript{2}. Vertical lines mark the renovation-cost benchmarks. Diamonds mark markets that move into Complete Adverse Selection under the bound. The bound is not an estimate.
    \end{minipage}
\end{figure}

\paragraph{What the exercises do not show.}
Three limits qualify the results. The first is that $\bar R=b$ requires the market to be outside Complete Adverse Selection, which holds at the calibrated hiding cost by construction (Section~\ref{ss:threshold}). The second is that the zero effect of salience on renovation is an implication of one hiding cost per market. With costs that vary across sellers, an informed seller who strictly prefers hiding would lose her rent by renovating and disclosing, so making hiding less attractive could change her decision; we do not model that margin, and Section~\ref{ss:att_saving} bounds what it could change. The third is that the benchmark renovation costs are illustrative. The close of the exercise is the contrast the model delivers: salience changes who hides and, to first order, not who renovates, whereas a subsidy changes who renovates and not who hides, until the stock of disclosed poor ratings is exhausted (Proposition~\ref{prop:subsidy_ge}).

\subsection{What Salience Could Save}
\label{ss:att_saving}

% [AUTHOR REWRITE] bound on the saving from salience
The result that salience leaves the subsidy unchanged follows from the single hiding cost, and it should not be read as a finding. With one cost per market a seller in Partial Pooling is indifferent between hiding and disclosing, so hiding earns her no rent and renovating destroys none. If hiding costs differ across sellers, those whose cost lies below the gain from hiding earn a strictly positive rent that renovation gives up. Their renovation threshold is below $b$ and they need a larger subsidy. Making silence salient lowers the gain from hiding, shrinks that rent, and lowers the subsidy needed to move them, so the saving to the state is the fall in these rents. We do not observe the distribution of hiding costs and we do not model it. The next proposition bounds the saving using only what we observe, the hiding rate and the price of silence.

\begin{proposition}[Upper Bound on the Saving from Salience]
\label{prop:att_saving}
Let hiding costs be distributed across informed sellers with any continuous distribution, independently of quality, and assume $e^{\rm def}\ge\bar\theta$ and that only low-quality sellers hide. Let $\hat\sigma_0$ be the observed hiding rate and $\lambda=(1-\hat\psi)\,b\,(e^{\rm def}-\hat e)$ the fall in the price of silence when $\psi\to1$. Per eligible seller, the fall in the minimum subsidy, in log points of the price, is at most
\[
    \hat\sigma_0\;\min\{\lambda,\;b,\;\hat R\}.
\]
The saving is zero in Cost-Deterred Hiding, for any programme that reaches only the sellers who disclose today, and when all sellers share one hiding cost.
\end{proposition}

\noindent\textit{Proof.} See Appendix~\ref{app:attention}. $\Box$

\medskip

The cap at $b$ reflects that a renovated seller keeps her hiding cost ($c_1=0$): the rent that renovation destroys cannot exceed the rating premium. If the subsidy required her to display the new rating, the cap would be absent and the bound would be $\hat\sigma_0\min\{\lambda,\hat R\}$. At the estimates the bound is \euro{}\attSavingUbTwoFiftyApt{} billion per year in apartments, or \euro{}\attSavingUbDTwoFiftyApt{} billion at the light benchmark and \euro{}\attSavingUbDFiveHundredApt{} billion at the medium one if the renovated seller must display. In houses it is \euro{}\attSavingUbTwoFiftyHouse{} billion, because $\hat\psi$ is close to one and the price of silence barely falls (\attSalienceHouse{} log points). Together the bound is \attSavingUbTotTwoFiftyLo{} to \attSavingUbTotTwoFiftyHi{} percent of MaPrimeR\'enov's annual budget, and the largest apartment market accounts for \attSavingTopShareApt{} percent of it. Using the observed wedge in place of the model's price of silence gives \euro{}\attSavingUbObsTwoFiftyApt{} billion for apartments. The hiding rate behind these figures is the maintained one, \attSigmaObsApt{} percent, which assigns every informed hider to low quality and so also counts hiders whose certificate is A to D; with the rate computed from E to G certificates alone, \attSigmaRegMeanApt{} percent, the bound is \euro{}\attSavingUbRegTwoFiftyApt{} billion.

The bound is not a point estimate, and the data do not pin the saving down further. Three things would. Hiding rates by EPC letter identify more than one point of the cost distribution: the registry shows them rising from E to G in apartments, \attHideLetterEApt{}, \attHideLetterFApt{} and \attHideLetterGApt{} percent, as a gain from hiding that rises with consumption predicts. Year-by-year hiding rates, once certificate dates are re-pulled, would exploit variation in the gain over time. A direct measure of what concealment costs sellers would give the distribution itself. Until then the finding is that, even at the bound, the saving is small against the budget, and the apartment figure rests mostly on one market.

% ============================================================
\section{Conclusion}
% ============================================================

Four years after France mandated energy performance disclosure, silence is still rewarded. Properties listed without an EPC score trade at a premium of 6.9 percent above disclosed low-efficiency properties for apartments and 3.4 percent for houses. Standard unraveling theory predicts the opposite. This paper resolves the puzzle.

The resolution requires two departures from the classical framework, each empirically motivated and individually insufficient. First, not all sellers are informed. At listing time, only 51.8\% of apartment sellers and 41.7\% of house sellers hold a certified energy audit. When buyers observe a missing EPC score, they face genuine ambiguity. The non-disclosing seller may be concealing an unfavorable rating or may simply not know her property's energy performance. Uninformed sellers therefore provide strategic hiders with involuntary cover, preventing silence from being read as the worst possible signal. Second, not all buyers price energy performance into their offers. When the share of energy-sensitive buyers is moderate, the revenue penalty from concealment is diluted, making hiding profitable even against a skeptical pool. Neither departure alone is sufficient. Uninformed sellers keep silence ambiguous, but buyer heterogeneity ensures that exploiting this ambiguity is financially worthwhile. In apartment markets a third ingredient is needed, because no belief about silence can price a silent listing above a verified efficient one: buyers attend to energy much less when a listing shows no rating than when it shows one. % [AUTHOR REWRITE]


These two forces generate three equilibrium regimes. In the partial-pooling equilibrium that describes the bulk of French housing markets, the fraction of informed low-quality sellers that hides adjusts endogenously until the marginal hider is exactly indifferent between concealment and disclosure. This indifference condition delivers the paper's central theoretical result. The non-disclosure premium equals exactly the concealment cost in partial pooling, invariant to the share of informed sellers and the share of energy-sensitive buyers. We identify the concealment cost directly from the disclosed-price gap and the registry-measured belief, without needing buyer energy-sensitivity to be separately estimated, and test the indifference prediction against it rather than assume it. The two segments concur less than an earlier reading of these facts suggested. Concealment costs are 2.4\% of property value for apartments and 3.5\% for houses, of the same order across two segments that differ in buyer energy-sensitivity (49.5\% for apartments, 45.8\% for houses, identified from an external energy-cost anchor rather than from the price data used to construct the premium itself); but the indifference prediction itself holds for houses and is decisively rejected for apartments, where the non-disclosure premium exceeds what the concealment cost and the registry-measured belief predict by a wide and precisely estimated margin. One attention parameter per segment, measuring how much buyers attend to energy on listings that show no rating, accounts for both the rejection in apartments and the fit in houses (Section~\ref{ss:attention}). % [AUTHOR REWRITE]
Strategic hiding is itself uncommon. 15.2\% of informed apartment sellers with unfavorable ratings withhold them, versus 6.8\% for houses. The non-disclosure premium survives not because strategic hiding is pervasive but because genuine uninformedness is. The uninformed seller, who pays no concealment cost yet receives the same premium as the strategic hider, is the residual claimant on the information rent; the informed hider, in the segment where the model fits, merely breaks even.

The policy implications depart from textbook prescriptions in one consequential way. Although the renovation threshold mechanically incorporates the concealment cost, the seller's indifference condition causes that cost to cancel exactly in partial pooling. The renovation incentive reduces to the verified-quality premium, the price gap between high-quality and low-quality disclosed properties, and is invariant to both the concealment cost and the share of informed sellers. Tightening enforcement against non-disclosure therefore does not move the renovation threshold in the empirically dominant regime. The regulator faces three distinct instruments: renovation subsidies (closing the gap between the verified-quality premium and the renovation cost), enforcement against strategic hiding (relevant only outside partial pooling), and audit-infrastructure investment (scaling how many sellers the other two instruments reach, not the per-seller incentive itself). These act on distinct margins. Limited attention adds a further instrument: making silence as salient as a rating, which acts on buyers and lowers hiding without moving the renovation threshold, whereas a higher weight on displayed ratings could raise it, within a bound the data do not pin down (Section~\ref{ss:att_policy}). % [AUTHOR REWRITE]
 A uniform national subsidy exploits none of this structure. The minimum cost-sharing requirement varies widely, from markets where the private incentive already suffices to markets where the verified-quality premium is small, and aggregate programme costs among the sellers who can renovate run to \euro{}\attScCostBaseTotTwoFifty{} billion a year at the light renovation benchmark and \euro{}\attScCostBaseTotFiveHundred{} billion at the medium benchmark, between \attScCostBaseTotTwoFiftyLo{} and \attScCostBaseTotFiveHundredHi{} percent of the current annual budget of MaPrimeR\'enov'. Efficient targeting must condition on local market composition, directing public expenditure to markets where the private renovation gap is largest, with audit investment scaling how many sellers any subsidy reaches.

The broader lesson extends beyond energy labels. Disclosure mandates discipline strategic silence only when sellers are actually informed. When a substantial share genuinely does not know their product's quality, the mandate produces ambiguity rather than transparency. Involuntary non-disclosers shelter strategic hiders, and silence cannot be read as bad news. Where buyers do not attend to what a listing leaves out, silence is not read at all, and making absence visible is a complement to the other instruments. % [AUTHOR REWRITE]
 The policy challenge is not to punish opacity more severely but to produce knowledge in the first place. In housing energy markets, as in any market where non-disclosure mixes strategy with ignorance, audit-infrastructure investment is the instrument that expands the reach of every other tool in the regulator's toolkit.

% ============================================================
% APPENDIX
% ============================================================
\clearpage
\appendix

% =====================================================================
% APPENDIX — robustness to the quality threshold and to the sigma_1 = 0
% restriction. Referenced from §4.1.
% =====================================================================

\section{Robustness: The Quality Threshold and the $\sigma_1^*=0$ Restriction}
\label{app:robustness}

Two maintained choices in the baseline specification deserve separate scrutiny, because each
is an assumption rather than an estimate.

The first is where the binary quality space is cut. The baseline places the boundary between
D and E, so that $\theta=1$ corresponds to grades A--D. This follows the rental-ban
discontinuity of \emph{loi Climat et R\'esilience} discussed in
Section~\ref{ss:institutions}, and it is the cut for which the price-index contrast and the
listing-level quality measure refer to the same partition of grades. The natural alternative
places the boundary between C and D.

The second is the treatment of informed sellers who withhold. Because such a listing displays
no score, its quality cannot be read from the listing itself. The baseline assigns
$\theta=0$ by revealed preference, invoking Corollary~\ref{cor:hq_disclose}. That assignment
is an assumption, and it has a consequence worth stating plainly: it makes
$\hat\sigma_{1,m}$ identically zero by construction, so the corollary cannot be assessed
against the baseline estimates. The alternative reads the withheld grade from the ADEME
registry, which records it for every matched listing, leaving $\hat\sigma_{1,m}$ measured.

Table~\ref{tab:robustness} crosses the two choices.

\begin{table}[H]
\centering
\caption{Sensitivity to the quality threshold and to the $\sigma_1$ assignment}
\label{tab:robustness}
\small\renewcommand{\arraystretch}{1.25}
\begin{tabular*}{\textwidth}{@{\extracolsep{\fill}}l*{6}{c}@{}}
\toprule\toprule
\textbf{Specification} & $N$ & $\hat\rho_m$ & $\hat c_{0,m}$ & $\hat\sigma_{0,m}$ & $\hat\sigma_{1,m}$ & $\hat{\bar\theta}_m$ \\
\midrule
\multicolumn{7}{l}{\textit{Panel A: Apartments}} \\[2pt]
A--D vs E/F/G, $\sigma_1$ imposed \textbf{(baseline)} & 45 & 90.9\% & 3.42\% & 14.0\% & --- & 62.5\% \\
A--D vs E/F/G, $\sigma_1$ from registry               & 45 & 90.8\% & 3.75\% &  5.5\% & 5.1\% & 65.9\% \\
A--C vs E/F/G, $\sigma_1$ imposed                     & 40 & 95.0\% & 1.62\% &  7.1\% & --- & 26.5\% \\
A--C vs E/F/G, $\sigma_1$ from registry               & 40 & 95.0\% & 1.75\% &  4.8\% & 6.5\% & 28.4\% \\
\midrule
\multicolumn{7}{l}{\textit{Panel B: Houses}} \\[2pt]
A--D vs E/F/G, $\sigma_1$ imposed \textbf{(baseline)} & 97 & 68.1\% & 4.45\% &  5.7\% & --- & 60.5\% \\
A--D vs E/F/G, $\sigma_1$ from registry               & 97 & 67.7\% & 4.59\% &  2.2\% & 2.1\% & 61.9\% \\
A--C vs E/F/G, $\sigma_1$ imposed                     & 100 & 71.8\% & 1.80\% &  3.0\% & --- & 28.5\% \\
A--C vs E/F/G, $\sigma_1$ from registry               & 100 & 71.7\% & 1.88\% &  1.9\% & 2.5\% & 29.3\% \\
\bottomrule\bottomrule
\end{tabular*}
\vspace{0.5em}
\begin{minipage}{\linewidth}
\footnotesize\textit{Notes.} Market-level means. ``$\sigma_1$ imposed'' assigns $\theta=0$ to
informed non-disclosers by revealed preference, which forces $\hat\sigma_{1,m}=0$ in every
market; the column is left blank there because the quantity is assumed rather than estimated.
``$\sigma_1$ from registry'' binarises the withheld grade from the ADEME record. The quality
threshold also governs the price-index contrast: A--D vs E/F/G merges the \emph{Efficient} and
\emph{Balanced} index series by transaction volume, while A--C vs E/F/G contrasts
\emph{Efficient} against \emph{Intensive} and omits \emph{Balanced}. Market counts differ
across rows because the minimum-volume filter is applied to the high-group series, which is
larger under the A--D contrast.
\end{minipage}
\end{table}

Three results follow.

First, the structural estimates are insensitive to the $\sigma_1$ assignment. Holding the
threshold fixed, relaxing the restriction moves $\hat\rho_m$ by less than a tenth of a
percentage point in both segments, and moves $\hat c_{0,m}$ from 3.42\% to 3.75\% for
apartments and from 4.45\% to 4.59\% for houses. The concealment cost is therefore not an
artefact of the revealed-preference assignment, even though $\hat\sigma_{0,m}$ itself is
sensitive to it: imposing $\theta=0$ places every informed non-discloser in the low-quality
group and roughly doubles the measured hiding rate.

Second, the restriction $\sigma_1^*=0$ is not satisfied exactly, but the model's predicted
\emph{ordering} holds under the maintained threshold. With quality read from the registry and
the cut at the rental-ban boundary, $\hat\sigma_{1,m} = 5.1\%$ against
$\hat\sigma_{0,m} = 5.5\%$ for apartments, and $2.1\%$ against $2.2\%$ for houses: informed
sellers of low-quality properties withhold more often than informed sellers of high-quality
ones, as the model requires, though the gap is narrow. Under the C--D cut the ordering
reverses in both segments ($6.5\%$ against $4.8\%$, and $2.5\%$ against $1.9\%$). The
threshold that the institutional setting selects is thus also the one under which the
behavioural restriction survives, which we read as support for the baseline cut rather than
as independent confirmation of the corollary.

Third, the threshold matters a great deal for the level of the concealment cost, and this is
the paper's most consequential specification choice. Moving the cut from C--D to D--E raises
$\hat c_{0,m}$ from roughly 1.6--1.8\% of property value to 3.4--4.6\%. The mechanism is
transparent: with A--D counted as high quality, the average quality of the non-disclosing
pool $\hat e^*_m$ roughly doubles, and $c_{0,m} = \rho_m\beta_{1,m}e^*_m$ scales with it.
Readers who prefer the narrower definition of ``high quality'' should therefore read the
concealment cost as the smaller figure throughout; the qualitative results --- the ordering
across segments, the boundary saturation in $\hat\rho_m$, and the regime classification ---
are unaffected.


\section{Figures}\label{sec:appendix_figures}

\begin{figure}[htbp]
    \centering
    \caption{Visual representation of the model dynamics}
    \includegraphics[width=1\textwidth]{fig/visuals/model_mechanics_overview_h.png} 
    \label{fig:vis_model}
\end{figure}

\section{Data and Model Objects}

Our dataset includes residential property listings from various French housing markets. For each listing $i$ in market $m$, we observe:

\begin{itemize}
\item $p$: listing price
\item $X$: vector of property characteristics (such as size, age, location, etc.)
\item $d \in {0,1}$: indicator of whether the Energy Performance Certificate (EPC) score is displayed in the listing
\item $\text{EPC}^{\text{listing}}$: the EPC score displayed on the listing (if $d_i = 1$)
\end{itemize}

What sets this approach apart is our ability to match these listings with data from ADEME (the French Environment and Energy Management Agency), which maintains an official registry of certified energy audits. These audits are reported by third-party assessors who are presumed to be honest and impartial. We match listings to the ADEME registry based on key characteristics that are registered in both databases, such as addresses, floor, item type, and area\footnote{We provide a full explanation of our matching strategy in the appendix}. This allows us to assess the accuracy of the EPC scores displayed in listings by comparing them to the official certified energy audits. This allows us to observe notably:

\begin{itemize}
    \item $K \in \{0,1\}$: whether a verified EPC exists at listing time
    \item $\theta$: the true (verified) EPC score (if $K = 1$)
\end{itemize}

This matching generates four mutually exclusive groups:

\begin{table}[htbp]
\centering
\caption{Classification of Listings}
\label{tab:listing_types}
\renewcommand{\arraystretch}{1.4}
\begin{tabular}{lccl}
\toprule
\toprule
\textbf{Group} & \textbf{Displays EPC} ($d$) & \textbf{Verified EPC Exists} ($K$) & \textbf{Interpretation} \\
\midrule
Truthful Disclosure & 1 & 1 & Informed seller reveals $\theta$ \\
Strategic Hiding & 0 & 1 & Informed seller conceals $\theta$ \\
Uninformed & 0 & 0 & Seller lacks information \\
Falsification & 1 & 0 or mismatch & Misreporting (excluded) \\
\bottomrule
\bottomrule
\end{tabular}
\end{table}

The key identification advantage is that we directly observe the latent variable $K$ (information status), which the model treats as unobservable to buyers. This allows us to separately identify the behavior of informed and uninformed sellers.


\section{Proofs}
\label{app:proofs}

All results below adopt the benchmark $\alpha_k=\beta_k$.

%------------------------------------------------------------
\subsection*{Preliminary: Optimal Price Sensitivity to Buyer Valuation}

\begin{proof}[Optimal price sensitivity to shifts in buyer valuation]

The optimal price $p^*$ satisfies the first-order condition:
\begin{equation}
    \frac{\gamma_1}{\gamma_0 + \gamma_1 \cdot p^*} = 1 - G(v_1 - p^*)
\end{equation}

Define the left-hand side as $\mathrm{LHS}(p) = \frac{\gamma_1}{\gamma_0 + \gamma_1 \cdot p}$ and the right-hand side as $\mathrm{RHS}(p; v_1) = 1 - G(v_1 - p)$. The optimal price is determined by the intersection $\mathrm{LHS}(p^*) = \mathrm{RHS}(p^*; v_1)$.

\textbf{Monotonicity of each side in $p$.}

The left-hand side is strictly decreasing in $p$:
\begin{equation}
    \frac{\partial \mathrm{LHS}}{\partial p} = -\frac{\gamma_1^2}{(\gamma_0 + \gamma_1 \cdot p)^2} < 0
\end{equation}

The right-hand side is strictly increasing in $p$. Since $G$ is the logistic CDF, $G'(x) = G(x)(1 - G(x)) > 0$ for all $x$. Therefore:
\begin{equation}
    \frac{\partial \mathrm{RHS}}{\partial p} = G'(v_1 - p) > 0
\end{equation}

A strictly decreasing function intersects a strictly increasing function at most once, so $p^*$ is unique.

\textbf{Effect of $v_1$ on the right-hand side.}

At any fixed $p$, the right-hand side is strictly decreasing in $v_1$:
\begin{equation}
    \frac{\partial \mathrm{RHS}}{\partial v_1} = -G'(v_1 - p) < 0
\end{equation}

An increase in $v_1$ shifts the RHS curve downward.

\textbf{Comparative static.}

Using the chain rule, differentiating the equilibrium condition with respect to $v_1$:
\begin{equation}
    \frac{\partial \mathrm{LHS}}{\partial p} \cdot \frac{\partial p^*}{\partial v_1} = \frac{\partial \mathrm{RHS}}{\partial p} \cdot \frac{\partial p^*}{\partial v_1} + \frac{\partial \mathrm{RHS}}{\partial v_1}
\end{equation}

Rearranging:
\begin{equation}
    \frac{\partial p^*}{\partial v_1} = \frac{\frac{\partial \mathrm{RHS}}{\partial v_1}}{\frac{\partial \mathrm{LHS}}{\partial p} - \frac{\partial \mathrm{RHS}}{\partial p}} = \frac{-G'(v_1 - p^*)}{-\frac{\gamma_1^2}{(\gamma_0 + \gamma_1 \cdot p^*)^2} - G'(v_1 - p^*)} > 0
\end{equation}

The numerator is negative. The denominator is negative (sum of two negative terms). Therefore $\frac{\partial p^*}{\partial v_1} > 0$.
\end{proof}

%------------------------------------------------------------
\subsection*{Lemma A.1 (Profit is Increasing in Sensitive-Buyer Valuation)}
\label{lem:profit}

\textit{Holding $v_0$ and $\rho$ fixed, the gross profit function $\tilde{\Pi}(v_1; v_0, \rho)$ is strictly increasing in $v_1$ whenever $\rho > 0$.}

\begin{proof}[Proof of Lemma A.1]
Let $p_{\mathrm{list}}^*(v_1)$ denote the optimal price given sensitive-buyer valuation $v_1$. By the envelope theorem, the effect of a small increase in $v_1$ on maximized profit is:
\begin{equation}
    \frac{\partial \tilde{\Pi}}{\partial v_1} = \phi(p_{\mathrm{list}}^*) \cdot \rho \cdot G'(v_1 - p_{\mathrm{list}}^*)
\end{equation}
where we evaluate the partial derivative of the objective holding $p$ fixed at the optimum.

All three terms are positive: $\phi(p^*) > 0$ (expected transaction price is positive), $\rho \in [0,1] \geq 0$ (by assumption), and $G'(\cdot) > 0$ (logistic density is everywhere positive). Therefore $\frac{\partial \tilde{\Pi}}{\partial v_1} > 0$.
\end{proof}

Since $\tilde{\Pi}$ is increasing in $v_1$, the profit gain from disclosure is positive if and only if $v_1(d=1) > v_1(d=0)$: sensitive buyers value the property more when quality is revealed than when it is hidden. In this case, disclosure raises profit, and the seller discloses (absent hiding costs). Conversely, the profit gain is negative if $v_1(d=1) < v_1(d=0)$: sensitive buyers value the property more under hiding. This occurs when buyers' inference about non-disclosers is more favorable than the truth. In this case, hiding raises profit, and the seller hides (absent sufficiently high hiding costs).

%------------------------------------------------------------
\subsection*{Proof of Proposition \ref{prop:premium} (Disclosure Incentives and Seller Sorting)}

\begin{proof}
By Lemma A.1, $p^*_0 > p^*_1$ if and only if $v_1(d=0) > v_1(d=1, \theta)$. Since $v_1(d=0) - v_1(d=1, \theta) = \beta_1 \cdot (\mathbb{E}[\theta \mid d=0] - \theta)$ and $\beta_1 > 0$, this holds if and only if $\mathbb{E}[\theta \mid d=0] > \theta$.

\begin{itemize}
    \item \textbf{For a low-quality seller} ($\theta = 0$): the condition is $\mathbb{E}[\theta \mid d=0] > 0$, which holds whenever the non-disclosing pool contains any high-quality sellers.
    \item \textbf{For a high-quality seller} ($\theta = 1$): the condition is $\mathbb{E}[\theta \mid d=0] > 1$, which is impossible since $\theta \in \{0,1\}$.
\end{itemize}
\end{proof}

%------------------------------------------------------------
\subsection*{Proof of Corollary \ref{cor:hq_disclose} (High-Quality Sellers Always Disclose)}

\begin{proof}
From Proposition \ref{prop:premium}, the non-disclosure premium exists if and only if $\mathbb{E}[\theta \mid d=0] > \theta$.

For a high-quality seller ($\theta = 1$), this condition requires $\mathbb{E}[\theta \mid d=0] > 1$. But $\theta \in \{0, 1\}$, so $\mathbb{E}[\theta \mid d=0] \leq 1$ always. Therefore the condition never holds: high-quality sellers never earn a premium from hiding.

In fact, the inequality is strict whenever the non-disclosing pool contains any low-quality sellers. In that case $\mathbb{E}[\theta \mid d=0] < 1$, so:
\begin{equation}
    v_1(d=1, \theta=1) - v_1(d=0) = \beta_1 \cdot (1 - \mathbb{E}[\theta \mid d=0]) > 0
\end{equation}

Disclosure yields a strictly higher valuation, and therefore strictly higher profit. Adding hiding costs $c(\theta=1) = c_0 + c_1 \geq 0$ only reinforces this: disclosure dominates hiding.
\end{proof}

%------------------------------------------------------------
\subsection*{Proof of Proposition \ref{prop:equilibrium} (Equilibrium Existence and Uniqueness)}

\begin{proof}
By Corollary~\ref{cor:hq_disclose}, $\sigma_1^* = 0$ in any equilibrium. The strategic problem reduces to characterising $\sigma_0^* \in [0,1]$, the hiding probability of informed low-quality sellers. The incentive to hide is
\[
    \Phi(\sigma_0) \;\equiv\; \Delta\!\left(e(\sigma_0)\right) - c_0
    \;=\; \frac{\rho\,\beta_1\,(1-\pi)\bar{\theta}}{\pi(1-\bar{\theta})\sigma_0 + (1-\pi)} - c_0,
\]
where $e(\sigma_0)$ is the equilibrium pool belief (Equation~\eqref{eq:belief_simplified}). Since $e(\sigma_0)$ is strictly decreasing in $\sigma_0$ (Section~\ref{sec:beliefs}) and $\Delta$ is strictly increasing in $e$, $\Phi$ is \emph{strictly decreasing} in $\sigma_0$, with boundary values
\[
    \Phi(0) = \Delta(\bar{\theta}) - c_0, \qquad
    \Phi(1) = \Delta\!\left(\tfrac{(1-\pi)\bar{\theta}}{1-\pi\bar{\theta}}\right) - c_0.
\]

\medskip
\textbf{Case 1: Cost-Deterred Hiding.} Suppose $c_0 > \Delta(\bar{\theta})$, i.e., $\Phi(0) < 0$.

Since $\Phi$ is strictly decreasing, $\Phi(\sigma_0) \leq \Phi(0) < 0$ for all $\sigma_0 \in [0,1]$. Hiding is strictly unprofitable at every hiding rate, including the most favourable case $\sigma_0 = 0$. Disclosure strictly dominates for every informed low-quality seller, so the unique equilibrium is $\sigma_0^* = 0$ with equilibrium belief $e^* = \bar{\theta}$.

\textit{Interpretation.} All informed low-quality sellers disclose; the non-disclosing pool consists exclusively of uninformed sellers. Buyers correctly infer quality $\bar{\theta}$ from non-disclosure. Unlike the classical Grossman--Milgrom outcome, silence is uninformative rather than damning: non-disclosure carries no stigma because it reveals nothing beyond the seller's lack of a certified audit.

\medskip
\textbf{Case 2: Complete Adverse Selection.} Suppose $c_0 < \Delta(e(1))$, i.e., $\Phi(1) > 0$.

Since $\Phi$ is strictly decreasing, $\Phi(\sigma_0) \geq \Phi(1) > 0$ for all $\sigma_0 \in [0,1]$. Hiding is strictly profitable at every hiding rate, including the least favourable case $\sigma_0 = 1$. Hiding is therefore strictly dominant for every informed low-quality seller regardless of the strategies of others, so the unique equilibrium is $\sigma_0^* = 1$ with equilibrium belief $e^* = (1-\pi)\bar{\theta}/(1-\pi\bar{\theta})$.

\textit{Interpretation.} All informed low-quality sellers pool with uninformed sellers. The non-disclosure premium $\rho\cdot\beta_1\cdot e(1)$ is the maximum attainable across all regimes (Equation~\eqref{eq:premium}).

\medskip
\textbf{Case 3: Partial Pooling.} Suppose $\Delta(e(1)) \leq c_0 \leq \Delta(\bar{\theta})$, i.e., $\Phi(1) \leq 0 \leq \Phi(0)$.

$\Phi$ is continuous and strictly decreasing on $[0,1]$, with $\Phi(0) \geq 0$ and $\Phi(1) \leq 0$. By the intermediate value theorem, there exists $\sigma_0^* \in [0,1]$ with $\Phi(\sigma_0^*) = 0$; strict monotonicity of $\Phi$ implies uniqueness. When the boundary inequalities hold strictly, $\sigma_0^*$ lies in the open interval $(0,1)$; at the boundaries, $\sigma_0^* \in \{0,1\}$ and the regime collapses to Case 1 or Case 2.

At $\sigma_0^*$, informed low-quality sellers are exactly indifferent between hiding and disclosing; mixing at rate $\sigma_0^*$ is self-confirming through the belief function. This regime generates the closed-form hiding rate and is the empirically dominant regime in the estimation sample.

\medskip
\textbf{Uniqueness and exhaustiveness.} In Cases 1 and 2, uniqueness follows from strict dominance: $\Phi$ has constant sign on $[0,1]$ so no other $\sigma_0$ is consistent with optimality. In Case 3, uniqueness follows from the unique zero of $\Phi$. The three conditions $c_0 > \Delta(\bar{\theta})$, $c_0 < \Delta(e(1))$, and $\Delta(e(1)) \leq c_0 \leq \Delta(\bar{\theta})$ partition $\mathbb{R}_{+}$ since $\Delta(\bar{\theta}) \geq \Delta(e(1))$ (as $e(0) = \bar{\theta} \geq e(1)$ and $\Delta$ is increasing). Equilibrium exists and is unique for all parameter values.
\end{proof}

%------------------------------------------------------------
\subsection*{Proof of Non-Disclosure Premium Characterisation \eqref{eq:premium}}

\begin{proof}
Under the benchmark $\alpha_k=\beta_k$, the equilibrium price of a listing equals its expected buyer valuation. For a low-quality seller who discloses ($d=1$, $\theta=0$), both buyer types value the property at $\beta_0$ (sensitive buyers: $v_1(d=1,\theta=0)=\beta_0+\beta_1\cdot 0=\beta_0$; insensitive buyers: $v_0=\beta_0$), so $p_1^*=\beta_0$. For a non-disclosing listing, sensitive buyers assign expected quality $e(\sigma_0^*)$ and value the property at $\beta_0+\beta_1 e(\sigma_0^*)$, while insensitive buyers value it at $\beta_0$. Weighting by buyer composition yields $p_0^*=\beta_0+\rho\,\beta_1\,e(\sigma_0^*)$, so the non-disclosure premium is:
\begin{equation}
    p_0^*-p_1^* = \rho\,\beta_1\,e(\sigma_0^*). \label{eq:premium_derived}
\end{equation}
The three cases of~\eqref{eq:premium} follow by substituting the equilibrium $\sigma_0^*$ from each regime of Proposition~\ref{prop:equilibrium} into~\eqref{eq:premium_derived}.

\medskip
\textbf{Case 1: Cost-Deterred Hiding ($\sigma_0^*=0$).}
Evaluating~\eqref{eq:belief_simplified} at $\sigma_0=0$:
\[
    e(0) = \frac{(1-\pi)\bar{\theta}}{(1-\pi)} = \bar{\theta}.
\]
Substituting into~\eqref{eq:premium_derived}: $p_0^*-p_1^* = \rho\,\beta_1\,\bar{\theta}$.

\medskip
\textbf{Case 2: Complete Adverse Selection ($\sigma_0^*=1$).}
With $\sigma_1^*=0$ by Corollary~\ref{cor:hq_disclose}, evaluating~\eqref{eq:belief_simplified} at $\sigma_0=1$:
\[
    e(1) = \frac{(1-\pi)\bar{\theta}}{\pi(1-\bar{\theta})+(1-\pi)} = \frac{(1-\pi)\bar{\theta}}{1-\pi\bar{\theta}}.
\]
Substituting into~\eqref{eq:premium_derived}: $p_0^*-p_1^* = \rho\,\beta_1\,\dfrac{(1-\pi)\bar{\theta}}{1-\pi\bar{\theta}}$.

\medskip
\textbf{Case 3: Partial Pooling ($\sigma_0^*\in(0,1)$).}
By Proposition~\ref{prop:equilibrium}, $\sigma_0^*$ solves the indifference condition $\Delta(e(\sigma_0^*))=c_0$, i.e., $\rho\,\beta_1\,e(\sigma_0^*)=c_0$. But the left-hand side is $p_0^*-p_1^*$ by~\eqref{eq:premium_derived}, so $p_0^*-p_1^*=c_0$.

\medskip
\textbf{Ordering.}
Since $e$ is strictly decreasing by Lemma~\ref{lem:belief_monotone}, we have $e(1)<e(\sigma_0^*)<e(0)$ for all $\sigma_0^*\in(0,1)$, hence $\rho\,\beta_1\,e(1)<c_0<\rho\,\beta_1\,\bar{\theta}$, confirming the strict ordering stated after~\eqref{eq:premium}.
\end{proof}

%------------------------------------------------------------
\subsection*{Proof of Proposition~\ref{prop:comparative_statics} (Comparative statics of the non-disclosure premium)}

\begin{proof}
From~\eqref{eq:premium_derived}, the non-disclosure premium is $\Delta\equiv p_0^*-p_1^*=\rho\,\beta_1\,e(\sigma_0^*)$, where the belief $e(\sigma_0)$ is given by~\eqref{eq:belief_simplified}. The equilibrium hiding rate $\sigma_0^*\in[0,1]$ is determined by the regime characterised in Proposition~\ref{prop:equilibrium}. We derive $\partial\Delta/\partial\pi$ and $\partial\Delta/\partial\rho$ in each of the three regimes, in turn.

\medskip
\textbf{Case 1: Cost-Deterred Hiding ($\sigma_0^*=0$).}
No informed low-quality seller hides, so the non-disclosing pool consists entirely of uninformed sellers. From~\eqref{eq:belief_simplified} at $\sigma_0=0$:
\begin{equation}
    e(0) = \frac{(1-\pi)\bar{\theta}}{(1-\pi)} = \bar{\theta},
\end{equation}
which is independent of $\pi$. The premium is therefore:
\begin{equation}
    \Delta = \rho\,\beta_1\,\bar{\theta}.
\end{equation}
Since $\bar{\theta}$ is a structural parameter, $\partial\Delta/\partial\pi=0$ and $\partial\Delta/\partial\rho=\beta_1\bar{\theta}>0$.

\medskip
\textbf{Case 2: Complete Adverse Selection ($\sigma_0^*=1$).}
All informed low-quality sellers hide. From~\eqref{eq:belief_simplified} at $\sigma_0=1$:
\begin{equation}
    e(1) = \frac{(1-\pi)\bar{\theta}}{\pi(1-\bar{\theta})+(1-\pi)} = \frac{(1-\pi)\bar{\theta}}{1-\pi\bar{\theta}},
\end{equation}
so $\Delta=\rho\,\beta_1\,e(1)$. Differentiating with respect to $\pi$ by the quotient rule:
\begin{equation}
    \frac{\partial e(1)}{\partial\pi}
        = \frac{-\bar{\theta}(1-\pi\bar{\theta})-(1-\pi)\bar{\theta}\cdot(-\bar{\theta})}{(1-\pi\bar{\theta})^2}
        = \frac{\bar{\theta}(\bar{\theta}-1)}{(1-\pi\bar{\theta})^2}
        = \frac{-\bar{\theta}(1-\bar{\theta})}{(1-\pi\bar{\theta})^2} < 0,
\end{equation}
hence $\partial\Delta/\partial\pi = \rho\,\beta_1\cdot\partial e(1)/\partial\pi < 0$.
Differentiating with respect to $\rho$ (with $e(1)$ determined solely by $\pi$ and $\bar\theta$):
$\partial\Delta/\partial\rho=\beta_1\,e(1)>0$.

\medskip
\textbf{Case 3: Partial Pooling ($\sigma_0^*\in(0,1)$).}
The interior equilibrium satisfies $c_0=\rho\,\beta_1\,e(\sigma_0^*)$, so the premium is pinned at the hiding cost:
\begin{equation}
    \Delta = c_0.
\end{equation}
Changes in $\pi$ or $\rho$ induce compensating adjustments in $\sigma_0^*$ that restore the indifference condition, leaving the premium invariant: $\partial\Delta/\partial\pi=0$ and $\partial\Delta/\partial\rho=0$.

To verify: from~\eqref{eq:belief_simplified}, $\partial e(\sigma_0)/\partial\pi = -\bar{\theta}(1-\bar{\theta})\sigma_0/[\pi(1-\bar{\theta})\sigma_0+(1-\pi)]^2 < 0$ for $\sigma_0>0$. Total differentiation of $c_0=\rho\,\beta_1\,e(\sigma_0^*(\pi),\pi)$ with respect to $\pi$ gives:
\begin{equation}
    0 = \rho\,\beta_1\left(\frac{\partial e}{\partial\sigma_0}\,\frac{\mathrm{d}\sigma_0^*}{\mathrm{d}\pi} + \frac{\partial e}{\partial\pi}\right),
\end{equation}
so $\mathrm{d}\sigma_0^*/\mathrm{d}\pi = -(\partial e/\partial\pi)/(\partial e/\partial\sigma_0) < 0$: as $\pi$ rises, informed sellers populate the non-disclosing pool more heavily, so fewer low-quality sellers need to hide to sustain the same buyer belief, and $\sigma_0^*$ falls.

\medskip
\textbf{Summary.} The non-disclosure premium is \emph{weakly} decreasing in $\pi$ and \emph{weakly} increasing in $\rho$, with strict monotonicity only in the Complete Adverse Selection regime. In Case~1, the premium is determined by the population quality $\bar{\theta}$ of the uninformed pool and is insensitive to market fundamentals $\pi$ and $\rho$. In Case~3, the premium is pinned at the primitive $c_0$ by the indifference condition; it is the hiding rate $\sigma_0^*$, not the premium, that absorbs variation in $\pi$ and $\rho$.
\end{proof}

%------------------------------------------------------------
\subsection*{Proof of Corollary~\ref{cor:rho} (Effect of $\rho$ on the non-disclosure premium)}

\begin{proof}
\textbf{Part~(i).}
From~\eqref{eq:premium_derived}, the non-disclosure premium is $\Delta = \rho\,\beta_1\,e(\sigma_0^*)$. We differentiate with respect to $\rho$ in each regime, holding the regime-specific equilibrium $\sigma_0^*$ fixed where appropriate.

\medskip
\textit{Cost-Deterred Hiding ($\sigma_0^* = 0$).} Here $e(0) = \bar{\theta}$ by~\eqref{eq:belief_simplified}, independent of $\rho$. Therefore $\partial\Delta/\partial\rho = \beta_1\bar{\theta} > 0$.

\textit{Complete Adverse Selection ($\sigma_0^* = 1$).} Here $e(1) = (1-\pi)\bar{\theta}/(1-\pi\bar{\theta})$, also independent of $\rho$. Therefore $\partial\Delta/\partial\rho = \beta_1\,e(1) > 0$.

\textit{Partial Pooling ($\sigma_0^*\in(0,1)$).} The indifference condition $\rho\,\beta_1\,e(\sigma_0^*) = c_0$ pins $\Delta = c_0$. Changes in $\rho$ induce compensating adjustments in $\sigma_0^*$, as established in Proposition~\ref{prop:comparative_statics}(ii): $\partial\sigma_0^*/\partial\rho > 0$. The premium satisfies $\partial\Delta/\partial\rho = 0$.

\medskip
\textbf{Part~(ii).}
From Proposition~\ref{prop:equilibrium}, the \emph{Partial Pooling} regime requires $\Delta(e(1)) \leq c_0 \leq \Delta(\bar{\theta})$, i.e., $\rho\,\beta_1\,e(1) \leq c_0 \leq \rho\,\beta_1\,\bar{\theta}$. Differentiating both bounds with respect to $\rho$: $\partial[\rho\,\beta_1\,e(1)]/\partial\rho = \beta_1\,e(1) > 0$ and $\partial[\rho\,\beta_1\,\bar{\theta}]/\partial\rho = \beta_1\,\bar{\theta} > 0$. Both boundaries increase with $\rho$. The width of the partial-pooling interval is $\rho\,\beta_1\,[\bar{\theta} - e(1)]$, which is increasing in $\rho$ since $\bar{\theta} > e(1)$ whenever $\pi > 0$. The regime is therefore weakly expanding in $\rho$.
\end{proof}

%------------------------------------------------------------
\subsection*{Proof of Proposition \ref{prop:pi} (Equilibrium effect of $\pi$: equilibrium moves toward full unraveling)}

\begin{proof}
From Eq.~\eqref{eq:belief_simplified}, the equilibrium belief under partial separation is:
\begin{equation}
    e(1) = \frac{(1-\pi)\bar{\theta}}{1 - \pi\bar{\theta}}
\end{equation}

Taking the derivative:
\begin{equation}
    \frac{\partial e(1)}{\partial \pi} = \frac{-\bar{\theta}(1-\pi\bar{\theta}) - (1-\pi)\bar{\theta}(-\bar{\theta})}{(1-\pi\bar{\theta})^2} = \frac{-\bar{\theta}(1-\bar{\theta})}{(1-\pi\bar{\theta})^2} < 0
\end{equation}

As $\pi$ increases, $e(1)$ falls: the pool quality under partial separation deteriorates because fewer uninformed sellers are available to provide cover.

Since $\Delta(e)$ is increasing in $e$, $\Delta(e(1))$ falls. The condition for partial separation, $c_0 < \Delta(e(1))$, becomes harder to satisfy. For fixed $c_0$, an increase in $\pi$ can push the equilibrium from partial separation into the mixed or full unraveling region.
\end{proof}

When most sellers are informed ($\pi$ high), non-disclosure is suspicious. Buyers recognize that uninformed sellers are rare, so most non-disclosers must be strategic hiders. The favorable cover provided by uninformed sellers evaporates. Low-quality sellers find hiding less attractive, and the market approaches classical unraveling.

When most buyers care about quality ($\rho$ high), the stakes of disclosure are high. Low-quality sellers have much to lose from revealing their type. The non-disclosure premium is large, making hiding attractive even when pool quality is low. Partial separation is sustained, and by Corollary~\ref{cor:rho}(ii), the partial-pooling interval itself widens with $\rho$.

%------------------------------------------------------------
\subsection*{Proof of Proposition \ref{prop:c0} (Equilibrium effect of $c_0$: equilibrium moves toward full unraveling)}

\begin{proof}
An increase in $c_0$ shifts the function $\Phi(\sigma_0) = \Delta(e(\sigma_0)) - c_0$ downward. The boundaries become:
\begin{equation}
    \begin{split}
        \Phi(0) &= \Delta(\bar{\theta}) - c_0 \quad \text{(decreases)} \\
        \Phi(1) &= \Delta(e(1)) - c_0 \quad \text{(decreases)}
    \end{split}
\end{equation}

The condition for full unraveling, $c_0 > \Delta(\bar{\theta})$, becomes easier to satisfy. The condition for partial separation, $c_0 < \Delta(e(1))$, becomes harder to satisfy.
\end{proof}

Higher hiding costs make concealment less attractive. Regulatory enforcement, administrative burden, or reputational concerns that raise $c_0$ push the market toward disclosure.

% ============================================================
% Appendix: Limited Attention to Silence -- proofs and estimation details (Section 4.6)
% DRAFT inserted 2026-10-08 from playground/notes.md Part IV (plan harmonic-munching-lecun.md).
% ============================================================

\section{Limited Attention to Silence: Proofs and Estimation}
\label{app:attention}

% ============================================================
% Rewritten fragment of sections/A-attention.tex, lines 9--106
% (The Cost Form ... through the bias paragraph after the proof of
% Proposition att_identification). Readability rewrite only: no
% mathematical claim added, dropped, weakened or strengthened.
% Lines starting with "% [CHECK]" flag gaps found in the ORIGINAL
% proofs; they are kept as in the original and logged in
% proofs_A_ledger.md with a proposed fix.
% ============================================================

\subsection*{The Cost Form and Why the Normalisation Binds}

Valuations enter prices only through differences across listings. Any amount a buyer subtracts from every listing alike therefore shifts all her valuations by the same constant and drops out of every price gap. The bonus form of Section~\ref{ss:model_setup} uses this normalisation twice:
\begin{itemize}
    \item it folds the energy cost of an inefficient dwelling, $C_0$, into the intercept $\beta_0$;
    \item it writes the insensitive buyer's valuation as $v_0=\delta_0+\sum_n\delta_nX_n$, which in cost form amounts to subtracting the same constant $C^{\rm ins}$ from every listing.
\end{itemize}
Both choices are innocuous in Section~\ref{ss:model_setup}, because each buyer treats every listing the same way: a sensitive buyer always prices energy, and an insensitive one never does.

The normalisation stops being innocuous as soon as the same buyer may price energy on one listing and not on another. A buyer who ignores energy on a silent listing but prices it on a rated one subtracts different amounts from the two. The difference involves the \emph{level} of what she subtracts when she ignores energy, not only $\beta_1$. Formally, under Assumption~\ref{assum:attention}, the non-disclosure premium departs from the prediction of Section~\ref{sec:equilibrium}, $e$ times the energy-efficiency premium $\rho\omega_1\beta_1$, by
\begin{equation*}
    \Delta-e\,\rho\,\omega_1\beta_1\;=\;\rho\,(\omega_1-\omega_0)\,\big(\bar C^{\rm sil}-C^{\rm def}\big).
\end{equation*}
To see this, combine~\eqref{eq:premium_att} with $e\,\beta_1=C_0-\bar C^{\rm sil}$, which is~\eqref{eq:csil}, and with $\omega_0=\psi\omega_1$:
\begin{align*}
    \Delta-e\,\rho\,\omega_1\beta_1
    &\;=\;\rho\,\omega_1\big(C_0-C^*\big)-\rho\,\omega_1\big(C_0-\bar C^{\rm sil}\big)
    \;=\;\rho\,\omega_1\big(\bar C^{\rm sil}-C^*\big)\\
    &\;=\;\rho\,\omega_1\,(1-\psi)\big(\bar C^{\rm sil}-C^{\rm def}\big)
    \;=\;\rho\,(\omega_1-\omega_0)\,\big(\bar C^{\rm sil}-C^{\rm def}\big).
\end{align*}
This departure depends on the level of the silent pool's energy cost and on the default. It vanishes when rated and silent listings receive the same attention, whatever the default. The cost form is the representation in which this object exists, and it has three further properties.
\begin{itemize}
    \item It nests the bonus form exactly (Lemma~\ref{lem:cost_form}), so the extension can only add to the model of Section~\ref{sec:equilibrium}.
    \item Its primitives are physical rather than behavioural. $C_\theta$ is what a dwelling of quality $\theta$ costs to heat, while attention and the default describe what buyers do with that cost. This is what lets the registry discipline the extension through $\eta$.
    \item It turns an implicit normalisation of the bonus form into an assumption that can be stated and tested (Proposition~\ref{prop:att_incentives}(iii)).
\end{itemize}

\subsection*{Proof of Lemma~\ref{lem:cost_form} (Bonus and Cost Forms)}

\begin{proof}\upshape
\emph{Idea.} The gross value $\bar V$ adds back exactly the energy cost $C_0$ that the bonus form nets out of its intercept, so subtracting each buyer's energy cost from $\bar V$ returns her bonus-form valuation. Throughout, $\bar V=\beta_0+\sum_n\beta_nX_n+C_0$ by~\eqref{eq:gross_value} and $\beta_1=C_0-C_1$.

\emph{Step 1 (disclosed rating).} For $\theta\in\{0,1\}$, $C_0-C_\theta=\beta_1\theta$. Hence
\begin{equation*}
    \bar V-C_\theta\;=\;\beta_0+\sum_n\beta_nX_n+(C_0-C_\theta)\;=\;\beta_0+\beta_1\theta+\sum_n\beta_nX_n,
\end{equation*}
the valuation of a disclosed rating in Section~\ref{ss:model_setup}.

\emph{Step 2 (silent listing).} By~\eqref{eq:csil}, $C_0-\bar C^{\rm sil}=e\,(C_0-C_1)=\beta_1e$. Hence
\begin{equation*}
    \bar V-\bar C^{\rm sil}\;=\;\beta_0+\beta_1e+\sum_n\beta_nX_n,
\end{equation*}
the valuation of a silent listing in Section~\ref{ss:model_setup}.

\emph{Step 3 (insensitive buyer).} $\bar V-C^{\rm ins}=v_0(X)$ by the definition of $C^{\rm ins}$.

\emph{Step 4 (uniqueness).} Suppose instead $\bar V=\beta_0+\sum_n\beta_nX_n+K$ for some constant $K$. Agreement with the bonus form requires $K-C_\theta=\beta_1\theta$ for $\theta\in\{0,1\}$. The case $\theta=0$ gives $K=C_0$. The case $\theta=1$ is then satisfied, because $C_0-C_1=\beta_1$.
\end{proof}

\subsection*{Proof of Lemma~\ref{lem:att_prices} (Prices and the Perceived Cost of Silence)}

\begin{proof}\upshape
\emph{Idea.} A price is the buyer-share-weighted average of valuations. The gross value and the insensitive buyer's cost are common to all listings and cancel from every gap, so only what sensitive buyers subtract survives. Shorthand: $C^*=\psi\,\bar C^{\rm sil}+(1-\psi)\,C^{\rm def}$, which is the $C^*$ of~\eqref{eq:premium_att}, and $\omega_0=\psi\omega_1$.

\emph{Step 1 (weights).} By Assumption~\ref{assum:attention} and Table~\ref{tab:att_buyers}, a buyer values a listing at $\bar V$ minus the energy cost she subtracts. The three buyer types arrive with probabilities $1-\rho$, $\rho\omega_d$ and $\rho(1-\omega_d)$, which sum to one. Under the pricing rule used to derive~\eqref{eq:premium_derived}, a listing's price is the probability-weighted average of these valuations.

\emph{Step 2 (the two prices).} Expand term by term and collect the $\bar V$ terms, whose weights sum to one by Step~1. A listing displaying a rating with energy cost $C$ ($\omega_d=\omega_1$) sells at
\begin{equation*}
    p\;=\;\bar V-(1-\rho)\,C^{\rm ins}-\rho\,\big[\omega_1\,C+(1-\omega_1)\,C^{\rm def}\big].
\end{equation*}
A silent listing ($\omega_d=\omega_0$, attentive buyers subtract $\bar C^{\rm sil}$) sells at
\begin{equation*}
    p_{\rm sil}\;=\;\bar V-(1-\rho)\,C^{\rm ins}-\rho\,\big[\omega_0\,\bar C^{\rm sil}+(1-\omega_0)\,C^{\rm def}\big].
\end{equation*}

\emph{Step 3 (price gap).} Subtract the two lines of Step~2. The common term $\bar V-(1-\rho)C^{\rm ins}$ cancels, and the default terms collect to $\rho(\omega_1-\omega_0)C^{\rm def}$. Using $\omega_0=\psi\omega_1$ and the definition of $C^*$,
\begin{align*}
    p-p_{\rm sil}
    &\;=\;-\rho\,\omega_1C+\rho\,\omega_0\,\bar C^{\rm sil}+\rho\,(\omega_1-\omega_0)\,C^{\rm def}\\
    &\;=\;-\rho\,\omega_1\big[C-\psi\,\bar C^{\rm sil}-(1-\psi)\,C^{\rm def}\big]\;=\;-\rho\,\omega_1\,(C-C^*).
\end{align*}
This is the last statement of the lemma.

\emph{Step 4 (the two premia).} Apply Step~3 with $C=C_1$ and with $C=C_0$, and subtract. $C^*$ cancels, and
\begin{equation*}
    b\;=\;p_{\rm H}-p_{\rm L}\;=\;\rho\,\omega_1\,(C_0-C_1)\;=\;\rho\,\omega_1\beta_1,
\end{equation*}
which is~\eqref{eq:b_att}. Applying Step~3 with $C=C_0$ alone gives $\Delta=p_{\rm sil}-p_{\rm L}=\rho\,\omega_1\,(C_0-C^*)$, which is~\eqref{eq:premium_att}.

\emph{Step 5 (nesting).} With $\omega_1=\omega_0=1$ and $C^{\rm ins}=C_0$, the two prices of Step~2 reduce to $p_1^*$ and $p_0^*$ of~\eqref{eq:premium_derived}.
\end{proof}

\subsection*{Proof of Proposition~\ref{prop:att_incentives} (Disclosure Incentives under Limited Attention)}

\begin{proof}\upshape
\emph{Idea.} Each threshold compares the gross gain from silence in~\eqref{eq:gains} with the hiding cost; the attention threshold rewrites $\tilde e>1$ on the cost scale. Throughout, $b>0$, as the definition $\tilde e=\Delta/b$ presupposes.

\emph{Step 1 (low-quality seller, (i)).} An informed seller with $\theta=0$ gains $g_0=b\,\tilde e$ by hiding, by~\eqref{eq:gains}, and pays $c_0$. She strictly prefers non-disclosure iff $b\,\tilde e>c_0$, that is, iff $\tilde e>c_0/b$.

\emph{Step 2 (high-quality seller, (ii)).} An informed seller with $\theta=1$ gains $g_1=b\,(\tilde e-1)$ and pays $c_0+c_1$. She strictly prefers non-disclosure iff
\begin{equation*}
    b\,(\tilde e-1)\;>\;c_0+c_1\quad\Longleftrightarrow\quad\tilde e\;>\;1+\frac{c_0+c_1}{b}.
\end{equation*}
Since $c_0,c_1\ge0$, this requires $\tilde e>1$, that is, $\Delta=b\,\tilde e>b$: silence sells above a verified efficient rating.

\emph{Step 3 (attention threshold, (iii)).} By~\eqref{eq:b_att}--\eqref{eq:premium_att}, $\tilde e=(C_0-C^*)/(C_0-C_1)$, so $\tilde e>1$ iff $C^*<C_1$. Substitute $C^*=C^{\rm def}+\psi(\bar C^{\rm sil}-C^{\rm def})$ and divide by $\bar C^{\rm sil}-C^{\rm def}>0$:
\begin{equation*}
    C^*<C_1\quad\Longleftrightarrow\quad\psi\;<\;\frac{C_1-C^{\rm def}}{\bar C^{\rm sil}-C^{\rm def}}.
\end{equation*}

\emph{Step 4 (the default must be optimistic, (iii)).} $C^*$ is a convex combination of $C^{\rm def}$ and $\bar C^{\rm sil}$, and $\bar C^{\rm sil}\ge C_1$ by~\eqref{eq:csil}. If $C^{\rm def}\ge C_1$, both points are at least $C_1$, and so is $C^*$. Hence $C^*<C_1$ requires $C^{\rm def}<C_1$.

\emph{Step 5 (full attention, (iv)).} With $\psi=1$,~\eqref{eq:etilde} gives $\tilde e=e$. Steps~1--2 are then the comparisons behind Proposition~\ref{prop:equilibrium}, with $b=\rho\omega_1\beta_1$ in the role of $\rho\beta_1$. Without hiding costs they are those of Proposition~\ref{prop:premium}.
\end{proof}

\subsection*{Proof of Corollary~\ref{cor:att_hq} (High-Quality Sellers Disclose outside Complete Adverse Selection)}

\begin{proof}\upshape
\emph{Idea.} By~\eqref{eq:gains}, silence always pays a high-quality seller $b$ less than a low-quality one, while hiding costs her at least as much. So she can want to hide only once every low-quality seller does.

\emph{Step 1 (if $\sigma_0^*<1$, then $\sigma_1^*=0$).} If $\sigma_0^*<1$, some informed low-quality seller weakly prefers disclosure, so $g_0\le c_0$. By~\eqref{eq:gains}, $g_1=g_0-b$, hence
\begin{equation*}
    g_1\;=\;g_0-b\;<\;c_0\;\le\;c_0+c_1.
\end{equation*}
Every informed high-quality seller strictly prefers disclosure, so $\sigma_1^*=0$.

\emph{Step 2 (Complete Adverse Selection).} If $\sigma_0^*=1$, a high-quality seller hides only if $g_1=b[\tilde e(1,\sigma_1)-1]>c_0+c_1$ at the prevailing $\sigma_1$. As in Step~2 of the proof of Proposition~\ref{prop:att_incentives}, this requires $\tilde e>1$.
% [CHECK] A seller who mixes (0<sigma_1^*<1), or who hides at sigma_1^*=1 when \bar g_1(1)=c_0+c_1,
% [CHECK] is only indifferent: hiding requires g_1 >= c_0+c_1, i.e. \tilde e(1,\sigma_1) >= 1+(c_0+c_1)/b
% [CHECK] (weak), which gives \tilde e >= 1, strictly > 1 only if c_0+c_1>0. Proposed fix: ">=" here and
% [CHECK] in the statement of Corollary cor:att_hq. See ledger, ISSUE 1.
\end{proof}

\subsection*{Proof of Proposition~\ref{prop:att_equilibrium} (Equilibrium under Limited Attention)}

\begin{proof}\upshape
\emph{Idea.} When $\sigma_1=0$, the low type's net gain $\Phi_\psi(\sigma_0)$ is continuous and weakly decreasing, so the regimes follow from its sign at $\sigma_0=0$ and $\sigma_0=1$, as in Proposition~\ref{prop:equilibrium}. What is new is that, once every low type hides, high-quality hiding improves the pool and reinforces itself, so equilibria can be multiple. Objects used: $\Phi_\psi$ in~\eqref{eq:Phi_psi}, $\Delta_\psi$ in~\eqref{eq:Delta_psi}, the high type's gain $\bar g_1(\sigma_1)=b\,[\tilde e(1,\sigma_1)-1]$, and $e(1,1)=\bar\theta$.

\emph{Step 1 (preliminaries).} Whenever $\sigma_0^*<1$, $\sigma_1^*=0$ by Corollary~\ref{cor:att_hq}. $\Phi_\psi$ is continuous on $[0,1]$. By Lemma~\ref{lem:etilde_monotone}, it is strictly decreasing if $\psi>0$ and constant if $\psi=0$.

\emph{Step 2 (Cost-Deterred Hiding, $c_0>\Delta_\psi(\bar\theta)$).} This follows as in Proposition~\ref{prop:equilibrium}. By Step~1, $\Phi_\psi(\sigma_0)\le\Phi_\psi(0)=\Delta_\psi(\bar\theta)-c_0<0$ for every $\sigma_0$, so $\sigma_0^*=0$ and $e^*=\bar\theta$.
% [CHECK] Gap in the original: "as in Proposition prop:equilibrium" does not rule out equilibria with
% [CHECK] sigma_1^*>0, which cannot arise in the core model. Proposed fix: Step 6 below uses only
% [CHECK] c_0 >= Delta_psi(e(1)), which holds here, so it applies verbatim (the proof of
% [CHECK] Proposition prop:att_seller_types already says so). See ledger, ISSUE 2.

\emph{Step 3 (Complete Adverse Selection: every low type hides).} Let $c_0<\Delta_\psi(e(1))$. Since $\Phi_\psi$ is decreasing,
\begin{equation*}
    \Phi_\psi(\sigma_0)\;\ge\;\Phi_\psi(1)\;=\;\Delta_\psi(e(1))-c_0\;>\;0\qquad\text{for every }\sigma_0 .
\end{equation*}
Hiding is strictly profitable for every low-quality seller, a fortiori when $\sigma_1>0$ because $\tilde e$ is weakly increasing in $\sigma_1$ (Lemma~\ref{lem:etilde_monotone}). So $\sigma_0^*=1$.

\emph{Step 4 (Complete Adverse Selection: high-quality best responses).} Given $\sigma_0=1$, a high-quality seller hides iff $\bar g_1(\sigma_1)>c_0+c_1$. By Lemma~\ref{lem:etilde_monotone}, $\bar g_1$ is continuous and increasing in $\sigma_1$: high-quality hiding is self-reinforcing because it improves the pool. Hence
\begin{itemize}
    \item $\sigma_1=0$ is a best response to itself iff $\bar g_1(0)\le c_0+c_1$;
    \item $\sigma_1=1$ is a best response to itself iff $\bar g_1(1)\ge c_0+c_1$.
\end{itemize}
In both cases $\sigma_0=1$ remains a best response, because by~\eqref{eq:gains} a low-quality seller's gain exceeds a high-quality seller's by $b$.

\emph{Step 5 (Complete Adverse Selection: existence and multiplicity).} One of the two conditions of Step~4, or an interior crossing of $\bar g_1$ with $c_0+c_1$, always holds, so an equilibrium exists. Uniqueness can fail. With $\pi=0.5$, $\bar\theta=0.6$, $e^{\rm def}=2.3$ and $\psi=0.5$,
\begin{equation*}
    \bar g_1(0)/b\;=\;0.364,\qquad\bar g_1(1)/b\;=\;0.450,
\end{equation*}
so any $(c_0+c_1)/b$ between the two supports both. If $\psi=1$, then $\tilde e(1,\sigma_1)=e(1,\sigma_1)<1$, so $\bar g_1<0$ and only $\sigma_1^*=0$ survives.

\emph{Step 6 (Partial Pooling: no high-quality hiding).} Let $\Delta_\psi(e(1))\le c_0\le\Delta_\psi(\bar\theta)$. An equilibrium with $\sigma_1^*>0$ would need $\sigma_0^*=1$ and $\bar g_1(\sigma_1)\ge c_0+c_1\ge\Delta_\psi(e(1))$. By Lemma~\ref{lem:etilde_monotone} and $e(1,1)=\bar\theta$, $\bar g_1(\sigma_1)\le\bar g_1(1)=\Delta_\psi(\bar\theta)-b$, so it would need $\Delta_\psi(\bar\theta)-\Delta_\psi(e(1))\ge b$. But, by~\eqref{eq:Delta_psi},
\begin{equation*}
    \Delta_\psi(\bar\theta)-\Delta_\psi(e(1))\;=\;\psi\,b\,\big[\bar\theta-e(1)\big]\;<\;b.
\end{equation*}

\emph{Step 7 (Partial Pooling with $\psi>0$).} By Step~6, $\sigma_1^*=0$. In the band, $\Phi_\psi(0)=\Delta_\psi(\bar\theta)-c_0\ge0\ge\Delta_\psi(e(1))-c_0=\Phi_\psi(1)$, so the intermediate value theorem gives a zero of $\Phi_\psi$. It is unique by strict monotonicity (Step~1). Solving $\Delta_\psi(e^*)=c_0$ gives
\begin{equation*}
    e^*\;=\;\frac{c_0/b-(1-\psi)\,e^{\rm def}}{\psi},
\end{equation*}
the first expression in~\eqref{eq:att_pp}. Inverting~\eqref{eq:belief_simplified} at $e^*$ gives $\sigma_0^*$.

\emph{Step 8 (Partial Pooling with $\psi=0$).} $\Phi_0$ is constant. If it is non-zero, one of the corners obtains. If it is zero, every $\sigma_0$ is a best response.
\end{proof}

\subsection*{Proof of Proposition~\ref{prop:att_seller_types} (Prices and Payoffs by Seller Type)}

\begin{proof}\upshape
\emph{Idea.} The silent price is pinned at a corner by the extreme pool and in Partial Pooling by the indifference of the marginal hider, which gives the truncation formula; payoffs follow from who bears a hiding cost. Normalise $p_{\rm L}=0$.

\emph{Step 1 (prices).} By Lemma~\ref{lem:att_prices}, a listing displaying an efficient rating sells at $b=\rho\omega_1\beta_1$ and a silent listing at $\Delta=b\,\tilde e$, with $\tilde e=\psi e+(1-\psi)e^{\rm def}$ evaluated at the equilibrium pool by~\eqref{eq:etilde}. Displayed prices depend on no belief, so they are the same in every regime and are defined even where no seller displays a given rating. With $\sigma_1^*=0$ the pool belief is $e^*=e(\sigma_0^*)$ of~\eqref{eq:belief_simplified}, so $\Delta^*=\Delta_\psi(e^*)$ by~\eqref{eq:Delta_psi}.

\emph{Step 2 (the band is ordered).} Since $b>0$ and $\psi\ge0$, $\Delta_\psi$ is weakly increasing in $e$. Since $e(1)<\bar\theta$, $\Delta_\psi(e(1))\le\Delta_\psi(\bar\theta)$.

\emph{Step 3 (the silent price).} By Proposition~\ref{prop:att_equilibrium}:
\begin{itemize}
    \item \emph{Cost-Deterred Hiding:} $e^*=\bar\theta$ and $c_0>\Delta_\psi(\bar\theta)$, so $\Delta^*=\Delta_\psi(\bar\theta)$;
    \item \emph{Partial Pooling:} $\Delta_\psi(e^*)=c_0$, with $c_0$ in the band (if $\psi=0$, $\Delta_0\equiv b\,e^{\rm def}=c_0$ at every $e$);
    \item \emph{Complete Adverse Selection} with $\sigma_1^*=0$: $e^*=e(1)$ and $c_0<\Delta_\psi(e(1))$, so $\Delta^*=\Delta_\psi(e(1))$.
\end{itemize}
In each case $\Delta^*$ is $c_0$ truncated to $[\Delta_\psi(e(1)),\Delta_\psi(\bar\theta)]$, an interval by Step~2. This is~\eqref{eq:att_premium_trunc}. In Partial Pooling, Definition~\ref{def:effective_belief} gives $\tilde e^*=\Delta^*/b=c_0/b$.

\emph{Step 4 (payoffs).} A seller's net payoff is the price of her listing less the hiding cost, which only an informed seller who withholds a rating bears.
\begin{itemize}
    \item An informed high-quality seller displays when $\sigma_1^*=0$ and earns $b$.
    \item An informed low-quality seller earns $0$ by displaying and $\Delta^*-c_0$ by hiding. In Cost-Deterred Hiding $\Delta^*-c_0<0$, so she displays. In Partial Pooling $\Delta^*-c_0=0$, so both actions earn zero. In Complete Adverse Selection $\Delta^*-c_0>0$, so she hides. Her equilibrium payoff is therefore $\max\{\Delta^*-c_0,0\}$.
    \item An uninformed seller cannot display (No Falsification), bears no hiding cost, and earns $\Delta^*$ whatever her quality.
\end{itemize}
In Partial Pooling the four payoffs are $b$, $0$, $0$ and $c_0$. This gives the two orderings in the statement, and $\Delta^*>b$ if and only if $c_0>b$.

\emph{Step 5 (silence above an efficient rating: $e^{\rm def}\le1$).} Since $b>0$, $\Delta^*>b$ if and only if $\tilde e^*>1$. Also $e(1)\le e^*\le\bar\theta<1$. If $e^{\rm def}\le1$, then $\tilde e^*$ is a weighted average of $e^*<1$ and $e^{\rm def}\le1$, so $\tilde e^*\le1$.

\emph{Step 6 (silence above an efficient rating: $e^{\rm def}>1$).} Then $e^{\rm def}>e^*$ and, for any $e\neq e^{\rm def}$,
\begin{equation*}
    \tilde e-1\;=\;\psi e+(1-\psi)e^{\rm def}-1\;=\;\big(e^{\rm def}-e\big)\big(\bar\psi(e)-\psi\big).
\end{equation*}
Since $e^{\rm def}-e^*>0$, $\tilde e^*>1$ if and only if $\psi<\bar\psi(e^*)$.

\emph{Step 7 (the threshold in cost units).} The quality scale $q$ maps $C_1$ to $1$, $C^{\rm def}$ to $e^{\rm def}$ and $\bar C^{\rm sil}$ to $e$. Hence $e^{\rm def}-1=(C_1-C^{\rm def})/(C_0-C_1)$ and $e^{\rm def}-e=(\bar C^{\rm sil}-C^{\rm def})/(C_0-C_1)$, and at the equilibrium pool
\begin{equation*}
    \bar\psi(e^*)\;=\;\frac{C_1-C^{\rm def}}{\bar C^{\rm sil}-C^{\rm def}},
\end{equation*}
which is the threshold of Proposition~\ref{prop:att_incentives}(iii).

\emph{Step 8 (the uninformed seller's rank).} When $\Delta^*>b$, the uninformed seller earns $\Delta^*>b$. She also earns $\Delta^*\ge\max\{\Delta^*-c_0,0\}$, since $c_0\ge0$ and $\Delta^*>0$.

\emph{Step 9 (high-quality sellers when $\sigma_1^*=0$).} An informed high-quality seller gains $g_1=\Delta^*-b$ from silence by~\eqref{eq:gains}. In Cost-Deterred Hiding and Partial Pooling $\Delta^*\le c_0$, so
\begin{equation*}
    g_1-(c_0+c_1)\;\le\;-(b+c_1)\;<\;0 .
\end{equation*}
She strictly prefers to display even when $g_1>0$. In Partial Pooling the slack is exactly $b+c_1$.

\emph{Step 10 (equilibria with $\sigma_1^*>0$).} By Corollary~\ref{cor:att_hq}, $\sigma_1^*>0$ requires $\sigma_0^*=1$ and $\Delta^*-b\ge c_0+c_1$. It cannot arise with $c_0\ge\Delta_\psi(e(1))$: that would give $\Delta_\psi(\bar\theta)-\Delta_\psi(e(1))\ge b$, whereas the difference is $\psi\,b\,[\bar\theta-e(1)]<b$ (proof of Proposition~\ref{prop:att_equilibrium}). So it arises only in Complete Adverse Selection, where $\Delta^*\ge b+c_0+c_1$. At $\sigma_1^*=1$ every seller is silent, $e(1,1)=\bar\theta$ and $\Delta^*=\Delta_\psi(\bar\theta)$. At an interior $\sigma_1^*$ the high-quality seller is indifferent and $\Delta^*=b+c_0+c_1$.
\end{proof}

\subsection*{Proof of Corollary~\ref{cor:att_ordering} (Where Silence Sits among Ratings)}

\begin{proof}\upshape
\emph{Idea.} The sign of $\Delta_\psi(e)-b$ is that of $\bar\psi(e)-\psi$, and $\bar\psi$ is increasing; since every equilibrium silent price lies between its values at $e(1)$ and $\bar\theta$, comparing $\psi$ with $\bar\psi$ at the two ends classifies every equilibrium.

\emph{Step 1 (two facts when $e^{\rm def}>1$).} Let $e^{\rm def}>1$. By the identity in the proof of Proposition~\ref{prop:att_seller_types},
\begin{equation*}
    \Delta_\psi(e)-b\;=\;b\,\big(e^{\rm def}-e\big)\big(\bar\psi(e)-\psi\big),\qquad
    \bar\psi'(e)\;=\;\frac{e^{\rm def}-1}{(e^{\rm def}-e)^2}\;>\;0,
\end{equation*}
so $\bar\psi(e(1))<\bar\psi(\bar\theta)$. Any equilibrium with $\sigma_1^*=0$ has $\Delta^*\in[\Delta_\psi(e(1)),\Delta_\psi(\bar\theta)]$ by~\eqref{eq:att_premium_trunc}.

\emph{Step 2 (the three cases when $e^{\rm def}>1$).}
\begin{itemize}
    \item If $\psi<\bar\psi(e(1))$, then $\Delta^*\ge\Delta_\psi(e(1))>b$. Any equilibrium with $\sigma_1^*>0$ has $\Delta^*\ge b+c_0+c_1>b$.
    \item If $\psi\ge\bar\psi(\bar\theta)$, then $\Delta^*\le\Delta_\psi(\bar\theta)\le b$. Moreover $\bar g_1(\sigma_1)\le\bar g_1(1)=\Delta_\psi(\bar\theta)-b\le0<c_0+c_1$ by Lemma~\ref{lem:etilde_monotone} and $e(1,1)=\bar\theta$, so no equilibrium has $\sigma_1^*>0$.
    \item Otherwise $\Delta_\psi(e(1))\le b<\Delta_\psi(\bar\theta)$, so $b$ lies inside the band. By~\eqref{eq:att_premium_trunc}, $\Delta^*>b$ if and only if $c_0>b$.
\end{itemize}

\emph{Step 3 (the case $e^{\rm def}\le1$).} Then $\tilde e\le1$ at every pool, and the second case applies by the same argument.

\emph{Step 4 (the model of Section~\ref{sec:equilibrium}).} With $\psi=1$, $\bar\psi(\bar\theta)<1=\psi$. And $0<e(1)\le e^*\le\bar\theta<1$ gives $0<\Delta^*<b$.
\end{proof}

\subsection*{Proof of Proposition~\ref{prop:att_cs} (Comparative Statics under Limited Attention)}

\begin{proof}\upshape
\emph{Idea.} In Partial Pooling the indifference condition pins $e^*$ in closed form by~\eqref{eq:att_pp}, and $\sigma_0^*$ is a strictly decreasing function of $e^*$. Each sign in~(i) therefore follows from differentiating $e^*$. The corner premia are linear in $\psi$.

\emph{Step 1 (attention, $\psi$).} From~\eqref{eq:att_pp},
\begin{equation*}
    \frac{\partial e^*}{\partial\psi}\;=\;\frac{e^{\rm def}-c_0/b}{\psi^2}.
\end{equation*}

\emph{Step 2 (sign of Step 1).} In Partial Pooling $c_0/b=\tilde e^*=\psi e^*+(1-\psi)e^{\rm def}$, so $e^{\rm def}-c_0/b=\psi\,(e^{\rm def}-e^*)$. This is positive whenever $e^*<e^{\rm def}$, which holds when $C^{\rm def}\le\bar C^{\rm sil}$, since then $e^{\rm def}\ge q(\bar C^{\rm sil})=e^*$. So $\partial e^*/\partial\psi>0$.
% [CHECK] The weak inequalities only give e^def >= e^*, i.e. de^*/dpsi >= 0. The proposition assumes
% [CHECK] C^def < Csil (strict), which gives e^def > e^* and the strict sign. Proposed fix: replace
% [CHECK] "\le" and "\ge" by "<" and ">" in this sentence. See ledger, ISSUE 3.

\emph{Step 3 (the other parameters).} Differentiating~\eqref{eq:att_pp},
\begin{equation*}
    \frac{\partial e^*}{\partial c_0}\;=\;\frac{1}{\psi\,b}\;>\;0,\qquad
    \frac{\partial e^*}{\partial C^{\rm def}}\;=\;\frac{1-\psi}{\psi\,(C_0-C_1)}\;\ge\;0,
\end{equation*}
and $\partial e^*/\partial(\rho\omega_1)<0$ through $b=\rho\omega_1\beta_1$.
% [CHECK] de^*/dC^def = 0 at psi = 1, so dsigma_0^*/dC^def = 0 there, while the statement claims a strict
% [CHECK] "< 0". Proposed fix: "dsigma_0^*/dC^def <= 0, strictly if psi < 1" in Proposition prop:att_cs(i).
% [CHECK] See ledger, ISSUE 4.

\emph{Step 4 (hiding rate).} By~\eqref{eq:att_pp}, $\sigma_0^*$ is strictly decreasing in $e^*$, so Steps~1--3 give the signs in~(i) for $\psi$, $C^{\rm def}$, $c_0$ and $\rho\omega_1$. The effect of $\pi$ operates only through the inversion of~\eqref{eq:belief_simplified}, exactly as in Proposition~\ref{prop:comparative_statics}(ii).

\emph{Step 5 (corner premia and attention, (ii)).} The Cost-Deterred premium $b[\psi\bar\theta+(1-\psi)e^{\rm def}]$ has derivative $b[\bar\theta-e^{\rm def}]$ in $\psi$. It is negative iff
\begin{equation*}
    C^{\rm def}\;<\;C_0-\bar\theta\,(C_0-C_1)\;=\;\bar\theta C_1+(1-\bar\theta)C_0 .
\end{equation*}
The Complete Adverse Selection premium has derivative $b[e(1)-e^{\rm def}]<b[\bar\theta-e^{\rm def}]$, since $e(1)<\bar\theta$. The same condition therefore makes it decreasing in $\psi$.

\emph{Step 6 (Complete Adverse Selection premium and $\pi$, (ii)).} That premium has derivative $\psi\,b\,\partial e(1)/\partial\pi$ in $\pi$, with
\begin{equation*}
    \frac{\partial e(1)}{\partial\pi}\;=\;-\frac{\bar\theta\,(1-\bar\theta)}{(1-\pi\bar\theta)^2}\;<\;0 .
\end{equation*}
% [CHECK] psi*b*de(1)/dpi is < 0 only for psi > 0 (it is 0 at psi = 0). Proposed fix: add "for psi > 0"
% [CHECK] in Proposition prop:att_cs(ii). See ledger, ISSUE 5.
\end{proof}

\subsection*{Proof of Proposition~\ref{prop:att_identification} (Identification under Limited Attention)}

\begin{proof}\upshape
\emph{Idea.} Under Assumption~\ref{assum:no_pre_pricing}, the reform-induced changes in the two premia are the post-reform premia of Lemma~\ref{lem:att_prices}. Their ratio cancels $\rho\omega_1$ and leaves the perceived cost of silence on the quality scale. Shorthand: $\delta x\equiv x^{\rm post}-x^{\rm pre}$, $\psi_\delta\equiv\delta(\rho\omega_0)/\delta(\rho\omega_1)$ and $C^*_\delta\equiv C^{\rm def}+\psi_\delta(\bar C^{\rm sil}-C^{\rm def})$.

\emph{Step 1 (attention to displayed ratings, (i)).} With $\beta_1^{\rm anchor}=C_0-C_1=\beta_1$,~\eqref{eq:b_att} in differences gives
\begin{equation*}
    \delta b\;=\;\delta(\rho\omega_1)\,\beta_1,
\end{equation*}
which $\hat\beta_1$ estimates. Dividing by the anchor gives $\delta(\rho\omega_1)$.
% [CHECK] The clause "with omega_1 = 1 it is the rho_m of Section sec:identification" is not argued in the
% [CHECK] original. One-line fix: with omega_1 = 1 and Assumption assum:no_pre_pricing, delta(rho omega_1)
% [CHECK] = rho^post, the post-reform rho_m measured there. See ledger, ISSUE 6 (minor).

\emph{Step 2 (the non-disclosure premium in differences).} Under Assumption~\ref{assum:no_pre_pricing} the pre-reform energy terms vanish. By~\eqref{eq:premium_att}, in differences,
\begin{equation*}
    \delta\Delta\;=\;\delta(\rho\omega_1)\,(C_0-C^{\rm def})-\delta(\rho\omega_0)\,(\bar C^{\rm sil}-C^{\rm def}),
\end{equation*}
where $\bar C^{\rm sil}$ is evaluated at the post-reform pool.

\emph{Step 3 (the perceived cost of silence, (ii)).} Dividing Step~2 by Step~1,
\begin{equation*}
    \frac{\delta\Delta}{\delta b}\;=\;e^{\rm def}-\psi_\delta\,(e^{\rm def}-e)\;=\;\frac{C_0-C^*_\delta}{C_0-C_1}.
\end{equation*}
Since $C_0$ and $C_1$ are known, the ratio pins down $C^*_\delta$.

\emph{Step 4 (point identification under $C^{\rm def}=0$, (iii)).} Subtract $e$ from Step~3 and multiply by $\delta b$:
\begin{equation*}
    \delta\Delta-e\,\delta b\;=\;(1-\psi_\delta)\,\delta b\,\big(e^{\rm def}-e\big).
\end{equation*}
Setting $C^{\rm def}=0$ (so $e^{\rm def}=\eta$) and replacing $\delta b$ and $\delta\Delta$ by their estimates gives~\eqref{eq:att_lambda}.

\emph{Step 5 (the identified set, (iii)).} The identified set follows because $C^*_\delta$ is a convex combination of $C^{\rm def}\ge0$ and $\bar C^{\rm sil}$. In particular $\psi_\delta\bar C^{\rm sil}=C^*_\delta-(1-\psi_\delta)C^{\rm def}\le C^*_\delta$.
% [CHECK] The bound C^def <= C^*_delta holds only if C^def <= Csil (equivalently C^*_delta <= Csil when
% [CHECK] psi_delta < 1): with Csil = 1, C^def = 3, psi_delta = 0.5, C^*_delta = 2 < C^def. The bound on
% [CHECK] psi_delta holds unconditionally. Proposed fix: add "if C^*_delta <= Csil" (or assume
% [CHECK] C^def <= Csil, as in Propositions prop:att_incentives(iii) and prop:att_cs) to part (iii).
% [CHECK] See ledger, ISSUE 7.

\emph{Step 6 (the hiding cost, (iv)).} In Partial Pooling the indifference of the informed low-quality seller gives $\Delta=c_0$ (Proposition~\ref{prop:att_equilibrium}). Under Assumption~\ref{assum:no_pre_pricing} the post-reform premium equals $\delta\Delta$.
\end{proof}

Suppose instead that buyers already priced energy before the reform, contrary to Assumption~\ref{assum:no_pre_pricing}, with pre-reform attention to silent listings $\omega_0^{\rm pre}>0$, and that the pool belief moved from $e^{\rm pre}$ to $e^{\rm post}$. Then the estimator of $1-\psi_\delta$ is biased by
\begin{equation*}
    \frac{\rho\omega_0^{\rm pre}}{\delta(\rho\omega_1)}\cdot\frac{e^{\rm post}-e^{\rm pre}}{\eta-e^{\rm post}}.
\end{equation*}
The bias is zero under Assumption~\ref{assum:no_pre_pricing}, and small whenever pre-reform attention to silence was low.

% ============================================================
% REWRITE (readability only) of paper-real-estate/sections/A-attention.tex, lines 107--177:
% "Exact Logit Pricing", "The Ratio-of-Sums Estimator", and "Proofs for Section sec:policy"
% up to the end of the proof of Proposition prop:att_saving. No mathematical claim changed.
% Comments starting "% [CHECK]" mark gaps or errors found in the ORIGINAL; they are kept
% as stated and a fix is proposed in proofs_B_ledger.md.
% ============================================================

\subsection*{Exact Logit Pricing}
\label{app:att_exact}

All propositions of Section~\ref{ss:attention} are derived under the linear pricing approximation of~\eqref{eq:premium_derived}. Under the exact demand~\eqref{eq:demand_att}, buyers form three segments and the seller's optimal listing price is not a function of the mean sensitive valuation alone, so the propositions do not carry over automatically. For the two-segment demand~\eqref{eq:demand}, the optimal listing price is nondecreasing in the sensitive buyers' valuation whenever the cross-partial condition of the monotone-comparative-statics argument holds. We conjecture, but have not shown, that the orderings of Propositions~\ref{prop:att_incentives}--\ref{prop:att_cs} survive exact pricing at the estimated logit scale. With attention \emph{shares} rather than a representative weight, a seller may also price to the inattentive segment and accept a lower probability of sale, which predicts longer time on market for silent listings at higher prices.

\subsection*{The Ratio-of-Sums Estimator}
\label{app:att_estimator}

\noindent\emph{Idea.} The two estimated gaps share the sampling error of the E-to-G index, so a regression of one on the other is biased; summing before dividing averages the noise out, unless the sample is selected on the noisy regressor. \emph{Notation.} $\hat\beta_{1,m}$ and $\hat\Delta_m$ estimate $\delta b_m$ and $\delta\Delta_m$; $\mathbb E_m$ is the expectation across markets.

\emph{Step 1 (Error structure).} Let $\varepsilon^{\rm H}_m$, $\varepsilon^{\rm L}_m$ and $\varepsilon^{\rm sil}_m$ be the sampling errors in the reform-induced change of the A-to-D, E-to-G and silent price indices of market $m$. They are independent across groups and markets, with variances $s_H^2$, $s_L^2$ and $s_{\rm sil}^2$. Then
\begin{equation*}
    \hat\beta_{1,m}=\delta b_m+\xi^{\beta}_m,\qquad \hat\Delta_m=\delta\Delta_m+\xi^{\Delta}_m,\qquad \xi^{\beta}_m=\varepsilon^{\rm H}_m-\varepsilon^{\rm L}_m,\qquad \xi^{\Delta}_m=\varepsilon^{\rm sil}_m-\varepsilon^{\rm L}_m.
\end{equation*}
Both errors therefore contain the E-to-G error $\varepsilon^{\rm L}_m$.

\emph{Step 2 (Regression through the origin).} Regress $\hat\lambda_m$ through the origin on $x_m\equiv\hat\beta_{1,m}(\eta_m-\hat e^*_m)$. By Step~1, $\mathbb E_m(\xi^\beta)^2=s_H^2+s_L^2$ and $\mathbb E_m\xi^\beta\xi^\Delta=s_L^2$, so the probability limit of $1-\hat\psi_\delta$ is
\begin{equation*}
    (1-\psi_\delta)\,\frac{\mathbb E_m(\delta b)^2}{\mathbb E_m(\delta b)^2+s_H^2+s_L^2}\;+\;\frac{s_L^2-e\,(s_H^2+s_L^2)}{(\eta-e)\,\big(\mathbb E_m(\delta b)^2+s_H^2+s_L^2\big)}.
\end{equation*}
% [CHECK] The formula treats e and eta as common across markets and \hat e^*_m as measured without error; neither is stated. See ledger, ISSUE R1.
The first term attenuates $1-\psi_\delta$ toward zero. The second is negative whenever $e>s_L^2/(s_H^2+s_L^2)$. Both push the estimate toward full attention.

\emph{Step 3 (Ratio of sums, no selection).} The ratio-of-sums estimator is $1-\hat\psi_\delta=\sum_m\hat\lambda_m/\sum_mx_m$. Absent selection it is consistent, because the noise in each sum has mean zero.

\emph{Step 4 (Selection).} The estimation sample is, however, selected on $\hat\beta_{1,m}>0$. Since $\hat\beta_{1,m}$ and $\hat\Delta_m$ share the E-to-G error (Step~1), selection gives the retained errors non-zero means. This biases both estimators toward full attention.
% [CHECK] The direction of the selection bias is asserted, not derived; analytically it depends on (s_H, s_L, e) and the distribution of delta b. See ledger, ISSUE R2.

\emph{Step 5 (Monte Carlo).} The Monte Carlo is calibrated to the \attMCmarkets{} volume-filtered apartment markets: the mean and variance of $\hat\beta_{1,m}$ before selection, with sampling noise accounting for \attMCshareLo{} to \attMCshareHi{} percent of its variance, and the aggregate $e$ and $\eta$ of Table~\ref{tab:attention}. The true $\psi_\delta$ is \attMCtrue{}. Mean estimates are:
\begin{itemize}
    \item \emph{without selection,} between \attMCromLo{} and \attMCromHi{} for the ratio of sums and between \attMColsLo{} and \attMColsHi{} for the regression through the origin;
    \item \emph{with selection on $\hat\beta_{1,m}>0$,} \attMCromSelLo{} to \attMCromSelHi{} for the ratio of sums and \attMColsSelLo{} to \attMColsSelHi{} for the regression.
\end{itemize}
The apartment estimate is therefore, if anything, an upper bound on attention to silence.

\emph{Step 6 (What the aggregation fixes by construction).} Under the ratio-of-sums aggregation, the point estimate equals the upper end of the identified interval at $C^{\rm def}=0$, and the fitted premium equals the mean premium. Both are identities of the aggregation. The informative check is market by market (Section~\ref{ss:att_equilibrium}).

\subsection*{Proofs for Section~\ref{sec:policy}}

\noindent\emph{Notation used throughout.} Payoffs are measured relative to a disclosed low rating, and $c_1=0$. By Lemma~\ref{lem:att_prices}, a disclosed low rating, a disclosed high rating and a silent listing sell for
\begin{equation*}
    0,\qquad b=\rho\omega_1\beta_1,\qquad b\tilde e.
\end{equation*}
The net gain from hiding is $h=b\tilde e-c_0$. By~\eqref{eq:gains}, the gains from hiding of the low and the high type are
\begin{equation*}
    g_0=b\tilde e,\qquad g_1=b(\tilde e-1),\qquad\text{so}\quad h=g_0-c_0\quad\text{and}\quad h-b=g_1-c_0.
\end{equation*}
$\Delta_\psi(e)$ is the hiding gain~\eqref{eq:Delta_psi}, $e(1)=(1-\pi)\bar\theta/(1-\pi\bar\theta)$ the pool belief when every informed low-quality seller hides, and $e^*$, $\sigma_0^*$ the Partial Pooling belief and hiding rate~\eqref{eq:att_pp}.

\subsubsection*{Proof of Lemma~\ref{lem:renov_behaviour}}

\begin{proof}\upshape
\emph{Idea.} A renovated seller faces exactly the choice of an informed high-quality seller, so she hides only when that seller would. \emph{Notation.} $h=b\tilde e-c_0$, $g_1=b(\tilde e-1)$; $\hat R-S$ is the net renovation cost.

\emph{Step 1 (Strategy~4 against Strategy~3).} Strategy~3 pays $b-(\hat R-S)$ and Strategy~4 pays $h-(\hat R-S)$. The renovation cost cancels, so Strategy~4 is weakly dominated if and only if
\begin{equation*}
    h\le b\quad\iff\quad b\tilde e-c_0\le b\quad\iff\quad g_1\le c_0.
\end{equation*}
This is the condition under which an informed high-quality seller discloses.

\emph{Step 2 (The condition by regime).}
\begin{itemize}
    \item \emph{Cost-Deterred Hiding and Partial Pooling.} The condition holds by Corollary~\ref{cor:att_hq}.
    % [CHECK] Corollary cor:att_hq assumes sigma_0^*<1; the Partial Pooling boundary c_0=Delta_psi(e(1)) has sigma_0^*=1. Its argument (g_0<=c_0) still applies. See ledger, ISSUE L1.
    \item \emph{Complete Adverse Selection with $\sigma_1^*=0$.} The condition is $\bar g_1(0)\le c_0+c_1$ of Proposition~\ref{prop:att_equilibrium}(ii), which holds because $\sigma_1^*=0$ is an equilibrium.
    \item \emph{Complete Adverse Selection with $\sigma_1^*>0$.} The same proposition requires $g_1\ge c_0$. By Step~1, $h\ge b$, so Strategy~4 is weakly preferred.
    % [CHECK] Proposition att_equilibrium(ii) states the sigma_1^*=1 case; an interior sigma_1^* gives g_1=c_0, hence h=b (indifference), which the statement reads as "also hides". See ledger, ISSUE L2.
\end{itemize}
\end{proof}

\subsubsection*{Proof of Proposition~\ref{prop:threshold} (Renovation Threshold under Limited Attention)}

\begin{proof}\upshape
\emph{Idea.} Renovation buys the rating premium $b$ and forfeits the rent from hiding; that rent is zero unless every informed low-quality seller hides, so $\bar R=b$ outside Complete Adverse Selection. \emph{Notation.} $h=g_0-c_0$, $\tilde e(1,0)=\psi e(1)+(1-\psi)e^{\rm def}$, $S^*=\max\{\hat R-\bar R,0\}$.

\emph{Step 1 (The threshold formula).} A renovating seller chooses between Strategies~3 and~4 and earns $\max\{b,h\}-(\hat R-S)$. A seller who does not renovate earns $\max\{0,h\}$. The threshold is therefore
\begin{equation*}
    \bar R=\max\{b,h\}-\max\{0,h\}=b-\min\big\{\max\{h,0\},\,b\big\}\;\in\;[0,b].
\end{equation*}

\emph{Step 2 (The bound in every regime).} By Step~1, $\bar R\le b$ in every regime. Hence $S^*=\max\{\hat R-\bar R,0\}\ge\max\{\hat R-b,0\}$, which is the bound stated before the regimes.
% [CHECK] The original proof calls this bound "part (iii)"; the statement has no part (iii). See ledger, ISSUE T1.

\emph{Step 3 (Cost-Deterred Hiding and Partial Pooling: $\bar R=b$).}
\begin{itemize}
    \item \emph{Cost-Deterred Hiding.} At the equilibrium pool $c_0>\Delta_\psi(\bar\theta)=g_0$, so $h<0$ and Step~1 gives $\bar R=b$.
    \item \emph{Partial Pooling.} The marginal informed low-quality seller is indifferent between hiding and disclosing, $g_0=c_0$ by~\eqref{eq:att_pp}. So $h=0$ and Step~1 gives $\bar R=b$.
\end{itemize}

\emph{Step 4 (Properties of $b$).} The expression $b=\rho\omega_1\beta_1$ contains neither $c_0$, $\pi$, $\psi$ nor $C^{\rm def}$, and $\partial b/\partial(\rho\omega_1)=\beta_1$.
% [CHECK] The statement also signs this derivative (>0); the proof does not say beta_1=C_0-C_1>0. See ledger, ISSUE T2.

\emph{Step 5 (Complete Adverse Selection, $\sigma_1^*=0$).} In Complete Adverse Selection $c_0<\Delta_\psi(e(1))$, so $h=g_0-c_0>0$ at the pool $e(1)$. With $\sigma_1^*=0$, $g_1\le c_0$ gives $h\le b$ (since $h-b=g_1-c_0$). Step~1 then gives
\begin{equation*}
    \bar R=b-h=b\,[1-\tilde e(1,0)]+c_0,\qquad \tilde e(1,0)=\psi e(1)+(1-\psi)e^{\rm def}.
\end{equation*}
$h>0$ gives $\bar R<b$, and $h\le b$ gives $\bar R\ge0$.

\emph{Step 6 (Complete Adverse Selection, $\sigma_1^*>0$).} Now $g_1\ge c_0$, which gives $h\ge b$. So $\max\{b,h\}=h$, and Step~1 gives $\bar R=0$.

\emph{Step 7 (Derivatives in Table~\ref{tab:threshold_cs}).} Write the threshold of Step~5 as
\begin{equation*}
    \bar R=b\,[1-\psi e(1)-(1-\psi)e^{\rm def}]+c_0.
\end{equation*}
Differentiating,
\begin{align*}
    \partial\bar R/\partial c_0&=1, & \partial\bar R/\partial\psi&=b\,[e^{\rm def}-e(1)],\\
    \partial\bar R/\partial\pi&=-b\,\psi\,e'(1),\quad e'(1)=\frac{\bar\theta(\bar\theta-1)}{(1-\pi\bar\theta)^2}<0, & \partial\bar R/\partial(\rho\omega_1)&=\beta_1(1-\tilde e).
\end{align*}
The last uses that $e(1)$ and $e^{\rm def}$ do not depend on $\rho\omega_1$.
% [CHECK] The table's C^def row, rho omega_1 (1-psi), is not derived here. See ledger, ISSUE T3.

\emph{Step 8 (Full attention).} At $\psi=\omega_1=1$, $\tilde e=e$, and only $\sigma_1^*=0$ survives in Complete Adverse Selection (Proposition~\ref{prop:att_equilibrium}(ii)). This is the full-attention statement.
\end{proof}

\subsubsection*{Proof of Proposition~\ref{prop:subsidy_ge} (Subsidies and Disclosure)}

\begin{proof}\upshape
\emph{Idea.} In Partial Pooling the indifference of the marginal hider pins the pool belief and hence the number of hiders, so renovators come out of the disclosers until these run out; after that the pool improves and the threshold falls. \emph{Notation.} $F$: share of informed low-quality sellers who renovate; $\sigma$: share of the unrenovated ones who hide; $e(1;F)$ as in the statement.

\emph{Step 1 (The pool belief).} Renovators become informed high-quality sellers who disclose (Lemma~\ref{lem:renov_behaviour}). The silent pool therefore consists of the uninformed, of whom $(1-\pi)\bar\theta$ per listing are efficient, and of the unrenovated informed low-quality sellers who hide. If a share $\sigma$ of the $\pi(1-\bar\theta)(1-F)$ unrenovated informed low-quality sellers hide, the pool belief is
\begin{equation*}
    e\;=\;\frac{(1-\pi)\bar\theta}{(1-\pi)+\pi(1-\bar\theta)(1-F)\,\sigma}.
\end{equation*}

\emph{Step 2 (Phase~1: the hider mass does not move).} In Partial Pooling the premium equals the cost, $\Delta_\psi(e^*)=c_0$, which pins $e=e^*$. Solving the display of Step~1 for the mass of hiders gives
\begin{equation*}
    H^*\;\equiv\;\pi(1-\bar\theta)(1-F)\sigma\;=\;\frac{(1-\pi)(\bar\theta-e^*)}{e^*},
\end{equation*}
which does not depend on $F$.

\emph{Step 3 (Phase~1: hiding share and threshold).} By Step~2 the share that hides is $\sigma=\sigma_0^*/(1-F)$. It is feasible ($\sigma\le1$) if and only if $F\le1-\sigma_0^*$. In that range $h=g_0-c_0=0$, so $\bar R=b$ by Proposition~\ref{prop:threshold}.

\emph{Step 4 (Phase~2: every remaining seller hides).} Let $F>1-\sigma_0^*$. Even $\sigma=1$ then gives a hider mass $\pi(1-\bar\theta)(1-F)<H^*$, so the belief $e(1;F)$ exceeds $e^*$. Hiding therefore pays strictly, $\Delta_\psi(e(1;F))>c_0$, and every remaining informed low-quality seller hides.

\emph{Step 5 (Phase~2: the threshold falls).} Using $\Delta_\psi(e^*)=c_0$ from Step~2,
\begin{equation*}
    h(F)\;=\;\Delta_\psi(e(1;F))-c_0\;=\;\psi b\,[e(1;F)-e^*]\;>\;0.
\end{equation*}
Since $e(1;F)\le\bar\theta$ and $\psi\le1$, $h(F)\le b$. So a renovated seller displays (Lemma~\ref{lem:renov_behaviour}) and $\bar R(F)=b-h(F)$. The belief $e(1;F)$ is increasing in $F$, so $\bar R(F)$ is decreasing.

\emph{Step 6 (Equilibrium: uniqueness).} The renovating share solves $F=F_R(\bar R(F)+S)$, with $F_R$ a continuous distribution function. The map $F\mapsto F_R(\bar R(F)+S)$ is non-increasing, because $\bar R(F)$ is non-increasing (Steps~3 and~5) and $F_R$ is non-decreasing. Hence
\begin{equation*}
    G(F)\;\equiv\;F-F_R(\bar R(F)+S)
\end{equation*}
is strictly increasing. Since $G(0)\le0\le G(1)$, the root is unique.
% [CHECK] Existence of the root (intermediate value theorem) also needs Rbar(F) continuous, including at F=1-sigma_0^*, where h=0; the proof uses continuity of F_R without saying so. See ledger, ISSUE S1.

\emph{Step 7 (Common cost).} Under a common cost, all renovate if and only if $\hat R-S\le\bar R(1)$. Since $e(1;1)=\bar\theta$, Step~5 gives
\begin{equation*}
    \bar R(1)\;=\;b-\psi b\,(\bar\theta-e^*).
\end{equation*}
\end{proof}

\subsubsection*{Proof of Proposition~\ref{prop:instruments} (Instruments)}

\begin{proof}\upshape
\emph{Idea.} Every instrument except full attention leaves the rating premium $b$ unchanged, so the threshold stays at $b$ as long as the market stays outside Complete Adverse Selection; the work is to check which regime the market lands in. \emph{Notation.} As in the opening paragraph; hats denote values at the estimates, and $e^{\rm def}\ge\bar\theta$ throughout.

\medskip\noindent\emph{Raising the share of informed sellers $\pi$ to one.}

\emph{Step 1 (Where $\pi$ enters).} The threshold depends on $\pi$ only through $e(1)$. It therefore depends on $\pi$ only in Complete Adverse Selection (Proposition~\ref{prop:threshold}).

\emph{Step 2 (Price of silence at $\pi=1$).} At $\pi=1$ the silent pool consists of hiders only. If the belief after silence is zero, a silent listing sells for
\begin{equation*}
    b\,(1-\psi)\,e^{\rm def}.
\end{equation*}

\emph{Step 3 (The belief after silence).} That belief is the limit of the pool belief as $\pi\to1$. It is the only one that survives the intuitive criterion, since a low type gains more from silence than a high type.
% [CHECK] (a) Along the Partial Pooling path e^* does not depend on pi (eq. att_pp), so the equilibrium pool belief tends to e^*, not to zero; zero is the limit of e(sigma_0) at fixed sigma_0>0. (b) The intuitive criterion also needs the high type's deviation to be equilibrium-dominated. See ledger, ISSUES I1-I2.

\emph{Step 4 (Hiding disappears and $\bar R=b$).} With this belief, hiding is weakly unprofitable when $c_0\ge b(1-\psi)e^{\rm def}$. A market that starts in Partial Pooling has
\begin{equation*}
    c_0\;=\;b\,[\psi e^*+(1-\psi)e^{\rm def}]\;>\;b\,(1-\psi)\,e^{\rm def},
\end{equation*}
so $\sigma_0^*=0$, $h<0$ and $\bar R=b$.
% [CHECK] The strict inequality needs psi e^*>0; at psi=0 it holds with equality, h=0, and Rbar=b still. See ledger, ISSUE I3.

\emph{Step 5 (Premium and hiding rate in the limit).} The premium is discontinuous at $\pi=1$. For $\pi<1$ near one it equals $c_0$, and at $\pi=1$ it falls to $b(1-\psi)e^{\rm def}$. The hiding rate, however, is continuous in the limit. By Step~4, the margin by which $c_0$ exceeds that gain is $b\psi e^*$, about \attPiOneMarginApt{} percent of $c_0$ in apartments and \attPiOneMarginHouse{} percent in houses.

\emph{Step 6 (Population of potential renovators).} The population of potential renovators rises from the informed low-quality share $\pi(1-\bar\theta)$ to $1-\bar\theta$.

\medskip\noindent\emph{Making silence as salient as a rating.}

\emph{Step 7 (What changes).} The premium $b=\rho\omega_1\beta_1$ does not depend on $\omega_0$. Raising $\psi$ to one leaves $c_0$ fixed and makes the gain from hiding
\begin{equation*}
    b\,e(\sigma_0)\;\le\;b\bar\theta.
\end{equation*}
% [CHECK] The fall in the price of silence, lambda=(1-\hat\psi) b (e^def-e^*), stated in the proposition, is not derived in the proof. See ledger, ISSUE I4.

\emph{Step 8 (Hiding falls weakly).} In Partial Pooling the belief that sustains $c_0$ at $\psi=1$ is
\begin{equation*}
    c_0/b\;=\;\hat\psi e^*+(1-\hat\psi)e^{\rm def}\;\ge\;e^*,
\end{equation*}
using $e^{\rm def}\ge\bar\theta\ge e^*$. A higher belief means a lower hiding rate, so hiding falls weakly. If $c_0\ge b\bar\theta$, the market is in Cost-Deterred Hiding (Proposition~\ref{prop:att_equilibrium}(i)), so $\sigma_0^*=0$ and $h<0$ or $h=0$.
% [CHECK] Proposition att_equilibrium(i) needs c_0>b thetabar strictly; at c_0=b thetabar the market is at the Partial Pooling boundary with sigma_0^*=0. The conclusion is unaffected. See ledger, ISSUE I5.

\emph{Step 9 (No market enters Complete Adverse Selection, so nothing on the renovation side moves).} The gain from hiding stays below $c_0$ in every market at the calibrated cost, because $c_0\ge b\,e(1)$ in Partial Pooling (using $e^{\rm def}\ge\bar\theta\ge e(1)$). So no market enters Complete Adverse Selection, and $\bar R=b$ before and after (Proposition~\ref{prop:threshold}). Hence the subsidy $S^*=\max\{\hat R-\bar R,0\}$, the flow $\pi(1-\bar\theta)F_R(\bar R+S)$ and the share of renovators at any $S$ are unchanged.
% [CHECK] "Stays below" is supported only weakly (c_0 >= b e(1)); that is all "no market enters Complete Adverse Selection" needs. See ledger, ISSUE I6.

\emph{Step 10 (Correction at full take-up).} In the equilibrium of Proposition~\ref{prop:subsidy_ge} the correction at full take-up is $\hat\psi b(\bar\theta-e^*)$ before the policy. After it, the pool belief that sustains $c_0$ is $c_0/b$ (Step~8). So the correction is $b\bar\theta-c_0$ in Partial Pooling and zero in Cost-Deterred Hiding. The fall is therefore
\begin{equation*}
    \hat\psi b(\bar\theta-e^*)-(b\bar\theta-c_0)^+.
\end{equation*}

\medskip\noindent\emph{Full attention.}

\emph{Step 11 (The new premium).} Replace $b$ by
\begin{equation*}
    b_{\rm new}\;=\;\rho\beta_1\;=\;b/\omega_1,\qquad \omega_1\in[\hat\rho,1],
\end{equation*}
because $\rho\le1$ and $\hat\rho=\rho\omega_1$.

\emph{Step 12 (Threshold by regime after the policy).}
\begin{itemize}
    \item \emph{Outside Complete Adverse Selection.} The threshold is $b_{\rm new}$, so $\Delta\bar R=b_{\rm new}-b$.
    \item \emph{A market that moves into Complete Adverse Selection.} Here $h=b_{\rm new}e(1)-c_0>0$ and
    \begin{equation*}
        \bar R_{\rm new}\;=\;b_{\rm new}-h\;=\;b_{\rm new}[1-e(1)]+c_0.
    \end{equation*}
\end{itemize}
% [CHECK] Using Rbar=b-h requires h<=b_new; this holds (psi=1, e(1)<1) but is not stated. See ledger, ISSUE I7.

\emph{Step 13 (The new threshold is at least $b$).} The market was in Partial Pooling at the calibrated cost, so
\begin{equation*}
    c_0\;\ge\;b\,[\hat\psi e(1)+(1-\hat\psi)e^{\rm def}]\;\ge\;b\,e(1),
\end{equation*}
using $e^{\rm def}\ge e(1)$. Substituting into Step~12 gives $\bar R_{\rm new}\ge b_{\rm new}[1-e(1)]+b\,e(1)\ge b$.

\emph{Step 14 (Bounds on $\Delta\bar R$).} The threshold never exceeds $b_{\rm new}$, so by Steps~12--13
\begin{equation*}
    0\;\le\;\Delta\bar R\;\le\;b_{\rm new}-b.
\end{equation*}
Finally $b_{\rm new}=\rho\beta_1\le\beta_1=\beta_1^{\rm anchor}$ and $b=\hat b$, which gives the upper bound when $\hat\rho<1$. If $\hat\rho\ge1$ the model requires $\rho=\omega_1=1$, so $b_{\rm new}=b$ and $\Delta\bar R=0$.

\medskip\noindent\emph{Enforcement.}

\emph{Step 15 (Enforcement moves who hides, not $\bar R$).} A rise in $c_0$ does not enter $b$, so it leaves $\bar R=b$ outside Complete Adverse Selection. It lowers $\Phi_\psi$, which moves the market toward Cost-Deterred Hiding (Proposition~\ref{prop:att_equilibrium}).
\end{proof}

\subsubsection*{Proof of Proposition~\ref{prop:att_saving} (Upper Bound on the Saving from Salience)}

\begin{proof}\upshape
\emph{Idea.} Salience can save only what hiders lose, their rent $g_0-c_i$; that rent falls by at most $\lambda$, and the subsidy moves one for one with the rent only over a window of length $\min\{b,\hat R\}$. \emph{Notation.} $c_i$, $g_0$, $r_i$, $\bar R_i$ and $S_i^*$ are defined in Steps~1--2; $\lambda=(1-\hat\psi)b(e^{\rm def}-\hat e)$.

\emph{Step 1 (Gain and hiding rate).} Let $c_i\sim G$ be seller $i$'s hiding cost and
\begin{equation*}
    g_0\;=\;b\,[\hat\psi\hat e+(1-\hat\psi)e^{\rm def}]
\end{equation*}
the gain from hiding at the observed pool. A seller with $c_i<g_0$ hides, so $\hat\sigma_0=G(g_0)$.

\emph{Step 2 (Individual threshold and subsidy).} A hider earns the rent $g_0-c_i$, which renovation destroys up to the premium $b$. With $r_i=(g_0-c_i)^+$, her threshold and minimum subsidy are
\begin{equation*}
    \bar R_i\;=\;b-\min\{(g_0-c_i)^+,b\},\qquad S_i^*\;=\;(\hat R-b+\min\{r_i,b\})^+.
\end{equation*}

\emph{Step 3 (The gain after $\psi\to1$).} Raising $\psi$ to one lowers the gain from hiding to $g_0'=b\,e(\sigma_0')$. Here $\sigma_0'=G(g_0')\le\hat\sigma_0$, because the hiding rate falls (Proposition~\ref{prop:att_equilibrium}). Also $g_0'\ge b\hat e$, because $e(\sigma_0')\ge\hat e$.
% [CHECK] Proposition att_equilibrium treats a single hiding cost; with a cost distribution G the fall in the hiding rate follows from the fixed point sigma=G(b e(sigma)). See ledger, ISSUE A1.

\emph{Step 4 (Each rent falls by at most $\lambda$).} By Step~3,
\begin{equation*}
    g_0-g_0'\;\le\;g_0-b\hat e\;=\;\lambda.
\end{equation*}
By cost:
\begin{itemize}
    \item \emph{$c_i<g_0'$:} the rent falls by $g_0-g_0'\le\lambda$;
    \item \emph{$g_0'\le c_i<g_0$:} the rent falls from $g_0-c_i\le g_0-g_0'\le\lambda$ to zero;
    \item \emph{$c_i\ge g_0$:} the seller has no rent.
\end{itemize}

\emph{Step 5 (Each subsidy falls by at most $\min\{\lambda,b,\hat R\}$).} The function $r\mapsto(\hat R-b+\min\{r,b\})^+$ is non-decreasing, flat for $r\ge b$ and for $r\le b-\hat R$, with slope one in between. So a fall in $r$ of at most $\lambda$ (Step~4) lowers it by at most $\min\{\lambda,b,\hat R\}$: the slope region has length $\min\{b,\hat R\}$.

\emph{Step 6 (The bound).} Only the share $\hat\sigma_0$ of sellers has a positive rent (Steps~1 and~4). Averaging Step~5 over sellers gives the bound $\hat\sigma_0\min\{\lambda,b,\hat R\}$.

\emph{Step 7 (Cases with zero saving).}
\begin{itemize}
    \item \emph{Cost-Deterred Hiding.} $\hat\sigma_0=0$.
    \item \emph{A programme that reaches only the share $1-\hat\sigma_0$ of sellers who disclose today.} It pays them $\hat R-b$ before and after, because their rent is zero.
    % [CHECK] Strictly (\hat R-b)^+; the two coincide when \hat R>=b. See ledger, ISSUE A2.
    \item \emph{One common cost, in Partial Pooling.} The rent is zero for every seller.
\end{itemize}
\end{proof}

\subsection*{Alternative Readings of the Policy Exercise}

\paragraph{Exact pricing.}
The exact-pricing solver of Section~\ref{sec:equilibrium} takes the renovation threshold from the full-attention equilibrium rather than setting it to the rating premium. It gives mean thresholds of \euro{}\attRbarCoreEurTwoFiftyApt{} (apartments) and \euro{}\attRbarCoreEurTwoFiftyHouse{} (houses) per m\textsuperscript{2}, against \euro{}\attScBaseRbarApt{} and \euro{}\attScBaseRbarHouse{} under linear pricing, and a share of markets requiring a subsidy at $\hat R=250$ of \attSharePosCoreTwoFiftyApt{} and \attSharePosCoreTwoFiftyHouse{} percent, against \attScBaseShareTwoFiftyApt{} and \attScBaseShareTwoFiftyHouse{}. The thresholds differ by \attRbarGapApt{} (apartments) and \attRbarGapHouse{} (houses) log points on average. The gap is linear versus exact pricing, because attention does not enter the threshold in the regimes the data support. Figure~\ref{fig:policy_core_vs_att} plots the two thresholds against each other.

\paragraph{Hiding cost read off the premium.}
Reading the hiding cost off each market's premium, $c_0=\hat\Delta_m$, instead of calibrating it to the observed hiding rate, places \attRawNCdApt{} apartment markets in Cost-Deterred Hiding, \attRawNPpApt{} in Partial Pooling and \attRawNCasApt{} in Complete Adverse Selection, \attRawNFlagApt{} of them on the branch where high-quality sellers may also hide and the threshold depends on which equilibrium is selected; the mean threshold is \attRawRbarRatioApt{} of $b$. For houses the counts are \attRawNCdHouse{}, \attRawNPpHouse{} and \attRawNCasHouse{}, with \attRawNFlagHouse{} on that branch, and a mean ratio of \attRawRbarRatioHouse{}. This reading also returns negative costs in \attRawNNegApt{} apartment and \attRawNNegHouse{} house markets, so we do not use it.

\paragraph{Additional exhibits.}
Figure~\ref{fig:policy_sgrid} shows the distribution of the minimum subsidy, Figure~\ref{fig:policy_map500} the map at the medium benchmark, and Figure~\ref{fig:policy_regimes} the regime of each market under each scenario of Table~\ref{tab:att_scenarios}.

\begin{figure}[htbp]
    \centering
    \caption{Distribution of the minimum incentive-compatible subsidy $S^*_m$ across markets}
    \label{fig:policy_sgrid}
    \includegraphics[width=\linewidth]{fig/visuals/fig_att_policy_sgrid.pdf}
    \vspace{0.4em}
    \begin{minipage}{0.97\linewidth}
        \footnotesize\textit{Notes.} Rows are property types and columns the renovation-cost benchmark $\hat R\in\{250,500\}$~\euro{}/m\textsuperscript{2}. Each cell shows the cross-market density of $S^*_m=\max\{\hat R-\hat b_m P_m,0\}$ (euros per m\textsuperscript{2}). The dashed line marks the median conditional on $S^*_m>0$; markets with $S^*_m=0$ appear in the left spike.
    \end{minipage}
\end{figure}

\begin{figure}[htbp]
    \centering
    \caption{Spatial distribution of the minimum incentive-compatible subsidy, $\hat R=500$}
    \label{fig:policy_map500}
    \includegraphics[width=\linewidth]{fig/visuals/fig_att_policy_map_R500.pdf}
    \vspace{0.4em}
    \begin{minipage}{0.97\linewidth}
        \footnotesize\textit{Notes.} As Figure~\ref{fig:policy_map} at $\hat R=500$~\euro{}/m\textsuperscript{2}, in the same fixed bins.
    \end{minipage}
\end{figure}

\begin{figure}[htbp]
    \centering
    \caption{Equilibrium regime of each market, by scenario}
    \label{fig:policy_regimes}
    \includegraphics[width=\linewidth]{fig/visuals/fig_att_policy_regimes.pdf}
    \vspace{0.4em}
    \begin{minipage}{0.97\linewidth}
        \footnotesize\textit{Notes.} Shares of markets in Cost-Deterred Hiding, Partial Pooling and Complete Adverse Selection under each scenario of Table~\ref{tab:att_scenarios}, with the mean hiding rate. The hatched house segment under full attention is an artefact of holding a single hiding cost fixed.
    \end{minipage}
\end{figure}

\begin{figure}[htbp]
    \centering
    \caption{Renovation threshold: exact-pricing solver against the rating premium}
    \label{fig:policy_core_vs_att}
    \includegraphics[width=\linewidth]{fig/visuals/fig_att_policy_rbar_core_vs_att.pdf}
    \vspace{0.4em}
    \begin{minipage}{0.97\linewidth}
        \footnotesize\textit{Notes.} Each point is a market. Horizontal axis: threshold of the exact-pricing solver of Section~\ref{sec:equilibrium}; vertical axis: $\bar R=\hat b$. The dashed line is the 45-degree line. The gap is linear versus exact pricing, not attention.
    \end{minipage}
\end{figure}

\subsection*{Belief and Attention Ceilings}

Figure~\ref{fig:att_ceilings} adds the attention ceiling of Proposition~\ref{prop:att_incentives} to the belief ceiling of Figure~\ref{fig:att_belief}.

\begin{figure}[htbp]
    \centering
    \caption{Non-disclosure premium against the energy-efficiency premium, by market}
    \label{fig:att_ceilings}
    \includegraphics[width=\linewidth]{fig/visuals/fig_att_F2_ceilings.pdf}
    \begin{minipage}{0.95\linewidth}
        \vspace{0.5em}
        \footnotesize\textit{Notes.} Each point is a market; both axes are log points. The solid ray $\hat\Delta=\hat b$ is the belief ceiling of Figure~\ref{fig:att_belief} (shaded region). The dashed ray $\hat\Delta=\bar\eta\,\hat b$ is the attention ceiling at the $\hat b$-weighted mean energy-cost ratio $\bar\eta$ (\attEtaApt{} for apartments, \attEtaHouse{} for houses); the grey fan spans the market-specific ceilings $\eta_m\hat b_m$. Diamonds mark markets above their own ceiling. Of \attNMarketsApt{} apartment markets, \attNDgtBApt{} lie above the belief ceiling and \attNViolApt{} above their own attention ceiling; of \attNMarketsHouse{} house markets, \attNDgtBHouse{} and \attNViolHouse{}. Sampling error in $\hat b_m$ and $\hat\Delta_m$ is not shown.
    \end{minipage}
\end{figure}

\subsection*{A Numerical Illustration}

Take round values close to the apartment means: $\rho=0.5$, $\omega_1=1$, $\omega_0=0.1$ (so $\psi=0.1$), $C_1=185$, $C_0=327$ and $\bar C^{\rm sil}=243$ in kWh-equivalent units, and $C^{\rm def}=0$. The perceived cost of silence is $C^*=24.3$; the energy-efficiency premium is $b=0.5\times(327-185)=71$; the non-disclosure premium is $\Delta=0.5\times(327-24.3)=151$, about twice $b$; and silence sells above a displayed A-to-D rating by $0.5\times(185-24.3)=80$. With $\omega_0=\omega_1$ instead, $C^*=243$ and $\Delta=42$, which is $e\,b$ with $e=0.59$: the prediction of Section~\ref{sec:equilibrium}.


% ============================================================
% BIBLIOGRAPHY
% ============================================================
\clearpage
\bibliographystyle{plainnat}
\bibliography{DPE_library}


\end{document}
